%% Research in Engineering Design -- Springer Nature template (sn-jnl), author-year references.
%DIF LATEXDIFF DIFFERENCE FILE
%DIF DEL old.tex    Mon Oct  5 21:09:22 2026
%DIF ADD main.tex   Mon Oct  5 21:09:22 2026
%DIF 2c2-3
%DIF < \documentclass[pdflatex,sn-basic]{sn-jnl}
%DIF -------
%% Line numbers for review (option lineno); remove the option for the final version. %DIF > 
\documentclass[pdflatex,sn-basic,lineno]{sn-jnl} %DIF > 
%DIF -------

\usepackage{graphicx}
\usepackage{multirow}
\usepackage{amsmath,amssymb,amsfonts}
\usepackage{amsthm}
\usepackage{booktabs}
\usepackage{array}
\usepackage{tikz}
\usetikzlibrary{arrows.meta,positioning,calc}
\usepackage{siunitx}
%DIF 13-14c14-15
%DIF < \usepackage{manyfoot}
%DIF < \usepackage{placeins}   % \FloatBarrier keeps floats from drifting past sections   % the class's begin-document footnote hook expects this package to be loaded
%DIF -------
\usepackage{manyfoot}   % the class's begin-document footnote hook expects this package to be loaded %DIF > 
\usepackage{placeins}   % \FloatBarrier keeps floats from drifting past sections %DIF > 
%DIF -------
\graphicspath{{figures/}}

\raggedbottom

%% Float placement: allow figures and tables to share pages with text and with each other.
\renewcommand{\topfraction}{0.9}
\renewcommand{\bottomfraction}{0.7}
\renewcommand{\textfraction}{0.07}
\renewcommand{\floatpagefraction}{0.8}
\setcounter{topnumber}{3}
\setcounter{totalnumber}{4}
\setlength{\textfloatsep}{10pt plus 3pt minus 2pt}
\setlength{\floatsep}{10pt plus 3pt minus 2pt}
\setlength{\intextsep}{10pt plus 3pt minus 2pt}
\setlength{\belowcaptionskip}{3pt}
%% Float pages fill from the top instead of centring their floats.
\makeatletter
\setlength{\@fptop}{0pt}
\setlength{\@fpsep}{12pt plus 2fil}
\setlength{\@fpbot}{0pt plus 1fil}
\makeatother
%% Table body size: the class's own 8 bp table font (its \footnotesize is 7 pt and its \scriptsize 9 pt).
\newcommand{\tabsize}{\fontsize{8}{9.5}\selectfont}
%% Ragged-right paragraph column for text tables.
\newcolumntype{L}[1]{>{\raggedright\arraybackslash\hspace{0pt}}p{#1}}
%DIF 40a41-43
%% Tables and figures of the Electronic Supplementary Material (numbers written by scripts/make_tables.py). %DIF > 
\newcommand{\esm}[1]{\ifcsname esm@#1\endcsname\csname esm@#1\endcsname\else\textbf{S?}\fi} %DIF > 
\InputIfFileExists{esm_labels}{}{} %DIF > 
%DIF -------

%% Set \anontrue to produce an anonymised copy (no author block, no identifying declarations).
\newcommand{\pending}[1]{\textit{[#1 pending]}}
\newif\ifanon
\anonfalse
%DIF PREAMBLE EXTENSION ADDED BY LATEXDIFF
%DIF UNDERLINE PREAMBLE %DIF PREAMBLE
\RequirePackage[normalem]{ulem} %DIF PREAMBLE
\RequirePackage{color}\definecolor{RED}{rgb}{1,0,0}\definecolor{BLUE}{rgb}{0,0,1} %DIF PREAMBLE
\providecommand{\DIFadd}[1]{{\protect\color{blue}\uwave{#1}}} %DIF PREAMBLE
\providecommand{\DIFdel}[1]{{\protect\color{red}\sout{#1}}}                      %DIF PREAMBLE
%DIF SAFE PREAMBLE %DIF PREAMBLE
\providecommand{\DIFaddbegin}{} %DIF PREAMBLE
\providecommand{\DIFaddend}{} %DIF PREAMBLE
\providecommand{\DIFdelbegin}{} %DIF PREAMBLE
\providecommand{\DIFdelend}{} %DIF PREAMBLE
\providecommand{\DIFmodbegin}{} %DIF PREAMBLE
\providecommand{\DIFmodend}{} %DIF PREAMBLE
%DIF FLOATSAFE PREAMBLE %DIF PREAMBLE
\providecommand{\DIFaddFL}[1]{\DIFadd{#1}} %DIF PREAMBLE
\providecommand{\DIFdelFL}[1]{\DIFdel{#1}} %DIF PREAMBLE
\providecommand{\DIFaddbeginFL}{} %DIF PREAMBLE
\providecommand{\DIFaddendFL}{} %DIF PREAMBLE
\providecommand{\DIFdelbeginFL}{} %DIF PREAMBLE
\providecommand{\DIFdelendFL}{} %DIF PREAMBLE
\newcommand{\DIFscaledelfig}{0.5}
%DIF HIGHLIGHTGRAPHICS PREAMBLE %DIF PREAMBLE
\RequirePackage{settobox} %DIF PREAMBLE
\RequirePackage{letltxmacro} %DIF PREAMBLE
\newsavebox{\DIFdelgraphicsbox} %DIF PREAMBLE
\newlength{\DIFdelgraphicswidth} %DIF PREAMBLE
\newlength{\DIFdelgraphicsheight} %DIF PREAMBLE
% store original definition of \includegraphics %DIF PREAMBLE
\LetLtxMacro{\DIFOincludegraphics}{\includegraphics} %DIF PREAMBLE
\newcommand{\DIFaddincludegraphics}[2][]{{\color{blue}\fbox{\DIFOincludegraphics[#1]{#2}}}} %DIF PREAMBLE
\newcommand{\DIFdelincludegraphics}[2][]{% %DIF PREAMBLE
\sbox{\DIFdelgraphicsbox}{\DIFOincludegraphics[#1]{#2}}% %DIF PREAMBLE
\settoboxwidth{\DIFdelgraphicswidth}{\DIFdelgraphicsbox} %DIF PREAMBLE
\settoboxtotalheight{\DIFdelgraphicsheight}{\DIFdelgraphicsbox} %DIF PREAMBLE
\scalebox{\DIFscaledelfig}{% %DIF PREAMBLE
\parbox[b]{\DIFdelgraphicswidth}{\usebox{\DIFdelgraphicsbox}\\[-\baselineskip] \rule{\DIFdelgraphicswidth}{0em}}\llap{\resizebox{\DIFdelgraphicswidth}{\DIFdelgraphicsheight}{% %DIF PREAMBLE
\setlength{\unitlength}{\DIFdelgraphicswidth}% %DIF PREAMBLE
\begin{picture}(1,1)% %DIF PREAMBLE
\thicklines\linethickness{2pt} %DIF PREAMBLE
{\color[rgb]{1,0,0}\put(0,0){\framebox(1,1){}}}% %DIF PREAMBLE
{\color[rgb]{1,0,0}\put(0,0){\line( 1,1){1}}}% %DIF PREAMBLE
{\color[rgb]{1,0,0}\put(0,1){\line(1,-1){1}}}% %DIF PREAMBLE
\end{picture}% %DIF PREAMBLE
}\hspace*{3pt}}} %DIF PREAMBLE
} %DIF PREAMBLE
\LetLtxMacro{\DIFOaddbegin}{\DIFaddbegin} %DIF PREAMBLE
\LetLtxMacro{\DIFOaddend}{\DIFaddend} %DIF PREAMBLE
\LetLtxMacro{\DIFOdelbegin}{\DIFdelbegin} %DIF PREAMBLE
\LetLtxMacro{\DIFOdelend}{\DIFdelend} %DIF PREAMBLE
\DeclareRobustCommand{\DIFaddbegin}{\DIFOaddbegin \let\includegraphics\DIFaddincludegraphics} %DIF PREAMBLE
\DeclareRobustCommand{\DIFaddend}{\DIFOaddend \let\includegraphics\DIFOincludegraphics} %DIF PREAMBLE
\DeclareRobustCommand{\DIFdelbegin}{\DIFOdelbegin \let\includegraphics\DIFdelincludegraphics} %DIF PREAMBLE
\DeclareRobustCommand{\DIFdelend}{\DIFOaddend \let\includegraphics\DIFOincludegraphics} %DIF PREAMBLE
\LetLtxMacro{\DIFOaddbeginFL}{\DIFaddbeginFL} %DIF PREAMBLE
\LetLtxMacro{\DIFOaddendFL}{\DIFaddendFL} %DIF PREAMBLE
\LetLtxMacro{\DIFOdelbeginFL}{\DIFdelbeginFL} %DIF PREAMBLE
\LetLtxMacro{\DIFOdelendFL}{\DIFdelendFL} %DIF PREAMBLE
\DeclareRobustCommand{\DIFaddbeginFL}{\DIFOaddbeginFL \let\includegraphics\DIFaddincludegraphics} %DIF PREAMBLE
\DeclareRobustCommand{\DIFaddendFL}{\DIFOaddendFL \let\includegraphics\DIFOincludegraphics} %DIF PREAMBLE
\DeclareRobustCommand{\DIFdelbeginFL}{\DIFOdelbeginFL \let\includegraphics\DIFdelincludegraphics} %DIF PREAMBLE
\DeclareRobustCommand{\DIFdelendFL}{\DIFOaddendFL \let\includegraphics\DIFOincludegraphics} %DIF PREAMBLE
%DIF COLORLISTINGS PREAMBLE %DIF PREAMBLE
\RequirePackage{listings} %DIF PREAMBLE
\RequirePackage{color} %DIF PREAMBLE
\lstdefinelanguage{DIFcode}{ %DIF PREAMBLE
%DIF DIFCODE_UNDERLINE %DIF PREAMBLE
  moredelim=[il][\color{red}\sout]{\%DIF\ <\ }, %DIF PREAMBLE
  moredelim=[il][\color{blue}\uwave]{\%DIF\ >\ } %DIF PREAMBLE
} %DIF PREAMBLE
\lstdefinestyle{DIFverbatimstyle}{ %DIF PREAMBLE
	language=DIFcode, %DIF PREAMBLE
	basicstyle=\ttfamily, %DIF PREAMBLE
	columns=fullflexible, %DIF PREAMBLE
	keepspaces=true %DIF PREAMBLE
} %DIF PREAMBLE
\lstnewenvironment{DIFverbatim}{\lstset{style=DIFverbatimstyle}}{} %DIF PREAMBLE
\lstnewenvironment{DIFverbatim*}{\lstset{style=DIFverbatimstyle,showspaces=true}}{} %DIF PREAMBLE
%DIF END PREAMBLE EXTENSION ADDED BY LATEXDIFF

\begin{document}

\title[Evaluating design-stage manufacturability decisions]{Evaluating design-stage manufacturability decisions against the resolution of production}

\ifanon
\author*[1]{\fnm{Anonymised} \sur{for review}}
\affil*[1]{\orgname{Affiliation withheld for review}}
\else
\author*[1]{\fnm{H\"useyin Tayyer} \sur{Canseven}}\email{huseyin.canseven@lut.fi}

\affil*[1]{\orgdiv{Department of Electrical Engineering, LUT School of Energy Systems}, \orgname{LUT University}, \orgaddress{\street{Yliopistonkatu 34}, \city{Lappeenranta}, \postcode{53850}, \country{Finland}}}
\fi

%% BEGIN ABSTRACT
\DIFdelbegin %DIFDELCMD < \abstract{Design-stage manufacturability decisions rest on simulation and machine learning, but whether a prediction is fit to decide is rarely evaluated against production. A framework is proposed that expresses this fitness in a unit set by production, the production floor: the difference that drift and scatter produce between two batches of one design. A rule returning meets, fails or trial is characterised by two operating characteristics in floors: its decisive distance, beyond which at least 95\,\% of its verdicts are correct, and its safe distance, beyond which at most 5\,\% are wrong, evaluated separately for each relation of a new design to the produced evidence: a variant of a produced setting, a new process setting or a new design family. It is evaluated on open data of 18 deep-drawn design--process alternatives of 500 parts each with matched simulations, for twelve rules from the nominal simulation to learned models, each calibrated rule also without the simulation. This relation governed decision fitness more than the model did: with a produced sibling the nearest produced setting decided within about 3 floors, for a new process setting no rule decided closer than 8.5 floors, and for a new family no rule was safe closer than 16 floors. The simulation contributed only to the new family, and learned models matched the nearest produced setting in accuracy while adding calibrated abstention where the produced alternatives resembled the new design. The findings show which design-stage decisions can be trusted and when a physical trial is unavoidable.}
%DIFDELCMD < %%%
\DIFdelend \DIFaddbegin \abstract{Design-stage manufacturability decisions increasingly rest on simulation and machine learning, yet whether a prediction is fit to decide is rarely evaluated against production. A framework is proposed whose unit, the production floor, is the difference that drift and scatter produce between two batches of one design within a run. A rule returning meets, fails or trial is characterised by a decisive distance in floors, beyond which at least 95\,\% of its verdicts are correct, and a safe distance, beyond which at most 5\,\% are wrong, for each relation of a new design to the produced evidence: a variant, a new process setting or a new family. A procedure turns them into a guard band for single verdicts. Thirteen rules were evaluated at a 95\,\% conformance level on open data of 18 deep-drawn design--process alternatives of 500 parts each, mostly process settings on one tool, with matched simulations; one floor of draw-in is 0.05\,mm. With a produced sibling the nearest produced setting decided within 3.4 floors; for a new process setting no rule decided closer than 8.5 floors; in two transfers to a new family none was safe closer than 16. Within a family the choice of rule mattered more than the relation, and requirements at a performance index of 1.33 lay inside every decisive distance. With six to nine produced settings per fit, the simulation, as used, helped only across families, and learned models added abstention rather than accuracy. The framework gives evaluations of AI-based manufacturability support a production-referenced standard.}
\DIFaddend %% END ABSTRACT

\DIFdelbegin %DIFDELCMD < \keywords{design for manufacturing, design decision making, verification and validation, operating characteristic, uncertainty quantification, deep drawing}
%DIFDELCMD < %%%
\DIFdelend \DIFaddbegin \keywords{design for manufacturing, machine learning, design decision making, verification and validation, operating characteristic, deep drawing}
\DIFaddend 

\maketitle

%% ----------------------------------------------------------------------
\section{Introduction}\label{sec:intro}

Whether a design alternative can be made to its requirements is decided long before it is made. In sheet-metal forming, as in most processes, the decision rests on a finite-element simulation \DIFdelbegin \DIFdel{of the process }\DIFdelend \citep{banabic2010sheet}: the \DIFdelbegin \DIFdel{design }\DIFdelend team simulates the alternative, compares the predicted characteristics with the requirements \DIFdelbegin \DIFdel{on the drawing and either releases the design}\DIFdelend \DIFaddbegin \DIFadd{and releases the alternative}\DIFaddend , changes it or sends it to a physical trial. Simulation has moved such decisions forward in the design process \DIFdelbegin \DIFdel{\mbox{%DIFAUXCMD
\citep{thomke1998simulation, thomke2000effect}}\hskip0pt%DIFAUXCMD
}\DIFdelend \DIFaddbegin \DIFadd{\mbox{%DIFAUXCMD
\citep{thomke1998simulation}}\hskip0pt%DIFAUXCMD
}\DIFaddend , and machine learning \DIFdelbegin \DIFdel{has made }\DIFdelend \DIFaddbegin \DIFadd{makes }\DIFaddend the loop faster and \DIFdelbegin \DIFdel{broader---surrogates }\DIFdelend \DIFaddbegin \DIFadd{broader: surrogates }\DIFaddend replace simulations, \DIFdelbegin \DIFdel{data-driven models learn from earlier products, and manufacturability checks are increasingly automated \mbox{%DIFAUXCMD
\citep{yan2026machine, liu2026afgnn}}\hskip0pt%DIFAUXCMD
---but neither }\DIFdelend \DIFaddbegin \DIFadd{models learn manufacturability from earlier designs, and checks are automated \mbox{%DIFAUXCMD
\citep{yan2026machine, ghadai2018learning, zhang2018featurenet}}\hskip0pt%DIFAUXCMD
. Neither }\DIFaddend has changed the question \DIFdelbegin \DIFdel{a designer }\DIFdelend \DIFaddbegin \DIFadd{the team }\DIFaddend has to answer: \emph{can this prediction be \DIFdelbegin \DIFdel{trusted }\DIFdelend \DIFaddbegin \DIFadd{relied on }\DIFaddend to decide?} A simulation that is right on average can \DIFdelbegin \DIFdel{still }\DIFdelend be wrong for the alternative at hand, and a model \DIFdelbegin \DIFdel{that learned from }\DIFdelend \DIFaddbegin \DIFadd{learned on }\DIFaddend earlier products may be asked about a design unlike any of them. Whitney argued in this journal that mechanical design cannot be verified by simulation the way integrated circuits are \citep{whitney1996why}\DIFdelbegin \DIFdel{, and variation in production remains the reason \mbox{%DIFAUXCMD
\citep{thornton2000more}}\hskip0pt%DIFAUXCMD
: the credibility of }\DIFdelend \DIFaddbegin \DIFadd{; production variation, with which development organisations still contend \mbox{%DIFAUXCMD
\citep{thornton2000more}}\hskip0pt%DIFAUXCMD
, adds a further reason. Whether }\DIFaddend a simulation-based decision is \DIFdelbegin \DIFdel{ultimately }\DIFdelend \DIFaddbegin \DIFadd{credible is in the end }\DIFaddend a property of production, and production evidence is \DIFdelbegin \DIFdel{exactly }\DIFdelend what is scarce at the design stage.

This paper treats the design-stage manufacturability decision itself as the object of study and asks three questions. How can the fitness of a design-stage prediction for a manufacturability decision be expressed \DIFdelbegin \DIFdel{independently of the predictor and of the }\DIFdelend \DIFaddbegin \DIFadd{on a common production-referenced scale, for any predictor and }\DIFaddend characteristic (RQ1)? How does that fitness depend on the production evidence \DIFdelbegin \DIFdel{available---on }\DIFdelend \DIFaddbegin \DIFadd{available, on }\DIFaddend its amount and on how the new design relates to what has been produced (RQ2)? And what do simulation and machine learning contribute to such a decision once production evidence exists (RQ3)? The questions \DIFdelbegin \DIFdel{sit between three strands of design research: }\DIFdelend \DIFaddbegin \DIFadd{matter for }\DIFaddend decision-based design \DIFdelbegin \DIFdel{, which treats design as a sequence of decisions under uncertainty \mbox{%DIFAUXCMD
\citep{hazelrigg1998framework, lewis2006decision}}\hskip0pt%DIFAUXCMD
; }\DIFdelend \DIFaddbegin \DIFadd{\mbox{%DIFAUXCMD
\citep{hazelrigg1998framework, lewis2006decision}}\hskip0pt%DIFAUXCMD
, for }\DIFaddend the margins that designers keep between predictions and requirements \citep{eckert2019design, brahma2020margin} \DIFdelbegin \DIFdel{; and }\DIFdelend \DIFaddbegin \DIFadd{and for }\DIFaddend the planning of verification and testing \DIFdelbegin \DIFdel{in product development }\DIFdelend \citep{thomke2001sequential, salado2018mathematical, tahera2019testing}\DIFdelbegin \DIFdel{. Each }\DIFdelend \DIFaddbegin \DIFadd{: each }\DIFaddend needs to know how far a prediction can be \DIFdelbegin \DIFdel{trusted }\DIFdelend \DIFaddbegin \DIFadd{relied on }\DIFaddend before a test is cheaper than a decision\DIFdelbegin \DIFdel{, and none has a production-referenced measure of it}\DIFdelend .

\DIFdelbegin \DIFdel{A framework that provides such a measure is proposed . Its unit is }\DIFdelend \DIFaddbegin \DIFadd{The framework proposed here has }\DIFaddend the \emph{production floor} \DIFaddbegin \DIFadd{as its unit}\DIFaddend : the difference \DIFaddbegin \DIFadd{that drift and scatter produce }\DIFaddend between two batches of consecutive parts of one \DIFdelbegin \DIFdel{and the same design that drift and scatter produce at a chosen confidence}\DIFdelend \DIFaddbegin \DIFadd{design within a run, a floor in the sense of a noise floor rather than of the shop floor}\DIFaddend . A rule that turns the design-stage evidence into \DIFdelbegin \DIFdel{one of three verdicts---the }\DIFdelend \DIFaddbegin \DIFadd{a verdict (the }\DIFaddend alternative \emph{meets} the requirement, \emph{fails} it, or goes to a physical \emph{trial}\DIFdelbegin \DIFdel{---is }\DIFdelend \DIFaddbegin \DIFadd{) is }\DIFaddend characterised by two operating characteristics \DIFdelbegin \DIFdel{expressed in floors: its }\emph{\DIFdel{decisive distance}}%DIFAUXCMD
\DIFdel{, beyond which at least 95\,\% of its verdicts are correct, }\DIFdelend \DIFaddbegin \DIFadd{in floors, its }\emph{\DIFadd{decisive distance}} \DIFaddend and its \emph{safe distance}, \DIFdelbegin \DIFdel{beyond which at most 5\,\% are wrong. Both are functions of how far the requirement lies from the truth, in the way the operating characteristic of an acceptance plan is a function of the true quality, and between them the rule abstains rather than errs. They are evaluated separately }\DIFdelend for each \emph{evidence relation} between a new design and the produced \DIFdelbegin \DIFdel{alternatives---a variant whose process setting has been produced with another, unresolvable variation; a new process setting within or outside the produced range; and a new design family---together with a calibration budget, an applicability guard and a six-step procedure}\DIFdelend \DIFaddbegin \DIFadd{alternatives, and a procedure turns them into a guard band on the margin of a single verdict. The decisions in scope are those of a team that has produced variants of a design family: the release of a further variant or process setting in process planning and tryout, and the transfer of production knowledge to a new family; early product design without production evidence of the family is covered only by the new-family relation}\DIFaddend .

The framework is evaluated on open data that make \DIFdelbegin \DIFdel{such an evaluation }\DIFdelend \DIFaddbegin \DIFadd{this }\DIFaddend possible: 9,000 deep-drawn and cut cups of dual-phase steel, produced as 18 design--process alternatives of 500 consecutive parts each \DIFdelbegin \DIFdel{, scanned after drawing and after cutting \mbox{%DIFAUXCMD
\citep{baum2026rddac}}\hskip0pt%DIFAUXCMD
, together }\DIFdelend \DIFaddbegin \DIFadd{and scanned \mbox{%DIFAUXCMD
\citep{baum2026rddac}}\hskip0pt%DIFAUXCMD
, }\DIFaddend with finite-element simulations of the same process \DIFdelbegin \DIFdel{and tools \mbox{%DIFAUXCMD
\citep{baum2025ddacs}}\hskip0pt%DIFAUXCMD
. Each produced }\DIFdelend \DIFaddbegin \DIFadd{\mbox{%DIFAUXCMD
\citep{baum2025ddacs}}\hskip0pt%DIFAUXCMD
. Within each of two cup geometries the alternatives differ in blank-holder force and lubrication on one tool, so two of the three relations concern process settings, and the two geometries give two transfers between families. Each }\DIFaddend alternative is judged as a new \DIFdelbegin \DIFdel{design---by leaving it out alone, or with every alternative that shares its process setting, or with its whole family---and }\DIFdelend \DIFaddbegin \DIFadd{design, and }\DIFaddend the verdicts of \DIFdelbegin \DIFdel{twelve rules }\DIFdelend \DIFaddbegin \DIFadd{thirteen rules at a 95\,\% conformance level }\DIFaddend are compared with what its 500 parts \DIFdelbegin \DIFdel{actually }\DIFdelend did. In the terms of the design research methodology \citep{blessing2009drm}, \DIFdelbegin \DIFdel{the paper }\DIFdelend \DIFaddbegin \DIFadd{this }\DIFaddend is a prescriptive study \DIFdelbegin \DIFdel{whose framework and procedure are given an initial evaluation on production data; that evaluation establishes }\DIFdelend \DIFaddbegin \DIFadd{with an initial support evaluation: it shows }\DIFaddend whether the framework discriminates between rules and evidence situations, not \DIFdelbegin \DIFdel{how design teams use it}\DIFdelend \DIFaddbegin \DIFadd{whether design teams decide better with it, and its findings on one dataset are hypotheses for confirmation}\DIFaddend .

The contributions are \DIFdelbegin \DIFdel{fourfold. The first is the frameworkitself (Section~\ref{sec:method}): the production floor with its effect-to-scatter ratio and minimum resolvable change, and the decisive and safe }\DIFdelend \DIFaddbegin \DIFadd{the framework, with the floor positioned as a precision limit of batch centres and the }\DIFaddend distances as operating characteristics \DIFdelbegin \DIFdel{, with their assumptions, properties and }\DIFdelend \DIFaddbegin \DIFadd{with }\DIFaddend an exact estimator \DIFdelbegin \DIFdel{. The second is }\DIFdelend \DIFaddbegin \DIFadd{(Section~\ref{sec:method}); }\DIFaddend an evaluation design \DIFdelbegin \DIFdel{for design-stage decision rules }\DIFdelend that separates the evidence relations by grouped cross-validation \citep{roberts2017cross} and isolates what the simulation \DIFdelbegin \DIFdel{contributes }\DIFdelend \DIFaddbegin \DIFadd{adds }\DIFaddend by evaluating every calibrated rule with and without it (Section~\DIFdelbegin \DIFdel{\ref{sec:evaldesign}). The third is the empirical answer to the three questions on deep drawing , including what the production floor contains and how the conclusions depend on the analysis choices }\DIFdelend \DIFaddbegin \DIFadd{\ref{sec:demo}); the empirical answers for deep drawing }\DIFaddend (Section~\ref{sec:results})\DIFdelbegin \DIFdel{. The fourth is a procedure and guidance for design practice that separate what is general from what holds for this dataset }\DIFdelend \DIFaddbegin \DIFadd{; and guidance and a reporting standard for evaluating simulation- and AI-based manufacturability support elsewhere }\DIFaddend (Section~\ref{sec:discussion}). \DIFdelbegin %DIFDELCMD < 

%DIFDELCMD < %%%
\DIFdelend Section~\ref{sec:background} positions the work \DIFdelbegin \DIFdel{; Section~\ref{sec:method} defines the framework; Section~\ref{sec:demo} describes the data, the characteristics and the evaluation design; Section~\ref{sec:results} reports the results; Section~\ref{sec:discussion} discusses what they establish and the threats to validity; }\DIFdelend and Section~\ref{sec:conclusion} concludes\DIFdelbegin \DIFdel{. }\DIFdelend \DIFaddbegin \DIFadd{; }\DIFaddend Appendix~\ref{app:appendix} \DIFdelbegin \DIFdel{gives computation details and additional tables}\DIFdelend \DIFaddbegin \DIFadd{summarises the notation, and the Electronic Supplementary Material (ESM) holds further results}\DIFaddend .

%% ----------------------------------------------------------------------
\section{Background and related work}\label{sec:background}

\subsection{Manufacturability evaluation and machine learning\DIFdelbegin \DIFdel{in design for manufacturing}\DIFdelend }

Design for manufacturing (DFM) \DIFdelbegin \DIFdel{has long codified }\DIFdelend \DIFaddbegin \DIFadd{codifies }\DIFaddend what a process can make into rules, guidelines and cost models applied at the design stage \DIFdelbegin \DIFdel{\mbox{%DIFAUXCMD
\citep{boothroyd2010product, bralla1999design}}\hskip0pt%DIFAUXCMD
}\DIFdelend \DIFaddbegin \DIFadd{\mbox{%DIFAUXCMD
\citep{boothroyd2010product}}\hskip0pt%DIFAUXCMD
}\DIFaddend , and quantitative manufacturability models have been used in this journal to drive design decisions \citep{latrobe2003design}. Process simulation \DIFdelbegin \DIFdel{now }\DIFdelend predicts the outcome of a specific alternative\DIFdelbegin \DIFdel{rather than of a class of features}\DIFdelend , and machine learning learns manufacturability from \DIFdelbegin \DIFdel{historical designs and their production records, from }\DIFdelend \DIFaddbegin \DIFadd{CAD models \mbox{%DIFAUXCMD
\citep{ghadai2018learning, zhang2018featurenet, williams2019design}}\hskip0pt%DIFAUXCMD
, from process data and from }\DIFaddend simulated design spaces \DIFdelbegin \DIFdel{, or from both \mbox{%DIFAUXCMD
\citep{yan2026machine, ballegeer2026bendfm}}\hskip0pt%DIFAUXCMD
}\DIFdelend \DIFaddbegin \DIFadd{\mbox{%DIFAUXCMD
\citep{yan2026machine, ballegeer2026bendfm}}\hskip0pt%DIFAUXCMD
, and generative models propose designs directly \mbox{%DIFAUXCMD
\citep{regenwetter2022deep}}\hskip0pt%DIFAUXCMD
}\DIFaddend . Work in this journal has used data and models to choose \DIFdelbegin \DIFdel{between manufacturing }\DIFdelend processes from shape \citep{benslama2024novel}\DIFdelbegin \DIFdel{, to quantify design robustness \mbox{%DIFAUXCMD
\citep{li2026deriving}}\hskip0pt%DIFAUXCMD
, to reason about failure causes in design FMEA \mbox{%DIFAUXCMD
\citep{liu2026afgnn} }\hskip0pt%DIFAUXCMD
and to show that measures left out of an alternative-selection problem change which alternative is chosen \mbox{%DIFAUXCMD
\citep{eaton2026omitted}}\hskip0pt%DIFAUXCMD
. A recurring difficulty is transfer: a model learned on one process, machine or product family is applied to another \mbox{%DIFAUXCMD
\citep{safdar2024transferability}}\hskip0pt%DIFAUXCMD
, and its reproducibility and reusability are open questions \mbox{%DIFAUXCMD
\citep{xie2025reproducible, xie2026reusability}}\hskip0pt%DIFAUXCMD
. What these studies rarely have is a production record }\DIFdelend \DIFaddbegin \DIFadd{. Transfer between processes, machines and product families is a recurring difficulty \mbox{%DIFAUXCMD
\citep{safdar2024transferability}}\hskip0pt%DIFAUXCMD
. Learned manufacturability models are typically evaluated by held-out accuracy or classification scores on labelled designs \mbox{%DIFAUXCMD
\citep{ghadai2018learning, zhang2018featurenet, williams2019design}}\hskip0pt%DIFAUXCMD
, not by the decisions they support against the production }\DIFaddend of the very alternatives \DIFdelbegin \DIFdel{the model is asked about, so their evaluation stops at prediction accuracy and does not reach the design decision}\DIFdelend \DIFaddbegin \DIFadd{they are asked about}\DIFaddend ; open datasets that would allow more are \DIFdelbegin \DIFdel{increasingly recognised as research }\DIFdelend \DIFaddbegin \DIFadd{recognised as }\DIFaddend contributions in their own right \citep{ahmed2025design}.

\subsection{Design decisions\DIFdelbegin \DIFdel{under uncertainty }\DIFdelend \DIFaddbegin \DIFadd{, design types }\DIFaddend and \DIFdelbegin \DIFdel{their evaluation}\DIFdelend \DIFaddbegin \DIFadd{production evidence}\DIFaddend }

Decision-based design casts design as \DIFdelbegin \DIFdel{choosing between alternatives whose outcomes are uncertain, with a preference over them \mbox{%DIFAUXCMD
\citep{hazelrigg1998framework, lewis2006decision}}\hskip0pt%DIFAUXCMD
. In practice }\DIFdelend \DIFaddbegin \DIFadd{a choice between alternatives with uncertain outcomes \mbox{%DIFAUXCMD
\citep{hazelrigg1998framework, lewis2006decision}}\hskip0pt%DIFAUXCMD
. How much a team knows about a new alternative depends on the kind of design: variant design changes parameters of a proven solution, adaptive design changes parts of it and original design creates a new one \mbox{%DIFAUXCMD
\citep{pahl2007engineering}}\hskip0pt%DIFAUXCMD
, and product-family and platform design organise these changes across products \mbox{%DIFAUXCMD
\citep{simpson2006product, jiao2007product}}\hskip0pt%DIFAUXCMD
. The evidence relations of this paper are their counterparts for production evidence: a new variant of a produced setting, a new process setting of a produced design, and a new family. Production knowledge enters design through process-capability data \mbox{%DIFAUXCMD
\citep{thornton2004variation, tata1999process}}\hskip0pt%DIFAUXCMD
, key characteristics \mbox{%DIFAUXCMD
\citep{thornton1999mathematical} }\hskip0pt%DIFAUXCMD
and tolerance analysis \mbox{%DIFAUXCMD
\citep{chase1991survey}}\hskip0pt%DIFAUXCMD
; a rule that returns the capability of the nearest produced setting is a capability-database look-up. Where }\DIFaddend the uncertainty of a prediction is \DIFaddbegin \DIFadd{not resolved it is }\DIFaddend absorbed by margins \DIFdelbegin \DIFdel{: buffers between what is predicted and what is required, which are seldom made explicit and whose value is rarely known \mbox{%DIFAUXCMD
\citep{eckert2019design, brahma2020margin}}\hskip0pt%DIFAUXCMD
. Testing is }\DIFdelend \DIFaddbegin \DIFadd{\mbox{%DIFAUXCMD
\citep{eckert2019design, brahma2020margin}}\hskip0pt%DIFAUXCMD
, and the part of a margin that covers uncertainty is distinct from the part that absorbs change \mbox{%DIFAUXCMD
\citep{eckert2017safety}}\hskip0pt%DIFAUXCMD
. Testing and prototyping are }\DIFaddend the alternative to a margin\DIFdelbegin \DIFdel{, and its planning---when }\DIFdelend \DIFaddbegin \DIFadd{: when }\DIFaddend to test, how many alternatives to test \DIFdelbegin \DIFdel{in parallel or in sequence, and which verification activities to perform---has }\DIFdelend \DIFaddbegin \DIFadd{and in which order has }\DIFaddend been modelled as a trade between the cost of tests and the cost of errors \DIFdelbegin \DIFdel{\mbox{%DIFAUXCMD
\citep{thomke1998simulation, thomke2001sequential, loch2001parallel, salado2018mathematical}}\hskip0pt%DIFAUXCMD
, with physical and virtual testing }\DIFdelend \DIFaddbegin \DIFadd{\mbox{%DIFAUXCMD
\citep{thomke1998simulation, thomke2001sequential, salado2018mathematical}}\hskip0pt%DIFAUXCMD
, with testing and prototyping strategies }\DIFaddend studied as design activities \DIFdelbegin \DIFdel{in their own right \mbox{%DIFAUXCMD
\citep{tahera2019testing}}\hskip0pt%DIFAUXCMD
. Variation risk management identifies the key characteristics whose variation matters and the processes that control them \mbox{%DIFAUXCMD
\citep{thornton1999mathematical, thornton2000more}}\hskip0pt%DIFAUXCMD
. The framework of this paper supplies a quantity these approaches assume but do not measure : how close to a requirement }\DIFdelend \DIFaddbegin \DIFadd{\mbox{%DIFAUXCMD
\citep{tahera2019testing, camburn2017design}}\hskip0pt%DIFAUXCMD
. When upstream information is fit to act on is the question of overlapped development and preliminary information \mbox{%DIFAUXCMD
\citep{krishnan1997model, terwiesch2002exchanging}}\hskip0pt%DIFAUXCMD
, and set-based concurrent engineering codifies production knowledge as feasibility limits for new designs \mbox{%DIFAUXCMD
\citep{sobek1999toyota}}\hskip0pt%DIFAUXCMD
. All of these assume a measure of how far }\DIFaddend a design-stage prediction can \DIFdelbegin \DIFdel{decide, }\DIFdelend \DIFaddbegin \DIFadd{be relied on; the framework supplies one }\DIFaddend in a unit set by production\DIFdelbegin \DIFdel{, and therefore how large a margin has to be, or when a test is unavoidable. The evaluation of design methods themselves has its own methodology: the design research methodology separates the development of support from its evaluation \mbox{%DIFAUXCMD
\citep{blessing2009drm}}\hskip0pt%DIFAUXCMD
, the validation square distinguishes the structural and performance validity of a method \mbox{%DIFAUXCMD
\citep{seepersad2006validation}}\hskip0pt%DIFAUXCMD
, and evaluation designs from medicine have been proposed for design methods \mbox{%DIFAUXCMD
\citep{frey2006validation}}\hskip0pt%DIFAUXCMD
}\DIFdelend .

\subsection{Simulation credibility, calibration and validation\DIFdelbegin \DIFdel{metrics}\DIFdelend }

Verification and validation \DIFdelbegin \DIFdel{methodology asks }\DIFdelend \DIFaddbegin \DIFadd{ask }\DIFaddend whether a model is accurate enough for its intended use \DIFdelbegin \DIFdel{, and answers it by comparing predictions with experiments under a stated validation metric \mbox{%DIFAUXCMD
\citep{oberkampf2010verification, roy2011comprehensive}}\hskip0pt%DIFAUXCMD
; standards for computational solid mechanics and for the credibility of decisions based on models formalise the same question \mbox{%DIFAUXCMD
\citep{asme2020standard, nasa2024standard}}\hskip0pt%DIFAUXCMD
. Validation metrics compare predicted and observed distributions in physical units \mbox{%DIFAUXCMD
\citep{ferson2008model, liu2011toward}}\hskip0pt%DIFAUXCMD
, and predictive maturity tracks how the predictive capability of a calibrated model improves with validation experiments \mbox{%DIFAUXCMD
\citep{hemez2010defining}}\hskip0pt%DIFAUXCMD
}\DIFdelend \DIFaddbegin \DIFadd{\mbox{%DIFAUXCMD
\citep{oberkampf2010verification}}\hskip0pt%DIFAUXCMD
, formalised in risk-informed credibility assessment, which scales the required evidence with the consequence of a wrong decision \mbox{%DIFAUXCMD
\citep{asme2018assessing}}\hskip0pt%DIFAUXCMD
. The distinction between the validation domain and the application domain of a model \mbox{%DIFAUXCMD
\citep{oberkampf2010verification} }\hskip0pt%DIFAUXCMD
is the question that the evidence relation asks of a design}\DIFaddend . Bayesian calibration adds a \DIFdelbegin \DIFdel{model of the discrepancy between simulator and reality, typically a Gaussian process over the inputs, together with calibration parameters of the simulator \mbox{%DIFAUXCMD
\citep{kennedy2001bayesian, higdon2008computer, bayarri2007framework}}\hskip0pt%DIFAUXCMD
; }\DIFdelend \DIFaddbegin \DIFadd{discrepancy model and calibration parameters to a simulator \mbox{%DIFAUXCMD
\citep{kennedy2001bayesian}}\hskip0pt%DIFAUXCMD
, multi-fidelity models scale a cheap model to an expensive one \mbox{%DIFAUXCMD
\citep{kennedy2000predicting}}\hskip0pt%DIFAUXCMD
, and }\DIFaddend omitting the discrepancy biases the \DIFdelbegin \DIFdel{calibrated parameters \mbox{%DIFAUXCMD
\citep{brynjarsdottir2014learning}}\hskip0pt%DIFAUXCMD
, and the identifiability of discrepancy and parameters limits what few experiments can teach \mbox{%DIFAUXCMD
\citep{arendt2012quantification}}\hskip0pt%DIFAUXCMD
}\DIFdelend \DIFaddbegin \DIFadd{parameters \mbox{%DIFAUXCMD
\citep{brynjarsdottir2014learning}}\hskip0pt%DIFAUXCMD
}\DIFaddend . Where a predictive model may be \DIFdelbegin \DIFdel{trusted }\DIFdelend \DIFaddbegin \DIFadd{relied on }\DIFaddend at all is the question of its valid input domain \citep{malak2010using}. Conformal prediction gives distribution-free intervals with a finite-sample coverage guarantee under exchangeability \DIFdelbegin \DIFdel{\mbox{%DIFAUXCMD
\citep{vovk2005algorithmic, lei2018distribution, barber2021predictive, angelopoulos2023conformal}}\hskip0pt%DIFAUXCMD
}\DIFdelend \DIFaddbegin \DIFadd{\mbox{%DIFAUXCMD
\citep{lei2018distribution, barber2021predictive}}\hskip0pt%DIFAUXCMD
}\DIFaddend , including for quantiles \citep{romano2019conformalized}\DIFdelbegin \DIFdel{and under a shift of the inputs \mbox{%DIFAUXCMD
\citep{tibshirani2019conformal} }\hskip0pt%DIFAUXCMD
}\DIFdelend \DIFaddbegin \DIFadd{, under covariate shift \mbox{%DIFAUXCMD
\citep{tibshirani2019conformal} }\hskip0pt%DIFAUXCMD
and for the control of decision risks \mbox{%DIFAUXCMD
\citep{angelopoulos2024conformal}}\hskip0pt%DIFAUXCMD
}\DIFaddend , and it has begun to appear in engineering design \DIFdelbegin \DIFdel{and quality assurance \mbox{%DIFAUXCMD
\citep{zong2026trustworthy, heddoub2026conformal}}\hskip0pt%DIFAUXCMD
. These methods produce intervals and accuracy statements; how close to a requirement such an interval can be trusted to decide, and how much production evidence that takes, remains open. }%DIFDELCMD < 

%DIFDELCMD < %%%
\DIFdelend \DIFaddbegin \DIFadd{\mbox{%DIFAUXCMD
\citep{zong2026trustworthy}}\hskip0pt%DIFAUXCMD
. }\DIFaddend In sheet-metal forming \DIFdelbegin \DIFdel{, }\DIFdelend the accuracy of finite-element \DIFdelbegin \DIFdel{predictions---springback above all---has been studied for decades \mbox{%DIFAUXCMD
\citep{wagoner2013advanced, chen2022robustness}}\hskip0pt%DIFAUXCMD
}\DIFdelend \DIFaddbegin \DIFadd{predictions has long been studied \mbox{%DIFAUXCMD
\citep{wagoner2013advanced}}\hskip0pt%DIFAUXCMD
}\DIFaddend , with friction \citep{hol2012advanced} and the unloading behaviour of the material \citep{yoshida2002model} among the main sources of error. \DIFdelbegin \DIFdel{Simulations are used to compensate springback in the tool design \mbox{%DIFAUXCMD
\citep{gan2004die}}\hskip0pt%DIFAUXCMD
, robust process design under material and process scatter is well established \mbox{%DIFAUXCMD
\citep{strano2006optimization, hou2010stochastic, wiebenga2014effect}}\hskip0pt%DIFAUXCMD
, often with metamodels \mbox{%DIFAUXCMD
\citep{jakumeit2005parameter, harsch2016process}}\hskip0pt%DIFAUXCMD
, and data-driven models predict forming outcomes from process data in series production \mbox{%DIFAUXCMD
\citep{havinga2020exploiting, purr2015stamping}}\hskip0pt%DIFAUXCMD
. }\DIFdelend Simulation datasets of deep drawing have been released for surrogate and transfer learning \citep{baum2025ddacs, heinzelmann2025comprehensive}, \DIFdelbegin \DIFdel{with surrogates trained on them \mbox{%DIFAUXCMD
\citep{heinzelmann2025datadriven, baum2025enhancing} }\hskip0pt%DIFAUXCMD
and }\DIFdelend the gap between synthetic and real data \DIFaddbegin \DIFadd{has been }\DIFaddend identified as the obstacle to \DIFdelbegin \DIFdel{their }\DIFdelend industrial use \citep{baum2026addressing}\DIFdelbegin \DIFdel{. The dataset's authors decomposed }\DIFdelend \DIFaddbegin \DIFadd{, and }\DIFaddend the deviation between simulation and parts \DIFdelbegin \DIFdel{and found that the part geometry explains 77--92\,\% of its variance and that gradient boosting with process measurements explains 81--92\,\% \mbox{%DIFAUXCMD
\citep{baum2026statistical}}\hskip0pt%DIFAUXCMD
; the present paper asks the complementary design question of }\DIFdelend \DIFaddbegin \DIFadd{has been decomposed statistically \mbox{%DIFAUXCMD
\citep{baum2026statistical}}\hskip0pt%DIFAUXCMD
. These methods produce intervals and accuracy statements; }\DIFaddend how close to a requirement \DIFdelbegin \DIFdel{a decision taken before production, without process measurements of the new alternative, can be trusted}\DIFdelend \DIFaddbegin \DIFadd{such an interval can be relied on to decide, and how much production evidence that takes, remains open}\DIFaddend .

\subsection{Decisions with abstention and their operating characteristics}

\DIFdelbegin \DIFdel{Three fields evaluate decisions that may abstain. }\DIFdelend In conformity assessment\DIFdelbegin \DIFdel{under measurement uncertainty}\DIFdelend , a measured value proves conformance, proves non-conformance or leaves the decision open, with guard bands set by the \DIFaddbegin \DIFadd{measurement }\DIFaddend uncertainty and the \DIFdelbegin \DIFdel{specific and global }\DIFdelend risks of false acceptance and false rejection as the measures of a decision rule \citep{iso2017decision, jcgm2012evaluation}. In \DIFdelbegin \DIFdel{analytical chemistry, the limit of detection is the smallest quantity a method detects with stated error rates \mbox{%DIFAUXCMD
\citep{currie1968limits, currie1995nomenclature}}\hskip0pt%DIFAUXCMD
, and in acceptance sampling and selection }\DIFdelend \DIFaddbegin \DIFadd{selection procedures and acceptance sampling }\DIFaddend the operating characteristic gives the probability of a decision as a function of the true quality\DIFdelbegin \DIFdel{\mbox{%DIFAUXCMD
\citep{bechhofer1954single}}\hskip0pt%DIFAUXCMD
. In machine learning, a }\DIFdelend \DIFaddbegin \DIFadd{, with an indifference zone inside which no decision is required \mbox{%DIFAUXCMD
\citep{bechhofer1954single}}\hskip0pt%DIFAUXCMD
. A }\DIFaddend classifier with a reject option trades errors against abstentions \citep{chow1970optimum}, summarised by risk--coverage curves \citep{geifman2017selective}\DIFdelbegin \DIFdel{, and interval forecasts are scored by their coverage and width \mbox{%DIFAUXCMD
\citep{gneiting2007strictly}}\hskip0pt%DIFAUXCMD
}\DIFdelend \DIFaddbegin \DIFadd{; learning to defer passes difficult cases to an expert \mbox{%DIFAUXCMD
\citep{madras2018predict}}\hskip0pt%DIFAUXCMD
; and appropriate reliance on automation requires that its reliability be visible to its users \mbox{%DIFAUXCMD
\citep{lee2004trust}}\hskip0pt%DIFAUXCMD
, which matters for design teams working with AI \mbox{%DIFAUXCMD
\citep{zhang2021cautionary}}\hskip0pt%DIFAUXCMD
}\DIFaddend . The decisive and safe distances are operating characteristics of this kind for design-stage manufacturability decisions\DIFdelbegin \DIFdel{; what the framework adds is a unit set by production, which makes them comparable across characteristics and predictors, and their evaluation by evidence relation}\DIFdelend \DIFaddbegin \DIFadd{, and the guard band of the procedure is the design-stage counterpart of a conformity-assessment guard band}\DIFaddend .

\subsection{Production variation and measurement\DIFdelbegin \DIFdel{capability}\DIFdelend }

Statistical quality control separates \DIFdelbegin \DIFdel{the variation of a process into }\DIFdelend short-term and long-term \DIFdelbegin \DIFdel{components and compares it }\DIFdelend \DIFaddbegin \DIFadd{variation and compares them }\DIFaddend with tolerances through capability \DIFdelbegin \DIFdel{indices \mbox{%DIFAUXCMD
\citep{montgomery2019introduction, kane1986process}}\hskip0pt%DIFAUXCMD
}\DIFdelend \DIFaddbegin \DIFadd{and performance indices \mbox{%DIFAUXCMD
\citep{montgomery2019introduction}}\hskip0pt%DIFAUXCMD
, with release criteria such as a performance index of at least 1.33 \mbox{%DIFAUXCMD
\citep{aiag2006production} }\hskip0pt%DIFAUXCMD
and time-dependent process models for drifting processes \mbox{%DIFAUXCMD
\citep{iso2026process}}\hskip0pt%DIFAUXCMD
. Precision experiments define the repeatability limit, the difference that two results under repeatability conditions exceed with a probability of 5\,\% \mbox{%DIFAUXCMD
\citep{iso2025repeatability}}\hskip0pt%DIFAUXCMD
}\DIFaddend , and measurement systems analysis \DIFaddbegin \DIFadd{separates repeatability from stability and }\DIFaddend decides whether a gauge can resolve a tolerance \DIFdelbegin \DIFdel{\mbox{%DIFAUXCMD
\citep{aiag2010measurement, iso2021measurement}}\hskip0pt%DIFAUXCMD
}\DIFdelend \DIFaddbegin \DIFadd{\mbox{%DIFAUXCMD
\citep{aiag2010measurement}}\hskip0pt%DIFAUXCMD
}\DIFaddend . The production floor is a \DIFdelbegin \DIFdel{quantity of the same family, measured between batches of one design: it contains short-term scatter only through the batch average, and drift within a series in full. }%DIFDELCMD < 

%DIFDELCMD < %%%
\DIFdelend \DIFaddbegin \DIFadd{repeatability-type limit of batch centres under within-run conditions; what is new is its use as the unit of a design decision rule. }\DIFaddend Table~\ref{tab:positioning} places the framework among these approaches\DIFdelbegin \DIFdel{. It }\DIFdelend \DIFaddbegin \DIFadd{: it }\DIFaddend does not replace calibration or conformal prediction \DIFdelbegin \DIFdel{; it uses them}\DIFdelend \DIFaddbegin \DIFadd{but scores them, }\DIFaddend as decision rules\DIFdelbegin \DIFdel{and scores them }\DIFdelend \DIFaddbegin \DIFadd{, }\DIFaddend against production in a common unit\DIFdelbegin \DIFdel{, which is what a design team needs to choose between them}\DIFdelend .

\DIFdelbegin %DIFDELCMD < \begin{table}[t]
%DIFDELCMD < \tabsize
%DIFDELCMD < \setlength{\tabcolsep}{3.5pt}
%DIFDELCMD < \caption{Positioning among methods that support or evaluate design-stage manufacturability decisions}\label{tab:positioning}
%DIFDELCMD < \begin{tabular}{@{}L{3.0cm}L{2.5cm}L{3.0cm}L{1.9cm}L{1.6cm}@{}}
%DIFDELCMD < \toprule
%DIFDELCMD < Approach & Evidence & Output & Scored against production & Abstains \\
%DIFDELCMD < \midrule
%DIFDELCMD < DFM rules and cost models \citep{boothroyd2010product} & codified process knowledge & guideline compliance, cost & no & no \\
%DIFDELCMD < Design margins \citep{eckert2019design} & experience, analysis & buffer to the requirement & rarely & no \\
%DIFDELCMD < V\&V, validation metrics \citep{oberkampf2010verification, ferson2008model} & physics model, validation tests & accuracy in physical units & on test cases & no \\
%DIFDELCMD < Bayesian calibration \citep{kennedy2001bayesian} & simulations, experiments & predictive distribution & not as decisions & implicitly \\
%DIFDELCMD < Conformal prediction \citep{lei2018distribution} & exchangeable calibration data & interval with coverage & coverage only & yes \\
%DIFDELCMD < Conformity assessment, OC curves \citep{iso2017decision, bechhofer1954single} & measurement and its uncertainty & accept, reject, undecided; risks & on the measured item & yes \\
%DIFDELCMD < Selective classification \citep{chow1970optimum, geifman2017selective} & labelled data & class or reject; risk--coverage & on held-out data & yes \\
%DIFDELCMD < ML manufacturability models \citep{yan2026machine} & historical or simulated data & class or value & on held-out data & rarely \\
%DIFDELCMD < This paper & simulations, produced alternatives & verdict; decisive and safe distance in floors by evidence relation & yes, per verdict & yes \\
%DIFDELCMD < \botrule
%DIFDELCMD < \end{tabular}
%DIFDELCMD < \end{table}
%DIFDELCMD < %%%
\DIFdelend \DIFaddbegin \begin{table}[t]
\tabsize
\setlength{\tabcolsep}{3.5pt}
\caption{Positioning among methods that support or evaluate design-stage manufacturability decisions}\label{tab:positioning}
\begin{tabular}{@{}L{3.0cm}L{2.5cm}L{3.0cm}L{1.9cm}L{1.6cm}@{}}
\toprule
Approach & Evidence & Output & Scored against production & Abstains \\
\midrule
DFM rules and cost models \citep{boothroyd2010product} & codified process knowledge & guideline compliance, cost & no & no \\
Design margins \citep{eckert2019design} & experience, analysis & buffer to the requirement & rarely & no \\
V\&V, credibility assessment \citep{oberkampf2010verification, asme2018assessing} & physics model, validation tests & accuracy in physical units; required evidence & on test cases & can require tests \\
Bayesian calibration \citep{kennedy2001bayesian} & simulations, experiments & predictive distribution & not as decisions & implicitly \\
Conformal prediction \citep{lei2018distribution} & exchangeable calibration data & interval with coverage & coverage; per decision possible & yes \\
Conformity assessment, OC curves \citep{iso2017decision, bechhofer1954single} & measurement and its uncertainty & accept, reject, undecided; risks & on the measured item & yes \\
Selective classification \citep{chow1970optimum, geifman2017selective} & labelled data & class or reject; risk--coverage & on held-out data & yes \\
ML manufacturability models \citep{ghadai2018learning, yan2026machine} & historical or simulated designs & class or value & on held-out data & rarely \\
This paper & simulations, produced alternatives & verdict; decisive and safe distance in floors by evidence relation & yes, per verdict & yes \\
\botrule
\end{tabular}
\end{table}
\DIFaddend 

%% ----------------------------------------------------------------------
\section{The framework}\label{sec:method}

\subsection{The decision problem \DIFaddbegin \DIFadd{and its scope}\DIFaddend }\label{sec:problem}

A design alternative $a$ is a combination of product geometry and process settings that is produced as a series of parts. A quality characteristic $c$ varies over the parts of a series. A requirement is an upper limit $R$ on the characteristic, or on its absolute value for characteristics whose target is zero, that a share $p$ of the parts must meet: alternative $a$ is \emph{adequate} for $c$ when the $p$-quantile of its parts satisfies $q_p(a,c)\le R$. \DIFdelbegin \DIFdel{The demonstration uses $p=0.95$ and checks $p=0.90$ and $0.99$; requirements on a lower limit }\DIFdelend \DIFaddbegin \DIFadd{Lower }\DIFaddend or two-sided limits follow by a change of sign or by the absolute deviation from the target\DIFaddbegin \DIFadd{; a tolerance of $\pm t$ on a wall angle about its design angle, for example, is the limit $R=t$ on the absolute deviation}\DIFaddend . At the design stage $q_p$ is unknown. A decision rule $r$ maps the available evidence to an interval $[L_r, U_r]$ for $q_p$ and to a verdict
\begin{equation}
V_r = \begin{cases} \text{meets} & U_r \le R,\\ \text{fails} & L_r > R,\\ \text{trial} & \text{otherwise.} \end{cases}
\label{eq:verdict}
\end{equation}
The three verdicts mirror the decision rules of conformity assessment \DIFdelbegin \DIFdel{, in which uncertainty creates a zone where neither conformance nor non-conformance can be proven \mbox{%DIFAUXCMD
\citep{iso2017decision}}\hskip0pt%DIFAUXCMD
; here the uncertainty is that of the }\DIFdelend \DIFaddbegin \DIFadd{\mbox{%DIFAUXCMD
\citep{iso2017decision}}\hskip0pt%DIFAUXCMD
, here for the uncertainty of a }\DIFaddend design-stage prediction. A rule that returns a point ($L_r=U_r$) always decides. The evidence consists of simulations of $a$ and of the production record of other alternatives, the \emph{calibration alternatives}\DIFdelbegin \DIFdel{. The decisions in scope are those of a design team that has produced variants of a design family and has to decide on }\DIFdelend \DIFaddbegin \DIFadd{; it is used by the process planner or tryout engineer for }\DIFaddend a further variant \DIFdelbegin \DIFdel{, a new process setting or a new family before producing it: the release of variants and process settings in process planning and tryout , and the transfer of production knowledge to }\DIFdelend \DIFaddbegin \DIFadd{or process setting of an existing tool, at tool tryout or before production part approval, and by the product designer for }\DIFaddend a new family\DIFdelbegin \DIFdel{in platform and variant design}\DIFdelend \DIFaddbegin \DIFadd{. The demonstration decides on $p=0.95$, which for a normal characteristic corresponds to a performance index of about 0.55. Release criteria of series production, a performance index of 1.33 or more \mbox{%DIFAUXCMD
\citep{aiag2006production}}\hskip0pt%DIFAUXCMD
, concern nonconforming fractions in parts per million that a series of 500 parts cannot estimate without a parametric tail; the framework reaches them through requirements anchored in capability (Section~\ref{sec:res-requirements})}\DIFaddend .

\subsection{Production floor, effect-to-scatter ratio and minimum resolvable change}\label{sec:floor}

Production evidence comes as series of consecutive parts\DIFdelbegin \DIFdel{. Each series is }\DIFdelend \DIFaddbegin \DIFadd{, each }\DIFaddend cut into batches of $B$ consecutive parts, and each batch is summarised by a robust centre $m$ (here the 10\,\% trimmed mean). The production floor of characteristic $c$ is
\begin{equation}
F_c = Q_\alpha\bigl\{\,|m_i - m_j| : i< j \text{ batches of the same alternative}\,\bigr\},
\label{eq:floor}
\end{equation}
the $\alpha$-quantile of the difference between two batches of one design\DIFdelbegin \DIFdel{, pooled over the produced alternatives (here $B=50$, }\DIFdelend \DIFaddbegin \DIFadd{. Under a stationary normal model with a standard deviation $\sigma_\mathrm{b}$ between batch means and $\sigma_\mathrm{w}$ between parts, $F_c = 1.96\sqrt{2}\sqrt{\sigma_\mathrm{b}^2+\sigma_\mathrm{w}^2/B}$ at }\DIFaddend $\alpha=0.95$\DIFdelbegin \DIFdel{and series of ten batches; floors per geometry are used }\DIFdelend \DIFaddbegin \DIFadd{: the repeatability limit of precision experiments \mbox{%DIFAUXCMD
\citep{iso2025repeatability, iso1994accuracy}}\hskip0pt%DIFAUXCMD
, applied to batch centres under within-run conditions. The reference protocol of this paper is $B=50$, ten batches of a 500-part run, all pairs of batches and $\alpha=0.95$, with per-geometry floors }\DIFaddend where geometries differ\DIFdelbegin \DIFdel{)}\DIFdelend . The floor rests on three assumptions: batches of one design differ only by drift and scatter within a run (A1); the runs from which it is estimated are like those of future production (A2); and the \DIFdelbegin \DIFdel{measurement adds little to it }\DIFdelend \DIFaddbegin \DIFadd{measuring system is repeatable and its resolution small against the floor }\DIFaddend (A3, checked in Section~\ref{sec:res-floor}).
\DIFdelbegin \DIFdel{It has three properties that }\DIFdelend \DIFaddbegin 

\DIFadd{Four properties }\DIFaddend matter for its use. It has the unit of the characteristic, so \DIFdelbegin \DIFdel{a margin expressed in floors does not depend on the unit or the origin of the characteristic. Under a stationary normal model with a standard deviation $\sigma_\mathrm{b}$ between batch means and $\sigma_\mathrm{w}$ between parts, $F_c = 1.96\sqrt{2}\sqrt{\sigma_\mathrm{b}^2+\sigma_\mathrm{w}^2/B}$ at $\alpha=0.95$, so the short-term capability enters through $\sigma_\mathrm{w}/\sqrt{B}$ while drift within a run enters in full; the floor obtained after the parts of each series are randomly permuted contains scatter only, and its ratio to the floor measures the drift. And the floor }\DIFdelend \DIFaddbegin \DIFadd{distances in floors are on a common scale for all characteristics and predictors, although equal distances do not mean equal nonconforming fractions, because the ratio of the floor to the scatter of the parts differs between characteristics (Table~\ref{tab:floors}). It }\DIFaddend is a property of a production protocol: \DIFaddbegin \DIFadd{under drift }\DIFaddend it grows with the \DIFdelbegin \DIFdel{length of the series under drift, }\DIFdelend \DIFaddbegin \DIFadd{time between the batches compared, and it }\DIFaddend shrinks with $B$\DIFaddbegin \DIFadd{; $\sigma_\mathrm{b}$}\DIFaddend , \DIFaddbegin \DIFadd{$\sigma_\mathrm{w}$ and the floor at every lag are reported so that it can be recomputed for another protocol (ESM Tables~}\esm{protocol} \DIFaddend and \DIFaddbegin \esm{lags}\DIFadd{). It contains whatever drifts between batches, the process and the measuring system alike, so without a reference artefact scanned through the series it is a production-and-measurement floor. And it }\DIFaddend excludes the variation between runs, which \DIFdelbegin \DIFdel{bounds it from below as a measure }\DIFdelend \DIFaddbegin \DIFadd{makes it a lower bound of the variation }\DIFaddend of series production. \DIFdelbegin \DIFdel{The }\DIFdelend \DIFaddbegin \DIFadd{For a new family, whose floor is unknown before it is produced, the floor of the produced family is used. Because the }\DIFaddend floor resolves batch centres while the requirement concerns a quantile\DIFdelbegin \DIFdel{; it is used as a unit, and }\DIFdelend \DIFaddbegin \DIFadd{, }\DIFaddend the reproducibility of the quantile \DIFdelbegin \DIFdel{itself }\DIFdelend is reported alongside.

Two alternatives $a$ and $b$ are resolvable for $c$ when their effect-to-scatter ratio
\begin{equation}
\mathrm{ESR}_c(a,b) = |\bar m_a - \bar m_b| \,/\, F_c
\label{eq:esr}
\end{equation}
is at least one, $\bar m$ being the centre of all parts of an alternative\DIFaddbegin \DIFadd{; this is a practical threshold, not a test}\DIFaddend . For a continuous design or process parameter $x$ with a local sensitivity $s_c = \partial \bar m/\partial x$, the minimum resolvable change is $F_c/|s_c|$: a change of $x$ by less moves the centre of the characteristic by less than two batches of one design differ. It is \DIFaddbegin \DIFadd{local, }\DIFaddend a statement about centres\DIFdelbegin \DIFdel{; }\DIFdelend \DIFaddbegin \DIFadd{, and }\DIFaddend a change below it can still move the tail of the distribution.

\subsection{Decision rules}\label{sec:rules}

\DIFdelbegin \DIFdel{Twelve }\DIFdelend \DIFaddbegin \DIFadd{Thirteen }\DIFaddend rules span what a design team can do with its evidence (Table~\ref{tab:rules}); each is defined for any set of calibration alternatives\DIFdelbegin \DIFdel{. Two use the simulation alone: the }\DIFdelend \DIFaddbegin \DIFadd{, and the first seven carry the results in the main text. The }\DIFaddend nominal simulation (M0) \DIFdelbegin \DIFdel{and the worst case of a process-window simulation }\DIFdelend \DIFaddbegin \DIFadd{uses the simulation alone. The bias-corrected simulation }\DIFaddend (\DIFdelbegin \DIFdel{M0w), as a robustness simulation is read in practice. Five combine the simulation with the production record. }\DIFdelend M1\DIFdelbegin \DIFdel{adds to the nominal simulation }\DIFdelend \DIFaddbegin \DIFadd{) adds }\DIFaddend the median offset between the calibration alternatives' $q_p$ and their nominal simulation\DIFdelbegin \DIFdel{. Mc calibrates a parameter of the simulator instead: the friction coefficient of each lubrication pattern is chosen so that the simulated draw-in matches the produced alternatives, in the sense of a calibration parameter of \mbox{%DIFAUXCMD
\citet{kennedy2001bayesian} }\hskip0pt%DIFAUXCMD
without a discrepancy model, and the median remaining offset is added. }\DIFdelend \DIFaddbegin \DIFadd{; the scaled simulation (M1s) fits $q_p = a + b\,s_0$ to them by least squares, a linear multi-fidelity correction \mbox{%DIFAUXCMD
\citep{kennedy2000predicting} }\hskip0pt%DIFAUXCMD
that rescales a simulated trend of the wrong size. The envelope rule (}\DIFaddend M2\DIFdelbegin \DIFdel{takes the upper end of the simulation envelope over plausible incoming conditions, }\DIFdelend \DIFaddbegin \DIFadd{) }\DIFaddend adds the median offset and \DIFdelbegin \DIFdel{a margin equal to }\DIFdelend the largest absolute residual of the calibration alternatives \DIFaddbegin \DIFadd{to the upper end of the simulation envelope over plausible incoming conditions}\DIFaddend . M5 models the offset between $q_p$ and the envelope as a Gaussian process \citep{rasmussen2006gaussian} \DIFdelbegin \DIFdel{over the alternatives' descriptors, so that the correction may vary with the design and process settings, and uses its }\DIFdelend \DIFaddbegin \DIFadd{with a }\DIFaddend 90\,\% predictive interval\DIFdelbegin \DIFdel{. M3 is a part-level model: gradient boosting \mbox{%DIFAUXCMD
\citep{friedman2001greedy} }\hskip0pt%DIFAUXCMD
learns the discrepancy between the measured characteristic and the matched simulation from the descriptors and the measured incoming conditions of the calibration parts, and the predicted parts of the new alternative give its $q_p$, with an offset and a margin from a leave-one-out over the calibration alternatives (a jackknife) . Five rules use the production record alone and isolate what the simulation adds: the median $q_p$ of the calibration alternatives (M1n) , the same with the residual margin (M2n), the Gaussian process of }\DIFdelend \DIFaddbegin \DIFadd{, whose noise variance is bounded below by the variance between the produced series of one geometry and force, so that a new series of a produced setting is not predicted more precisely than produced series reproduce each other; M5n is the same process of }\DIFaddend $q_p$ \DIFdelbegin \DIFdel{itself (M5n), the }\DIFdelend \DIFaddbegin \DIFadd{without the simulation. The nearest produced setting (NN) is a capability-database look-up that uses neither simulation nor learning. The six further rules include a friction-calibrated simulation \mbox{%DIFAUXCMD
\citep{kennedy2001bayesian} }\hskip0pt%DIFAUXCMD
(Mc) and }\DIFaddend part-level \DIFdelbegin \DIFdel{boosting model of the characteristic (M4), and the nearest produced setting (NN), which takes the mean $q_p$ of the calibration alternatives of the design's family at the nearest process setting}\DIFdelend \DIFaddbegin \DIFadd{gradient boosting \mbox{%DIFAUXCMD
\citep{friedman2001greedy} }\hskip0pt%DIFAUXCMD
with a jackknife+ interval \mbox{%DIFAUXCMD
\citep{barber2021predictive}}\hskip0pt%DIFAUXCMD
, with and without the simulation (M3, M3n). The nominal simulation and the envelope do not depend on the lubrication pattern, so a lubrication variant of a produced setting gives M0, M1, M1s, M2 and M5 no simulated information that distinguishes it from its siblings}\DIFaddend .

The margins are labelled by what they guarantee. With $n$ calibration alternatives, the largest-residual margin \DIFdelbegin \DIFdel{of M2 and M2n }\DIFdelend covers a new exchangeable residual with probability $n/(n+1)$ \DIFdelbegin \DIFdel{when the offset is fixed and }\DIFdelend \DIFaddbegin \DIFadd{for a fixed offset and at least }\DIFaddend $(n-1)/(n+1)$ \DIFdelbegin \DIFdel{when it is estimated from the same }\DIFdelend \DIFaddbegin \DIFadd{for an estimated one, because the new residual is then covered whenever it falls within the range of the calibration }\DIFaddend residuals; it is the split-conformal quantile \citep{lei2018distribution} at the highest \DIFdelbegin \DIFdel{level that $n$ residuals can attain}\DIFdelend \DIFaddbegin \DIFadd{attainable level}\DIFaddend , not a 90\,\% interval. \DIFdelbegin \DIFdel{The Gaussian-process interval is model-based, and the jackknife margin of M3 }\DIFdelend \DIFaddbegin \DIFadd{In general, a distribution-free interval cannot certify a miss rate below $1/(n+1)$, so a miss rate of 5\,\% needs at least 19 produced alternatives in the relevant relation and one of 10\,\% needs nine. Under exchangeability the jackknife+ interval covers with probability at least 0.8 at the level of 0.9 used here \mbox{%DIFAUXCMD
\citep{barber2021predictive}}\hskip0pt%DIFAUXCMD
, and the Gaussian-process interval is model-based. The level is a decision variable that trades decisiveness for safety, as the cost of errors would set the threshold of a reject option \mbox{%DIFAUXCMD
\citep{chow1970optimum}}\hskip0pt%DIFAUXCMD
; 0.9 is used here, and 0.8 }\DIFaddend and \DIFdelbegin \DIFdel{M4 carries no guarantee \mbox{%DIFAUXCMD
\citep{barber2021predictive}}\hskip0pt%DIFAUXCMD
. Their empirical coverage is reported in Section~\ref{sec:res-ml}}\DIFdelend \DIFaddbegin \DIFadd{0.95 are reported in the robustness analysis}\DIFaddend .

%DIF < % BEGIN GENERATED rules
\DIFdelbegin %DIFDELCMD < \begin{table}[t]
%DIFDELCMD < \tabsize
%DIFDELCMD < \setlength{\tabcolsep}{3.5pt}
%DIFDELCMD < \caption{The decision rules}\label{tab:rules}
%DIFDELCMD < \begin{tabular}{@{}L{3.2cm}L{5.1cm}L{2.0cm}L{1.6cm}@{}}
%DIFDELCMD < \toprule
%DIFDELCMD < Rule & Prediction of $q_p$ & Interval & Production evidence \\
%DIFDELCMD < \midrule
%DIFDELCMD < \multicolumn{4}{@{}l}{\emph{Simulation only}} \\
%DIFDELCMD < M0 nominal simulation & simulation at nominal sheet thickness and friction & point & none \\
%DIFDELCMD < M0w process-window worst case & lowest to highest simulated value over the incoming conditions & range & none \\
%DIFDELCMD < \multicolumn{4}{@{}l}{\emph{Simulation and production record}} \\
%DIFDELCMD < M1 bias-corrected simulation & M0 plus the median offset & point & $q_p$ \\
%DIFDELCMD < Mc friction-calibrated simulation & simulation with friction calibrated on the draw-in, plus the median offset & point & centres, $q_p$ \\
%DIFDELCMD < M2 envelope + residual margin & upper end of the envelope plus the median offset & largest residual & $q_p$ \\
%DIFDELCMD < M3 simulation + boosting & matched simulation plus a boosting model of the discrepancy, part level & jackknife & parts \\
%DIFDELCMD < M5 GP correction of the envelope & envelope plus a Gaussian-process offset over the descriptors & GP, 90\,\% & $q_p$ \\
%DIFDELCMD < \multicolumn{4}{@{}l}{\emph{Production record only}} \\
%DIFDELCMD < M1n median of the produced & median $q_p$ of the calibration alternatives & point & $q_p$ \\
%DIFDELCMD < NN nearest produced setting & mean $q_p$ at the nearest produced process setting of the family & point & $q_p$ \\
%DIFDELCMD < M2n median + residual margin & M1n with the largest-residual margin & largest residual & $q_p$ \\
%DIFDELCMD < M4 boosting, no simulation & boosting model of the characteristic, part level & jackknife & parts \\
%DIFDELCMD < M5n GP of $q_p$ & Gaussian process of $q_p$ over the descriptors & GP, 90\,\% & $q_p$ \\
%DIFDELCMD < \botrule
%DIFDELCMD < \end{tabular}
%DIFDELCMD < \footnotetext{Descriptors: geometry indicator, blank-holder force and lubrication rank (the order of the oil amount). Largest residual: with $n$ calibration alternatives, coverage $n/(n+1)$ at most. GP: the noise variance is bounded below by the sampling variance of the calibration alternatives' $q_p$.}
%DIFDELCMD < \end{table}
%DIFDELCMD < %%%
%DIF < % END GENERATED rules
\DIFdelend \DIFaddbegin \begin{table}[t]
\tabsize
\setlength{\tabcolsep}{3.5pt}
\caption{The decision rules; the first seven carry the results in the main text}\label{tab:rules}
\begin{tabular}{@{}L{3.3cm}L{5.0cm}L{2.0cm}L{1.6cm}@{}}
\toprule
Rule & Prediction of $q_p$ & Interval & Production evidence \\
\midrule
M0 nominal simulation & simulation at nominal sheet thickness and friction & point & none \\
M1 bias-corrected simulation & M0 plus the median offset & point & $q_p$ \\
M1s scaled simulation & $a+b\,$M0, least squares on the calibration alternatives & point & $q_p$ \\
M2 envelope + residual margin & upper end of the envelope plus the median offset & largest residual & $q_p$ \\
M5 GP correction of the envelope & envelope plus a Gaussian-process offset over the descriptors & GP, 90\,\% & $q_p$ \\
M5n GP of $q_p$ & Gaussian process of $q_p$ over the descriptors & GP, 90\,\% & $q_p$ \\
NN nearest produced setting & mean $q_p$ at the nearest produced process setting of the family & point & $q_p$ \\
\midrule
M0w process-window worst case & lowest to highest simulated value over the incoming conditions & range & none \\
Mc friction-calibrated simulation & simulation with friction calibrated on the draw-in, plus the median offset & point & centres, $q_p$ \\
M3 simulation + boosting & matched simulation plus a boosting model of the discrepancy, part level & jackknife+ & parts \\
M3n boosting, no simulation & boosting model of the characteristic, part level & jackknife+ & parts \\
M1n median of the produced & median $q_p$ of the calibration alternatives & point & $q_p$ \\
M2n median + residual margin & M1n with the largest-residual margin & largest residual & $q_p$ \\
\botrule
\end{tabular}
\footnotetext{Descriptors: geometry indicator, blank-holder force and lubrication rank (the order of the oil amount). GP: noise variance bounded below by the variance between the series of one geometry and force. M3n is named M4 in the released code; M6, a conformalised GP, appears only in the robustness analysis.}
\end{table}
\DIFaddend 

\subsection{Operating characteristics: decisive and safe distances}\label{sec:distances}

A rule is scored by judging each produced alternative as a new design at requirements placed at a known distance from the truth,
\begin{equation}
R = q_p(a,c) + d\,F_c ,
\label{eq:requirement}
\end{equation}
so that $a$ is adequate exactly when $d\ge 0$. For a case $(a,c)$ with interval $[L,U]$, let $u=(U-q_p)/F_c$ and $v=(q_p-L)/F_c$ be the distances by which the interval excludes the truth from above and below. At $d=+\delta$ the verdict is correct when $u\le\delta$ and \DIFdelbegin \DIFdel{wrong (}\DIFdelend a false reject \DIFdelbegin \DIFdel{) }\DIFdelend when $-v>\delta$; at $d=-\delta$ it is correct when $v<\delta$ and \DIFdelbegin \DIFdel{wrong (}\DIFdelend a false accept \DIFdelbegin \DIFdel{) }\DIFdelend when $-u\ge\delta$. \DIFdelbegin \DIFdel{Over a population of cases, let }\DIFdelend \DIFaddbegin \DIFadd{Let }\DIFaddend $\kappa(\delta)$ and $\omega(\delta)$ be the shares of correct and of wrong verdicts \DIFaddbegin \DIFadd{among all feasible verdicts }\DIFaddend at $|d|=\delta$\DIFdelbegin \DIFdel{, both signs weighted equally and requirements below zero excluded for characteristics that cannot be negative}\DIFdelend \DIFaddbegin \DIFadd{: every case gives one verdict on an adequate alternative and, where $q_p-\delta F_c\ge0$, one on a failing alternative, and every verdict has equal weight. As $\delta$ grows, failing-side requirements become infeasible for characteristics close to zero, so the failing side carries half the weight up to a few floors and a third at about 100 floors (ESM Table~}\esm{weights}\DIFadd{)}\DIFaddend . The \emph{decisive distance} $\delta_\mathrm{d}$ of a rule is the smallest $\delta$ beyond which $\kappa\ge 0.95$, and its \emph{safe distance} $\delta_\mathrm{s}$ the smallest $\delta$ beyond which $\omega\le 0.05$.

The \DIFdelbegin \DIFdel{two distances have five properties. They }\DIFdelend \DIFaddbegin \DIFadd{distances }\DIFaddend are exact: $\kappa$ and $\omega$ change only where $\delta$ crosses one of the case's $u$, $v$, $-u$ or $-v$, so the distances are order statistics of these exclusion distances\DIFdelbegin \DIFdel{and need no grid}\DIFdelend ; for a rule that always decides, $\delta_\mathrm{d}=\delta_\mathrm{s}$ is \DIFaddbegin \DIFadd{close to }\DIFaddend the 90\,\% quantile of its absolute \DIFdelbegin \DIFdel{prediction }\DIFdelend error in floors\DIFdelbegin \DIFdel{when no requirement is excluded}\DIFdelend . Always $\delta_\mathrm{s}\le\delta_\mathrm{d}$, and between the two the rule sends designs to trial rather than misleading. \DIFdelbegin \DIFdel{They do not depend on the unit of the characteristic. They split into one-sided distances }\DIFdelend \DIFaddbegin \DIFadd{The distances are reported for both sides together and one-sided, }\DIFaddend for failing alternatives \DIFdelbegin \DIFdel{, where the errors are false accepts, and for adequate ones , where they are false rejects}\DIFdelend \DIFaddbegin \DIFadd{(false accepts) and adequate ones (false rejects)}\DIFaddend , which matters when the \DIFdelbegin \DIFdel{two }\DIFdelend errors cost differently \DIFdelbegin \DIFdel{. And they are operating characteristics: }\DIFdelend \DIFaddbegin \DIFadd{(ESM Table~}\esm{onesided}\DIFadd{). They are points on operating characteristics, }\DIFaddend functions of the true margin over a population of design situations\DIFdelbegin \DIFdel{, like the operating characteristic of an acceptance plan \mbox{%DIFAUXCMD
\citep{bechhofer1954single} }\hskip0pt%DIFAUXCMD
or the risk--coverage curve of a selective classifier \mbox{%DIFAUXCMD
\citep{geifman2017selective}}\hskip0pt%DIFAUXCMD
, not statements about a single verdict}\DIFdelend \DIFaddbegin \DIFadd{: the decisive distance is the analogue of the indifference zone of a selection procedure \mbox{%DIFAUXCMD
\citep{bechhofer1954single} }\hskip0pt%DIFAUXCMD
at a probability of correct selection of 0.95, and the pair plays the role of the acceptable and rejectable quality levels of an acceptance plan}\DIFaddend . A design team uses them before a requirement is fixed, to choose a rule for a family and evidence relation, to judge whether a requirement \DIFdelbegin \DIFdel{placed }\DIFdelend near the expected capability can be decided at the design stage at all, and to plan trials and production evidence. \DIFdelbegin \DIFdel{For a single verdict the observable quantity is the margin between the predicted interval and the requirement; the share of correct verdicts as a function of that margin is reported with the results.
}%DIFDELCMD < 

%DIFDELCMD < %%%
\DIFdel{In decision-theoretic terms, let a false accept cost $c_\mathrm{FA}$, a false reject $c_\mathrm{FR}$ and a trial $c_\mathrm{T}$, with correct verdicts free and trials }\DIFdelend \DIFaddbegin \DIFadd{A trial is itself not }\DIFaddend error-free\DIFdelbegin \DIFdel{. Over the population at a true margin beyond $\delta_\mathrm{d}$, the expected cost of a verdict is at most $0.05\,\max(c_\mathrm{FA}, c_\mathrm{FR}, c_\mathrm{T})$, and beyond $\delta_\mathrm{s}$ at most $0.05\,\max(c_\mathrm{FA}, c_\mathrm{FR}) + u(\delta)\,c_\mathrm{T}$ with $u(\delta)$ the share sent to trial; the bound holds on average over adequate and failing alternatives weighted equally, and one-sided distances replace it when the base rates differ. Which rule to prefer then depends on the cost ratios; a worked comparison with realistic requirement margins is given in Section~\ref{sec:res-scenarios}}\DIFdelend \DIFaddbegin \DIFadd{: a short run reproduces the 95th percentile only to about one floor (Section~\ref{sec:res-floor}), so no decision can be relied on closer to the requirement than the trial that would replace it}\DIFaddend .

\subsection{Evidence relation, calibration budget and applicability guard}\label{sec:budget}

How a new design relates to the produced alternatives is part of the decision situation. With a \emph{sibling}, an alternative at the same process setting has been produced with another variation that production may not resolve; a \emph{new process setting} lies between produced settings (interpolation) or outside them (extrapolation); and a \emph{new family} shares no produced alternative. The \DIFdelbegin \DIFdel{two }\DIFdelend distances are evaluated for each relation. \DIFaddbegin \DIFadd{Beyond a factorial, the relation can be graded by a distance in descriptor space or by the effect-to-scatter ratio between the predicted characteristic of the new design and its nearest produced alternative (Section~\ref{sec:reuse}). }\DIFaddend The calibration budget is \DIFdelbegin \DIFdel{their dependence }\DIFdelend \DIFaddbegin \DIFadd{the dependence of the distances }\DIFaddend on the number $k$ of produced alternatives, over all subsets of size $k$ \DIFdelbegin \DIFdel{, and the calibration design asks which alternatives to produce first: subsets are classified by the relation of }\DIFdelend \DIFaddbegin \DIFadd{classified by their relation to }\DIFaddend the new design\DIFdelbegin \DIFdel{to them}\DIFdelend . An applicability guard returns \emph{trial} when the new design lies outside what the calibration has seen: its geometry is not among the calibration alternatives (family guard), a descriptor lies outside the calibrated range (range guard), or its simulated $q_p$ lies outside the range of the calibration alternatives' simulated values (novelty guard)\DIFdelbegin \DIFdel{; all three use design-stage information only}\DIFdelend , in the spirit of a valid input domain \citep{malak2010using}.

\subsection{Procedure}\label{sec:procedure}

Figure~\ref{fig:procedure} turns the framework into six steps\DIFdelbegin \DIFdel{. Steps~}\DIFdelend \DIFaddbegin \DIFadd{: steps }\DIFaddend 1 to 5 \DIFdelbegin \DIFdel{are carried out }\DIFdelend once per design family\DIFdelbegin \DIFdel{and }\DIFdelend \DIFaddbegin \DIFadd{, }\DIFaddend repeated when production evidence is added\DIFdelbegin \DIFdel{; step ~}\DIFdelend \DIFaddbegin \DIFadd{, and step }\DIFaddend 6 \DIFdelbegin \DIFdel{is carried out }\DIFdelend for every new design. Step \DIFdelbegin \DIFdel{~}\DIFdelend 1 frames the decision \DIFdelbegin \DIFdel{: the family, its alternatives, the characteristicsand the }\DIFdelend \DIFaddbegin \DIFadd{(family, alternatives, characteristics, }\DIFaddend form of the requirement \DIFdelbegin \DIFdel{. Step~}\DIFdelend \DIFaddbegin \DIFadd{and conformance level); step }\DIFaddend 2 measures the floor \DIFdelbegin \DIFdel{on the produced alternatives, and step ~}\DIFdelend \DIFaddbegin \DIFadd{with a stated protocol; step }\DIFaddend 3 states which differences between alternatives, and which local changes of a parameter, production can resolve; \DIFdelbegin \DIFdel{a design question below the floor is not one the design stage needs to answer finer. Step~}\DIFdelend \DIFaddbegin \DIFadd{and step }\DIFaddend 4 calibrates the candidate rules \DIFdelbegin \DIFdel{, }\DIFdelend with and without the simulation\DIFdelbegin \DIFdel{, and step~}\DIFdelend \DIFaddbegin \DIFadd{. Step }\DIFaddend 5 qualifies them by their \DIFdelbegin \DIFdel{decisive and safe }\DIFdelend distances for each evidence relation \DIFdelbegin \DIFdel{, their calibration budget and the guard; when }\DIFdelend \DIFaddbegin \DIFadd{and fixes for each rule and relation a }\emph{\DIFadd{guard band}} \DIFadd{$g$: the smallest observable margin between the predicted interval and the requirement beyond which at least 95\,\% of the decided verdicts were correct, for requirements within 20 floors of the truth. If }\DIFaddend the distances are too large for the requirements \DIFdelbegin \DIFdel{expected }\DIFdelend of coming designs, another variant is produced\DIFdelbegin \DIFdel{, chosen by the calibration design. Step ~}\DIFdelend \DIFaddbegin \DIFadd{. Step }\DIFaddend 6 \DIFdelbegin \DIFdel{decides: it }\DIFdelend establishes the relation of the new design\DIFdelbegin \DIFdel{to the produced evidence}\DIFdelend , applies the rule qualified for \DIFdelbegin \DIFdel{that relation and the guard, and returns a verdict }\DIFdelend \DIFaddbegin \DIFadd{it and the guards, widens the rule's interval by $g$ floors and returns meets or fails }\DIFaddend with the margin \DIFdelbegin \DIFdel{of the requirement }\DIFdelend in floors, or a trial when \DIFdelbegin \DIFdel{the guard firesor the margin lies below the rule's safe distance}\DIFdelend \DIFaddbegin \DIFadd{a guard fires, no rule is qualified or the requirement lies inside the widened interval}\DIFaddend .

\begin{figure}[t]
\centering
\input{fig_procedure}
\caption{The procedure: steps 1--5 once per design family, step 6 for every new design; dashed boxes are the evidence it uses}\label{fig:procedure}
\end{figure}

%% ----------------------------------------------------------------------
\section{Demonstration on deep drawing}\label{sec:demo}

\subsection{Alternatives and production data}\label{sec:data}

The production data are the Real Deep Drawing and Cutting dataset (RDDAC) of the University of Stuttgart \citep{baum2026rddac, baum2026statistical}. Square blanks of DP600 dual-phase steel, 210\,mm wide and about 1\,mm thick, were deep drawn into cups (operation OP10) and cut (OP20) on one tool set with two geometries, concave and convex side walls, which differ in the curvature of the base shape, the wall angle and the bottom radius at once. Each geometry was run at three blank-holder forces (100, 300 and 500\,kN) and three lubrication patterns (coarse, medium and fine oil application, with series medians of 1.15--1.48, 0.92--1.23 and 0.85--1.13\,g/m$^2$ of oil), giving 18 alternatives, each produced as one series of 500 consecutive parts\DIFdelbegin \DIFdel{. Within a geometry the new designs are therefore new process settings and lubrication variants on an existing tool}\DIFdelend ; the change of geometry is the only change of product design. For every part the dataset holds \DIFdelbegin \DIFdel{the press force, traverses of the sheet thickness and of the oil film before forming, }\DIFdelend \DIFaddbegin \DIFadd{press signals, sheet and oil-film traverses }\DIFaddend and 3D scans after drawing and \DIFdelbegin \DIFdel{after cutting(0.077}\DIFdelend \DIFaddbegin \DIFadd{cutting; 10 of the 9,000 parts carry no scan. The series start at punch temperatures of 20.5--29.6}\DIFaddend \,\DIFdelbegin \DIFdel{mm per pixel across the scan line,0.158}\DIFdelend \DIFaddbegin \DIFadd{$^\circ$C and warm up by 1.0--5.0\,K over their first 150 parts; their median sheet thickness ranges from 0.986 to 0.993}\DIFaddend \,mm\DIFdelbegin \DIFdel{per line and 0.008}\DIFdelend \DIFaddbegin \DIFadd{, and the 100\,kN series ran at a stroke speed of 47}\DIFaddend \,mm\DIFdelbegin \DIFdel{in height). Of the 9,000 parts , 10 carry no scan, leaving 8,990.
}\DIFdelend \DIFaddbegin \DIFadd{/s against 78--79\,mm/s for the others (ESM Table~}\esm{runs}\DIFadd{). The documentation states neither the order and dates of the series nor whether the parts were scanned in production order.
}\DIFaddend 

\subsection{Simulations}\label{sec:sims}

The simulations come from the Deep Drawing and Cutting Simulations dataset (DDACS) of the same group \citep{baum2025ddacs, heinzelmann2025comprehensive}: LS-DYNA models of the same process on a quarter of the part, with nodes about 1\,mm apart, a scaled DP600 flow curve, Coulomb friction with a constant coefficient \DIFdelbegin \DIFdel{between blank and tools, }\DIFdelend and springback after both operations; the documentation does not state the yield criterion or the unloading model, which govern springback \citep{yoshida2002model}\DIFaddbegin \DIFadd{, and the models do not represent the stroke speed}\DIFaddend . For the RDDAC tools, 396 simulations vary the sheet thickness (0.95--1.00\,mm\DIFdelbegin \DIFdel{in six steps}\DIFdelend ) and the friction coefficient (0.05--0.15\DIFdelbegin \DIFdel{in eleven steps}\DIFdelend ) at each geometry and blank-holder force. A part is matched to the simulation with the nearest sheet thickness and the friction coefficient that the dataset's linear mapping assigns to its \DIFdelbegin \DIFdel{measured oil film}\DIFdelend \DIFaddbegin \DIFadd{oil film, an assumption of the dataset's authors on which Mc, M3 and every matched simulation depend}\DIFaddend . The nominal simulation of an alternative is the one at 0.99\,mm, the grid value nearest the measured median thickness\DIFdelbegin \DIFdel{of 0.990\,mm}\DIFdelend , and friction 0.10, the middle of the range.
\DIFdelbegin \DIFdel{The simulated characteristics carry extraction noise from the 1\,mm node spacing; it is quantified in Section~\ref{sec:res-sim}.
}\DIFdelend 

\subsection{Quality characteristics and their measurement}\label{sec:qcs}

Seven characteristics were extracted \DIFaddbegin \DIFadd{with the same definitions }\DIFaddend from the scans and from the simulated nodes, all referred to the plane of the cup bottom (\DIFdelbegin \DIFdel{Table~\ref{tab:qcs})}\DIFdelend \DIFaddbegin \DIFadd{ESM Table~}\esm{qcs}\DIFadd{): the mid-side and corner draw-in and the flange waviness, process indicators of material flow and wrinkling on the flange that the cut removes, and four characteristics that a drawing would carry: the wall angle (design angle 20$^\circ$ concave and 10$^\circ$ convex), the cup depth, and the arm angle and bottom dome of the cut part, the last two with requirements on their absolute values}\DIFaddend . The four sides of a part measure several of them twice, and this redundancy exposes the measurement (\DIFdelbegin \DIFdel{Table~\ref{tab:scan}). The flange width in the direction }\DIFdelend \DIFaddbegin \DIFadd{ESM Table~}\esm{scan}\DIFadd{): the flange width }\DIFaddend along the scan lines scatters 4 to 18 times more than across them, \DIFdelbegin \DIFdel{uncorrelated with it and without serial correlation, and reads 2.3--3.5\,mm larger; }\DIFdelend \DIFaddbegin \DIFadd{so }\DIFaddend the mid-side draw-in is \DIFdelbegin \DIFdel{therefore taken from the widths }\DIFdelend \DIFaddbegin \DIFadd{taken }\DIFaddend across the scan lines only\DIFdelbegin \DIFdel{. The }\DIFdelend \DIFaddbegin \DIFadd{, and the }\DIFaddend north wall carries \DIFdelbegin \DIFdel{about 0.3$^\circ$ of uncorrelated scatter in every series and reads 0.9--1.6$^\circ$ low on the convex cups, }\DIFdelend \DIFaddbegin \DIFadd{uncorrelated scatter }\DIFaddend and the south wall is shadowed\DIFdelbegin \DIFdel{; }\DIFdelend \DIFaddbegin \DIFadd{, so }\DIFaddend the wall angle is \DIFdelbegin \DIFdel{therefore }\DIFdelend the mean of the east and west walls\DIFdelbegin \DIFdel{, which also cancels a residual tilt of the part (on the convex cups the two are anti-correlated). The corner }\DIFdelend \DIFaddbegin \DIFadd{. The mid-side }\DIFaddend draw-in \DIFdelbegin \DIFdel{uses the mean of the two flange diagonals, which differ by 1.6--2.2}\DIFdelend \DIFaddbegin \DIFadd{is quantised in steps of 0.019}\DIFaddend \,mm\DIFdelbegin \DIFdel{in every series, consistent with a small non-orthogonality of the scan axes to which their mean is insensitive. The quarter-model simulations are symmetric about the diagonal, so their values along both axes estimate the same quantity and are averaged. Differences that remain }\DIFdelend \DIFaddbegin \DIFadd{, 0.38 floors, by the pixel pitch; the other characteristics are effectively continuous. Differences }\DIFaddend between scans and \DIFdelbegin \DIFdel{simulations---the }\DIFdelend \DIFaddbegin \DIFadd{simulations, such as the }\DIFaddend scanned surface against the shell mid-surface, \DIFdelbegin \DIFdel{edge treatment and spatial resolution---affect }\DIFdelend \DIFaddbegin \DIFadd{enter }\DIFaddend the absolute offsets between the two but not the floors\DIFdelbegin \DIFdel{and the effects within production}\DIFdelend .

%DIF < % BEGIN GENERATED qcs
\DIFdelbegin %DIFDELCMD < \begin{table}[t]
%DIFDELCMD < \tabsize
%DIFDELCMD < \setlength{\tabcolsep}{4pt}
%DIFDELCMD < \caption{Quality characteristics, extracted with the same definitions from the scans and the simulated nodes}\label{tab:qcs}
%DIFDELCMD < \begin{tabular}{@{}L{2.6cm}L{0.7cm}L{8.4cm}@{}}
%DIFDELCMD < \toprule
%DIFDELCMD < Characteristic & Unit & Definition \\
%DIFDELCMD < \midrule
%DIFDELCMD < Draw-in, mid-side & mm & $(210 - W_x)/2$; $W_x$: flange width across the cup centre, across the scan lines \\
%DIFDELCMD < Draw-in, corner & mm & $(2\sqrt{2}\cdot 210 - L_1 - L_2)/4$; $L$: distance between opposite corner tips of the flange \\
%DIFDELCMD < Flange waviness & mm & robust standard deviation of the flange about a fitted quadratic surface \\
%DIFDELCMD < Wall angle & $^\circ$ & inclination of the wall from the punch axis between 10 and 20\,mm depth, mean of the east and west walls \\
%DIFDELCMD < Arm angle, cut & $^\circ$ & inclination of the flat segment of the cut arms, mean of four arms (requirement on its absolute value) \\
%DIFDELCMD < Cup depth & mm & median depth of the flange below the plane of the cup bottom after drawing \\
%DIFDELCMD < Bottom dome, cut & mm & sag of the bottom at 40\,mm radius from a fitted paraboloid (requirement on its absolute value) \\
%DIFDELCMD < \botrule
%DIFDELCMD < \end{tabular}
%DIFDELCMD < \end{table}
%DIFDELCMD < %%%
%DIF < % END GENERATED qcs
%DIFDELCMD < 

%DIFDELCMD < %%%
\DIFdelend \subsection{Evaluation design}\label{sec:evaldesign}

Every produced alternative is judged as a new design under \DIFdelbegin \DIFdel{four schemes of }\DIFdelend grouped cross-validation \citep{roberts2017cross}, one \DIFaddbegin \DIFadd{scheme }\DIFaddend per evidence relation. \emph{New variant}: the alternative alone is held out and the other eight alternatives of its geometry calibrate, so two siblings at its blank-holder force \DIFdelbegin \DIFdel{, produced with another lubrication pattern, }\DIFdelend are among them. \emph{New process setting}: \DIFdelbegin \DIFdel{all }\DIFdelend \DIFaddbegin \DIFadd{the }\DIFaddend three alternatives of one \DIFdelbegin \DIFdel{blank-holder }\DIFdelend force of a geometry are held out together and the six others calibrate; the 300\,kN alternatives are \DIFdelbegin \DIFdel{then }\DIFdelend interpolated and the 100 and 500\,kN alternatives extrapolated in force. \emph{New family}: the alternatives of one geometry calibrate the rules for the other\DIFdelbegin \DIFdel{. The pooled }\DIFdelend \DIFaddbegin \DIFadd{, which gives two transfers. Pooled }\DIFaddend versions of the first two add the other geometry's alternatives to the calibration set \DIFaddbegin \DIFadd{and test whether the production record of another family helps}\DIFaddend . Every calibrated rule is evaluated with and without the simulation (M1/M1n, M2/M2n, M3/\DIFdelbegin \DIFdel{M4}\DIFdelend \DIFaddbegin \DIFadd{M3n}\DIFaddend , M5/M5n), so that the paired difference isolates what the simulation contributes \DIFaddbegin \DIFadd{as it is used}\DIFaddend , and the nearest produced setting \DIFdelbegin \DIFdel{(NN) }\DIFdelend is the baseline without simulation and \DIFdelbegin \DIFdel{without }\DIFdelend learning. A scope comprises 126 cases (18 alternatives, 7 characteristics)\DIFdelbegin \DIFdel{and two transfer calibrations}\DIFdelend .

The distances are computed exactly\DIFdelbegin \DIFdel{from the cases}\DIFdelend . Their intervals come from 1000 resamples of the held-out alternatives within each geometry, \DIFdelbegin \DIFdel{every resample taking one }\DIFdelend \DIFaddbegin \DIFadd{each with a }\DIFaddend block-bootstrap replicate of the true $q_p$ of \DIFdelbegin \DIFdel{each alternative(batches of 50 parts resampled), so that the uncertainty of the reference enters; the resamples are }\DIFdelend \DIFaddbegin \DIFadd{every alternative, }\DIFaddend common to all rules \DIFdelbegin \DIFdel{, which gives paired intervals for their differences . These }\DIFdelend \DIFaddbegin \DIFadd{so that their differences get paired intervals (ESM Table~}\esm{distci}\DIFadd{); these }\DIFaddend intervals hold the fitted rules fixed\DIFdelbegin \DIFdel{; a }\DIFdelend \DIFaddbegin \DIFadd{. A }\DIFaddend second bootstrap resamples \DIFdelbegin \DIFdel{the alternatives of }\DIFdelend \DIFaddbegin \DIFadd{whole lubrication patterns within }\DIFaddend each geometry, \DIFaddbegin \DIFadd{which keeps every case in its evidence relation, and }\DIFaddend refits every alternative-level rule \DIFdelbegin \DIFdel{on each resample and repeats the evaluation, which adds the variability of the calibration sets}\DIFdelend \DIFaddbegin \DIFadd{(ESM Section~S1)}\DIFaddend .

\subsection{Computation and development record}\label{sec:computation}

The constants of the floor and the requirement grid were \DIFdelbegin \DIFdel{chosen }\DIFdelend \DIFaddbegin \DIFadd{fixed }\DIFaddend before the decision results were seen\DIFdelbegin \DIFdel{and are listed in Appendix~\ref{app:computation}; }\DIFdelend \DIFaddbegin \DIFadd{, as }\DIFaddend the repository history records\DIFdelbegin \DIFdel{their dates}\DIFdelend . The rules were not\DIFdelbegin \DIFdel{all fixed in advance}\DIFdelend : M0--M5 \DIFdelbegin \DIFdel{and the three original calibration scopes }\DIFdelend were defined in a first analysis\DIFdelbegin \DIFdel{, the conformalised Gaussian process M6 was added as a candidate during it, and the rules }\DIFdelend \DIFaddbegin \DIFadd{; }\DIFaddend M0w, Mc, M1n, NN, M2n\DIFdelbegin \DIFdel{and }\DIFdelend \DIFaddbegin \DIFadd{, }\DIFaddend M5n, the grouped schemes, the \DIFaddbegin \DIFadd{exact estimator and the }\DIFaddend noise floor of the Gaussian processes \DIFdelbegin \DIFdel{and the exact estimator of the distances were added in a second round of analysis, after the results of the first }\DIFdelend \DIFaddbegin \DIFadd{were added after the first results }\DIFaddend had been seen\DIFdelbegin \DIFdel{. Every variant tried is reported, and }\DIFdelend \DIFaddbegin \DIFadd{; and a review of that version added M1s, the jackknife+ intervals, }\DIFaddend the \DIFdelbegin \DIFdel{comparisons between rules carry }\DIFdelend \DIFaddbegin \DIFadd{between-run noise floor and the analyses of Sections~\ref{sec:res-drivers} and \ref{sec:res-procedure}. The comparison of rules is therefore exploratory and carries }\DIFaddend the optimism of a selection made on the evaluation data \citep{cawley2010over}\DIFaddbegin \DIFadd{; every variant tried is reported in the paper or the ESM}\DIFaddend . The pipeline reads single files from the public archives \DIFdelbegin \DIFdel{by HTTP range requests, keeps only the extracted features, }\DIFdelend and reproduces every result from the cached feature tables in about four hours on four processor cores.

%% ----------------------------------------------------------------------
\section{Results}\label{sec:results}

\subsection{What the production floor contains}\label{sec:res-floor}

\DIFdelbegin \DIFdel{Table~\ref{tab:floors} gives the floor of each characteristic. Two }\DIFdelend \DIFaddbegin \DIFadd{In 95\,\% of batch pairs, two }\DIFaddend batches of 50 consecutive parts of one alternative differ by up to 0.051\,mm in mid-side draw-in \DIFdelbegin \DIFdel{, 0.046\,mm in corner draw-in, }\DIFdelend \DIFaddbegin \DIFadd{and }\DIFaddend 0.078$^\circ$ in wall angle \DIFdelbegin \DIFdel{, 0.11$^\circ$ in arm angle, 0.011\,mm in cup depth and 0.007\,mm in bottom dome at 95\,\% confidence; resampling whole alternatives instead of batches, which respects the dependence of neighbouring batches under drift, widens the intervals most for the waviness (0.0021--0.0050\,mm)and the arm angle (0.089--0.15$^\circ$)}\DIFdelend \DIFaddbegin \DIFadd{(Table~\ref{tab:floors}), and one floor is 0.66 to 2.2 long-term standard deviations of the parts}\DIFaddend . The floors are set by drift \DIFdelbegin \DIFdel{: once }\DIFdelend \DIFaddbegin \DIFadd{within the run: without }\DIFaddend the time order of the parts \DIFdelbegin \DIFdel{is removed, }\DIFdelend they shrink by a factor of 2.1 \DIFdelbegin \DIFdel{(wall angle) }\DIFdelend to 4.5\DIFdelbegin \DIFdel{(corner draw-in), and over a series the punch temperature rises by a median of 4.3\,K (range }\DIFdelend \DIFaddbegin \DIFadd{, batches one position apart differ by 0.51--0.80 floors and nine positions apart by }\DIFaddend 1.1\DIFdelbegin \DIFdel{--6.4}\DIFdelend \DIFaddbegin \DIFadd{--1.5 floors, and the normal model of Section~\ref{sec:floor} reproduces the floors to within 12}\DIFaddend \,\DIFdelbegin \DIFdel{K), one of the sources of this drift.
}%DIFDELCMD < 

%DIFDELCMD < %%%
\DIFdel{Table~\ref{tab:components} shows what the floor contains. Part-to-part noisethat is independent between parts---which }\DIFdelend \DIFaddbegin \DIFadd{\% (ESM Tables~}\esm{protocol} \DIFadd{and }\esm{lags}\DIFadd{). Independent part-to-part noise, which }\DIFaddend includes the repeatability of the \DIFdelbegin \DIFdel{scanner---would }\DIFdelend \DIFaddbegin \DIFadd{scanner, would }\DIFaddend produce 15--33\,\% of each floor, \DIFdelbegin \DIFdel{so the floor is dominated by variation between batches, which a repeatable measurement cannot create. The redundancy of the scans had shown where the measurement did matter: with the flange widths along the scan lines and the north wall included, }\DIFdelend \DIFaddbegin \DIFadd{and the resolution of the measurement is small against the floor except for }\DIFaddend the mid-side draw-in \DIFdelbegin \DIFdel{had a floor of 0.085\,mm and a part standard deviation of 0.13\,mm, of which the redefinition removes 40\,\% and 70\,\%; the wall angle kept its floor (0.081 against 0.078$^\circ$)because the north wall's scatter averaged out within batches. Drift }\DIFdelend \DIFaddbegin \DIFadd{(ESM Tables~}\esm{components} \DIFadd{and }\esm{resolution}\DIFadd{); assumption A3 holds for repeatability and resolution, but the stability }\DIFaddend of the scanner \DIFdelbegin \DIFdel{itself cannot be separated from drift of the process }\DIFdelend \DIFaddbegin \DIFadd{cannot be checked }\DIFaddend without a reference artefact\DIFdelbegin \DIFdel{scanned through the series. The floor also depends on the protocol: from the first 250 parts of every series it is 0.49--0.95 of the full value, and it varies by a factor of 0.6--2.2 over batch sizes of 25--100 and quantiles of 0.90--0.99 (Section~\ref{sec:res-robust})}\DIFdelend \DIFaddbegin \DIFadd{. Nor do the process signals identify the source of the drift: batch centres of press force, punch temperature, sheet thickness and oil film predict it out of fold with $R^2$ of $-0.23$ to $0.20$, and removing the predicted drift changes the floors by $-16$ to $+12$\,\% (ESM Table~}\esm{inline}\DIFadd{). Variation between runs is not in the floor: the three lubrication series of one geometry and force differ by up to 1.2 (waviness) to 9.8 floors (dome), and a floor taken across these series is 1.5 to 9.1 times the reference}\DIFaddend . The decided statistic \DIFdelbegin \DIFdel{itself }\DIFdelend reproduces to about one floor: the 95th percentiles of two batches of 100 parts differ by 0.92--1.49 floors \DIFdelbegin \DIFdel{at }\DIFdelend \DIFaddbegin \DIFadd{in }\DIFaddend 95\,\% \DIFdelbegin \DIFdel{confidence, and the block-bootstrap standard error of an alternative's $q_{95}$ is 0.11 floors in the median (at most 0.39)}\DIFdelend \DIFaddbegin \DIFadd{of pairs}\DIFaddend . Distances below about one floor are therefore \DIFdelbegin \DIFdel{at the limit of what productioncan confirm. Variation between runs is not in the floor: the three lubrication series of one geometry and force, which production mostly cannot tell apart, differ by a median of 0.5--1.6 floors and, at 95\,\%, by 1.2 (waviness) to 9.8 floors (dome), an upper bound of what separate runs of one design add}\DIFdelend \DIFaddbegin \DIFadd{beyond what production, or a trial, can confirm}\DIFaddend .

%% BEGIN GENERATED floors
\DIFdelbegin %DIFDELCMD < \begin{table}[t]
%DIFDELCMD < \tabsize
%DIFDELCMD < \setlength{\tabcolsep}{2.6pt}
%DIFDELCMD < \caption{Production floor per characteristic, pooled and per geometry}\label{tab:floors}
%DIFDELCMD < \begin{tabular}{@{}lrrrrrrr@{}}
%DIFDELCMD < \toprule
%DIFDELCMD <  & & \multicolumn{2}{c}{95\,\% interval} & & & & \\
%DIFDELCMD < \cmidrule(lr){3-4}
%DIFDELCMD < Characteristic & $F$ & batches & alternatives & Concave & Convex & Part SD & Drift \\
%DIFDELCMD < \midrule
%DIFDELCMD < Draw-in, mid-side (mm) & 0.051 & 0.045--0.061 & 0.042--0.061 & 0.050 & 0.052 & 0.039 & 3.3 \\
%DIFDELCMD < Draw-in, corner (mm) & 0.046 & 0.035--0.052 & 0.032--0.053 & 0.053 & 0.030 & 0.027 & 4.5 \\
%DIFDELCMD < Flange waviness (mm) & 0.0028 & 0.0024--0.0031 & 0.0021--0.0050 & 0.0037 & 0.0021 & 0.0020 & 2.5 \\
%DIFDELCMD < Wall angle ($^\circ$) & 0.078 & 0.066--0.10 & 0.054--0.10 & 0.081 & 0.078 & 0.083 & 2.1 \\
%DIFDELCMD < Arm angle, cut ($^\circ$) & 0.11 & 0.094--0.12 & 0.089--0.15 & 0.12 & 0.100 & 0.091 & 2.6 \\
%DIFDELCMD < Cup depth (mm) & 0.011 & 0.0096--0.012 & 0.0090--0.012 & 0.0087 & 0.012 & 0.0069 & 3.9 \\
%DIFDELCMD < Bottom dome, cut (mm) & 0.0070 & 0.0061--0.0081 & 0.0052--0.0084 & 0.0034 & 0.0082 & 0.0036 & 3.9 \\
%DIFDELCMD < \botrule
%DIFDELCMD < \end{tabular}
%DIFDELCMD < \footnotetext{Batches of 50 consecutive parts, 95\,\% quantile of the difference of batch centres. Intervals: 2000 resamples of the batches within the alternatives, and of whole alternatives (clusters). Part SD: median within-alternative standard deviation of single parts. Drift factor: floor over the floor of randomly permuted series. The decision analysis uses the per-geometry floors.}
%DIFDELCMD < \end{table}
%DIFDELCMD < %%%
\DIFdelend \DIFaddbegin \begin{table}[t]
\tabsize
\setlength{\tabcolsep}{1.8pt}
\caption{Production floor per characteristic, pooled and per geometry}\label{tab:floors}
\begin{tabular}{@{}lrrrrrrrr@{}}
\toprule
 & & \multicolumn{2}{c}{95\,\% interval} & & & & & \\
\cmidrule(lr){3-4}
Characteristic & $F$ & batches & series & Concave & Convex & $F/\sigma_\mathrm{LT}$ & Drift & Step \\
\midrule
Draw-in, mid-side (mm) & 0.051 & 0.045--0.061 & 0.042--0.061 & 0.050 & 0.052 & 1.2/1.5 & 3.3 & 0.38 \\
Draw-in, corner (mm) & 0.046 & 0.035--0.052 & 0.032--0.053 & 0.053 & 0.030 & 1.7/1.3 & 4.5 & ${<}$0.01 \\
Flange waviness (mm) & 0.0028 & 0.0024--0.0031 & 0.0021--0.0050 & 0.0037 & 0.0021 & 2.1/0.8 & 2.5 & ${<}$0.01 \\
Wall angle ($^\circ$) & 0.078 & 0.066--0.10 & 0.054--0.10 & 0.081 & 0.078 & 1.6/0.7 & 2.1 & ${<}$0.01 \\
Arm angle, cut ($^\circ$) & 0.11 & 0.094--0.12 & 0.089--0.15 & 0.12 & 0.100 & 1.9/0.9 & 2.6 & ${<}$0.01 \\
Cup depth (mm) & 0.011 & 0.0096--0.012 & 0.0090--0.012 & 0.0087 & 0.012 & 2.2/1.7 & 3.9 & ${<}$0.01 \\
Bottom dome, cut (mm) & 0.0070 & 0.0061--0.0081 & 0.0052--0.0084 & 0.0034 & 0.0082 & 1.4/1.5 & 3.9 & ${<}$0.01 \\
\botrule
\end{tabular}
\footnotetext{Batches of 50 consecutive parts, 95\,\% quantile of the difference of batch centres. Intervals: 2000 resamples of batches within series, and of whole series. $F/\sigma_\mathrm{LT}$: floor over the median long-term standard deviation of the parts (concave/convex). Drift: floor over the floor of randomly permuted series. Step: measurement resolution in floors.}
\end{table}
\DIFaddend %% END GENERATED floors

%DIF < % BEGIN GENERATED components
\DIFdelbegin %DIFDELCMD < \begin{table}[t]
%DIFDELCMD < \tabsize
%DIFDELCMD < \setlength{\tabcolsep}{4.5pt}
%DIFDELCMD < \caption{What the production floor contains: components and variants in floors}\label{tab:components}
%DIFDELCMD < \begin{tabular}{@{}lrrrrrr@{}}
%DIFDELCMD < \toprule
%DIFDELCMD <  & & & & & \multicolumn{2}{c}{Between series} \\
%DIFDELCMD < \cmidrule(lr){6-7}
%DIFDELCMD < Characteristic & $\sigma_\mathrm{LT}/\sigma_\mathrm{ST}$ & White share (\%) & First half & $q_{95}$ floor & median & 95\,\% \\
%DIFDELCMD < \midrule
%DIFDELCMD < Draw-in, mid-side & 1.22 & 24 & 0.93 & 1.49 & 0.9 & 2.4 \\
%DIFDELCMD < Draw-in, corner & 1.24 & 17 & 0.63 & 1.08 & 0.7 & 2.4 \\
%DIFDELCMD < Flange waviness & 1.16 & 23 & 0.81 & 1.07 & 0.5 & 1.2 \\
%DIFDELCMD < Wall angle & 1.23 & 33 & 0.90 & 0.92 & 0.5 & 2.7 \\
%DIFDELCMD < Arm angle, cut & 1.35 & 23 & 0.71 & 0.96 & 1.0 & 4.3 \\
%DIFDELCMD < Cup depth & 1.41 & 15 & 0.95 & 1.11 & 0.5 & 4.8 \\
%DIFDELCMD < Bottom dome, cut & 1.10 & 17 & 0.49 & 1.02 & 1.6 & 9.8 \\
%DIFDELCMD < \botrule
%DIFDELCMD < \end{tabular}
%DIFDELCMD < \footnotetext{$\sigma_\mathrm{ST}$: standard deviation within subgroups of five consecutive parts; $\sigma_\mathrm{LT}$: of the whole series (medians over series). White share: the floor that independent part-to-part noise alone would produce, $1.96\sqrt{2/B}$ times the successive-difference standard deviation, over $F$; it bounds the share a scanner's repeatability can explain. First half: floor of the first 250 parts of each series. $q_{95}$ floor: 95\,\% quantile of the difference between the 95th percentiles of two batches of 100 parts. Between series: difference between the centres of the three lubrication series of one geometry and force, an upper bound of the variation between runs (lubrication effects included).}
%DIFDELCMD < \end{table}
%DIFDELCMD < %%%
%DIF < % END GENERATED components
%DIFDELCMD < 

%DIFDELCMD < %%%
\DIFdelend \subsection{What production resolves \DIFaddbegin \DIFadd{and how far the simulation misses it}\DIFaddend }\DIFdelbegin %DIFDELCMD < \label{sec:res-resolve}
%DIFDELCMD < %%%
\DIFdelend \DIFaddbegin \label{sec:res-sim}
\DIFaddend 

\DIFdelbegin \DIFdel{Figure~\ref{fig:alternatives} shows the 18 alternatives against the floor and Table~\ref{tab:resolve} counts the contrasts that production resolves. }\DIFdelend All 63 geometry contrasts exceed the floor \DIFdelbegin \DIFdel{, with median effect-to-scatter ratios from 2.8 (waviness) to 129 (wall angle). Blank-holder }\DIFdelend \DIFaddbegin \DIFadd{and blank-holder }\DIFaddend force is resolved in 105 of 126 contrasts\DIFdelbegin \DIFdel{; the exception is the wall angle, which force barely moves (4 of 18). The lubrication pattern is resolved }\DIFdelend \DIFaddbegin \DIFadd{, but the lubrication pattern only }\DIFaddend in 51 of 126 \DIFdelbegin \DIFdel{contrasts and in at most 11 of 18 for any characteristic, so that for most characteristics an alternative's }\DIFdelend \DIFaddbegin \DIFadd{(42 with a bootstrap probability of at least 0.95), so a }\DIFaddend lubrication variant is \DIFaddbegin \DIFadd{often }\DIFaddend a near-replicate of \DIFdelbegin \DIFdel{it. The force sensitivity is local: between 100 and 300\,kN, where the stroke speed changes as well, the minimum resolvable }\DIFdelend \DIFaddbegin \DIFadd{its siblings (ESM Table~}\esm{resolve}\DIFadd{). The minimum resolvable force }\DIFaddend change is 21--67\,kN \DIFdelbegin \DIFdel{for the cup depth, the waviness, the draw-ins and the arm angle, whereas between 300 and 500\,kN, at equal speed, it is 43--105}\DIFdelend \DIFaddbegin \DIFadd{between 100 and 300}\DIFaddend \,kN for \DIFdelbegin \DIFdel{the depth, the draw-ins and the waviness and the arm angle no longer responds; }\DIFdelend \DIFaddbegin \DIFadd{most characteristics and lies beyond either force interval for }\DIFaddend the wall angle and the dome\DIFdelbegin \DIFdel{need 366--731\,kN in either interval. }%DIFDELCMD < 

%DIFDELCMD < %%%
%DIF < % BEGIN GENERATED resolve
%DIFDELCMD < \begin{table}[t]
%DIFDELCMD < \tabsize
%DIFDELCMD < \setlength{\tabcolsep}{3.5pt}
%DIFDELCMD < \caption{Single-factor contrasts that production resolves and minimum resolvable force change}\label{tab:resolve}
%DIFDELCMD < \begin{tabular}{@{}lrrrrrrrr@{}}
%DIFDELCMD < \toprule
%DIFDELCMD <  & \multicolumn{2}{c}{Geometry} & \multicolumn{2}{c}{Blank-holder force} & \multicolumn{2}{c}{Lubrication} & \multicolumn{2}{c}{Force MRC (kN)} \\
%DIFDELCMD < \cmidrule(lr){2-3}\cmidrule(lr){4-5}\cmidrule(lr){6-7}\cmidrule(lr){8-9}
%DIFDELCMD < Characteristic & ESR\,$\geq$\,1 & Median & ESR\,$\geq$\,1 & Median & ESR\,$\geq$\,1 & Median & 100--300 & 300--500 \\
%DIFDELCMD < \midrule
%DIFDELCMD < Draw-in, mid-side & 9/9 & 18.7 & 18/18 & 6.8 & 9/18 & 0.9 & 32 & 105 \\
%DIFDELCMD < Draw-in, corner & 9/9 & 108.8 & 18/18 & 4.3 & 7/18 & 0.7 & 67 & 47 \\
%DIFDELCMD < Flange waviness & 9/9 & 2.8 & 15/18 & 6.4 & 3/18 & 0.5 & 21 & 78 \\
%DIFDELCMD < Wall angle & 9/9 & 129.0 & 4/18 & 0.6 & 5/18 & 0.5 & 420 & 366 \\
%DIFDELCMD < Arm angle, cut & 9/9 & 47.2 & 17/18 & 3.0 & 9/18 & 1.0 & 48 & 1083 \\
%DIFDELCMD < Cup depth & 9/9 & 34.6 & 18/18 & 10.2 & 7/18 & 0.5 & 21 & 43 \\
%DIFDELCMD < Bottom dome, cut & 9/9 & 9.6 & 15/18 & 2.4 & 11/18 & 1.6 & 543 & 731 \\
%DIFDELCMD < \midrule All & 63/63 & 34.6 & 105/126 & 4.3 & 51/126 & 0.7 & & \\
%DIFDELCMD < \botrule
%DIFDELCMD < \end{tabular}
%DIFDELCMD < \footnotetext{ESR: effect-to-scatter ratio of two alternatives that differ in one factor; counts of contrasts with ESR\,$\geq$\,1 and median ESR. MRC: floor over the sensitivity of the alternative centres to blank-holder force, between 100 and 300\,kN (the stroke speed also changes) and between 300 and 500\,kN (same speed); the sensitivity is local.}
%DIFDELCMD < \end{table}
%DIFDELCMD < %%%
%DIF < % END GENERATED resolve
%DIFDELCMD < 

%DIFDELCMD < \begin{figure}[p]
%DIFDELCMD < \centering
%DIFDELCMD < \includegraphics[width=\textwidth]{F1_alternatives}
%DIFDELCMD < \caption{Centre and 5--95\,\% range of each characteristic per alternative; the black bar is the production floor}\label{fig:alternatives}
%DIFDELCMD < \end{figure}
%DIFDELCMD < 

%DIFDELCMD < %%%
\subsection{\DIFdel{Simulation against production}}%DIFAUXCMD
\addtocounter{subsection}{-1}%DIFAUXCMD
%DIFDELCMD < \label{sec:res-sim}
%DIFDELCMD < 

%DIFDELCMD < %%%
\DIFdelend \DIFaddbegin \DIFadd{. }\DIFaddend The nominal simulation misses production by far more than a floor\DIFdelbegin \DIFdel{(Table~\ref{tab:offsets}). Its }\DIFdelend \DIFaddbegin \DIFadd{: its }\DIFaddend median offset reaches 109 floors\DIFdelbegin \DIFdel{for the cup depth of the concave cup and 108 for the corner draw-in of the convex cup, and within one geometry the offset varies by 8.5 to 39 floors between alternatives. In units, the simulation overestimates the corner draw-in by about 3\,mm and the mid-side draw-in by 0.6--1.2\,mm and underestimates the wall angle by 1.0--1.2$^\circ$. Its }\DIFdelend \DIFaddbegin \DIFadd{, its }\DIFaddend blank-holder force effect on the draw-in is \DIFdelbegin \DIFdel{2.6 }\DIFdelend \DIFaddbegin \DIFadd{2.7 }\DIFaddend to 5.1 times the measured one, \DIFdelbegin \DIFdel{which points to the friction model---a constant Coulomb coefficient where friction falls with contact pressure \mbox{%DIFAUXCMD
\citep{hol2012advanced}}\hskip0pt%DIFAUXCMD
---and to the modelled blank-holder contact rather than to the solver. The signed comparison exposes a sign error that the decision scale hides: }\DIFdelend \DIFaddbegin \DIFadd{and }\DIFaddend the cut arms of the concave cups bend \DIFdelbegin \DIFdel{upwards in every part and }\DIFdelend downwards in 89--100\,\% of the simulations \DIFdelbegin \DIFdel{, so their absolute values agree better than their physics. The offset of the cup depth, $-0.90$ to $-0.95$}\DIFdelend \DIFaddbegin \DIFadd{but upwards in every part (ESM Tables~}\esm{offsets} \DIFadd{and }\esm{simnoise}\DIFadd{). Part of the offset is definitional: the cup depth differs by a median 0.90 and 0.95}\DIFaddend \,mm \DIFdelbegin \DIFdel{for all alternatives, is nearly constant and is consistent with a difference of reference---the shell mid-surface against }\DIFdelend \DIFaddbegin \DIFadd{in the two geometries, consistent with }\DIFaddend the scanned surface \DIFdelbegin \DIFdel{, and the drawing depth of the model---rather than with forming physics. The simulated characteristics also carry extraction noise from the 1\,mm node spacing (Table~\ref{tab:simnoise}): 1.7--3.1 floors for the wall angle and 2.4--9.8 floors for the dome, which inflates the upper end of the simulation envelope of the dome by up to 21 }\DIFdelend \DIFaddbegin \DIFadd{against the shell mid-surface, and removing this offset alone brings the decisive distance of the nominal simulation from 108 to 65 }\DIFaddend floors.

%DIF < % BEGIN GENERATED offsets
\DIFdelbegin %DIFDELCMD < \begin{table}[t]
%DIFDELCMD < \tabsize
%DIFDELCMD < \setlength{\tabcolsep}{5pt}
%DIFDELCMD < \caption{Offset between the parts' 95th percentile and the nominal simulation, in floors}\label{tab:offsets}
%DIFDELCMD < \begin{tabular}{@{}lrrrrrr@{}}
%DIFDELCMD < \toprule
%DIFDELCMD <  & \multicolumn{2}{c}{Concave} & \multicolumn{2}{c}{Convex} & \multicolumn{2}{c}{Matched, signed} \\
%DIFDELCMD < \cmidrule(lr){2-3}\cmidrule(lr){4-5}\cmidrule(lr){6-7}
%DIFDELCMD < Characteristic & Offset & Spread & Offset & Spread & Concave & Convex \\
%DIFDELCMD < \midrule
%DIFDELCMD < Draw-in, mid-side (mm) & $-$16 & 37 & $-$20 & 32 & $-$0.60 & $-$1.19 \\
%DIFDELCMD < Draw-in, corner (mm) & $-$59 & 17 & $-$108 & 22 & $-$2.96 & $-$3.32 \\
%DIFDELCMD < Flange waviness (mm) & 7 & 11 & 2 & 13 & +0.03 & $-$0.01 \\
%DIFDELCMD < Wall angle ($^\circ$) & 16 & 19 & 12 & 13 & +1.20 & +1.03 \\
%DIFDELCMD < Arm angle, cut ($^\circ$) & $-$27 & 39 & 1 & 15 & +4.26 & $-$0.04 \\
%DIFDELCMD < Cup depth (mm) & $-$109 & 14 & $-$74 & 31 & $-$0.95 & $-$0.90 \\
%DIFDELCMD < Bottom dome, cut (mm) & $-$28 & 22 & 10 & 9 & $-$0.04 & +0.06 \\
%DIFDELCMD < \botrule
%DIFDELCMD < \end{tabular}
%DIFDELCMD < \footnotetext{Offset: median over the nine alternatives of a geometry of $(q_{95}-s_0)/F$, $s_0$ the nominal simulation, on the decision scale (absolute values for the arm angle and the dome); spread: its range over the alternatives. Matched, signed: median difference between the centre of the parts and of their matched simulations, in the characteristic's unit, with signs; the simulated arm angle of the concave cup has the opposite sign of the parts'.}
%DIFDELCMD < \end{table}
%DIFDELCMD < %%%
%DIF < % END GENERATED offsets
%DIFDELCMD < 

%DIFDELCMD < %%%
\DIFdelend \subsection{Decisions by evidence relation}\label{sec:res-decisions}

Table~\ref{tab:distances} and Fig.~\ref{fig:decisions} score the \DIFdelbegin \DIFdel{twelve }\DIFdelend rules. The nominal simulation \DIFdelbegin \DIFdel{(M0) }\DIFdelend needs requirements 108 floors from the truth (bootstrap interval 101--115) before 95\,\% of its verdicts are right\DIFdelbegin \DIFdel{, and }\DIFdelend \DIFaddbegin \DIFadd{; }\DIFaddend it errs mostly by rejecting adequate alternatives\DIFdelbegin \DIFdel{: its safe distance is }\DIFdelend \DIFaddbegin \DIFadd{, with a safe distance of }\DIFaddend 17 floors for false accepts but 109 for false rejects\DIFdelbegin \DIFdel{(Table~\ref{tab:onesided}). Reading the process-window simulation by its worst case (M0w) makes it neither safe nor decisive (94 and 119 floors), because the window does not contain the simulation's bias}\DIFdelend .

\emph{A new variant with a produced sibling.} With the other eight alternatives of the family produced, every calibrated rule improves \DIFdelbegin \DIFdel{by an order of magnitude}\DIFdelend \DIFaddbegin \DIFadd{on the nominal simulation}\DIFaddend , and the simplest \DIFdelbegin \DIFdel{one }\DIFdelend is the most decisive: the nearest produced setting\DIFdelbegin \DIFdel{(NN)}\DIFdelend , the mean $q_{95}$ of the two siblings at the same \DIFdelbegin \DIFdel{blank-holder }\DIFdelend force, is right beyond 3.4 floors (2.0--4.5)\DIFdelbegin \DIFdel{and at 10 floors decides every case correctly. The Gaussian process without the simulation (M5n) decides at 6.1 floorsand is safe at 1.0, and with the simulation (M5) at 6.7 and 1.4. The interval rules are safe at or below 1.4 floors, but the envelope with the residual margin (M2) }\DIFdelend \DIFaddbegin \DIFadd{, 0.17\,mm of mid-side draw-in. The scaled simulation follows at 4.9 floors, the Gaussian processes with and without the simulation decide at 7.2 floors and are safe below one floor, and the envelope rule is safe at 0.15 floors but }\DIFaddend sends a third of the designs to trial at 10 floors and \DIFdelbegin \DIFdel{is decisive }\DIFdelend \DIFaddbegin \DIFadd{decides }\DIFaddend only at 23 floors. Bias-correcting the nominal simulation (M1, 15 floors) is worse than taking the median of the produced alternatives without it (M1n, 9.1)\DIFdelbegin \DIFdel{, and calibrating the friction coefficient (Mc, 27 floors) is worse still}\DIFdelend .

\emph{A new process setting.} When every alternative at the new design's \DIFdelbegin \DIFdel{blank-holder }\DIFdelend force is held out, all rules lose. The nearest produced setting is right beyond 8.5 floors (5.5--10.2)\DIFdelbegin \DIFdel{and the median of the produced alternatives beyond 10}\DIFdelend ; the interval rules \DIFdelbegin \DIFdel{are decisive only at 13--18 floors and safe at 5.4--8.1 (M4 13/7.9, M2n 15/}\DIFdelend \DIFaddbegin \DIFadd{decide only at 13--31 floors and are safe at }\DIFaddend 5.9\DIFdelbegin \DIFdel{, M2 18/8.1, M5 18/7.4, M5n 18/5.4; M3 32/16)}\DIFdelend \DIFaddbegin \DIFadd{--16}\DIFaddend , and at 10 floors they send \DIFdelbegin \DIFdel{a fifth to a third }\DIFdelend \DIFaddbegin \DIFadd{18--37\,\% }\DIFaddend of the designs to trial. Interpolating in force (300\,kN from 100 and 500\,kN) is much easier than extrapolating: the nearest setting decides at 4.7 against 9.5 floors, and M2 and M5 are safe at 2.7 and \DIFdelbegin \DIFdel{2.9 }\DIFdelend \DIFaddbegin \DIFadd{2.8 }\DIFaddend against 9.9 and 8.9 floors. \DIFaddbegin \DIFadd{The scaled simulation decides at 19 floors and, extrapolating a fitted trend, falsely accepts failing alternatives as far as 164 floors from their limit.
}\DIFaddend 

\emph{A new design family.} \DIFdelbegin \DIFdel{Calibrated only on the other geometry}\DIFdelend \DIFaddbegin \DIFadd{In the two transfers}\DIFaddend , every rule without the simulation fails (131--133 floors), because nothing in one family's production record describes the other. With the simulation, the bias correction decides at 35 floors, M5 at 38 and \DIFaddbegin \DIFadd{the boosting model }\DIFaddend M3 at \DIFdelbegin \DIFdel{48}\DIFdelend \DIFaddbegin \DIFadd{49}\DIFaddend , and the envelope \DIFdelbegin \DIFdel{with its residual margin }\DIFdelend \DIFaddbegin \DIFadd{rule }\DIFaddend is the safest \DIFdelbegin \DIFdel{rule }\DIFdelend at 16 floors; no rule is safe closer than that. \DIFdelbegin \DIFdel{The simulation thus reduces the distance at which a new familycan be decided from about 130 floors to 35--51, but not to a distance at }\DIFdelend \DIFaddbegin \DIFadd{Scored in the floor of the produced family, }\DIFaddend which a design team would \DIFdelbegin \DIFdel{rely on it near the requirement}\DIFdelend \DIFaddbegin \DIFadd{have, M5 decides at 35 and is safe at 31 floors and the envelope rule at 41 and 26. Sixteen floors are about 0.8\,mm of mid-side draw-in or 1.3$^\circ$ of wall angle; ESM Tables~}\esm{qcdist} \DIFadd{and }\esm{units} \DIFadd{give the distances per characteristic in floors and in units}\DIFaddend .

%% BEGIN GENERATED distances
\DIFdelbegin %DIFDELCMD < \begin{table}[t]
%DIFDELCMD < \tabsize
%DIFDELCMD < \setlength{\tabcolsep}{3.5pt}
%DIFDELCMD < \caption{Decisive and safe distances (floors) by the relation of the new design to the produced evidence}\label{tab:distances}
%DIFDELCMD < \begin{tabular}{@{}lrrrrr@{}}
%DIFDELCMD < \toprule
%DIFDELCMD <  & New & \multicolumn{3}{c}{New process setting} & New \\
%DIFDELCMD < \cmidrule(lr){3-5}
%DIFDELCMD < Rule & variant & all & interpolation & extrapolation & family \\
%DIFDELCMD < \midrule
%DIFDELCMD < \multicolumn{6}{@{}l}{\emph{Simulation only}} \\
%DIFDELCMD < M0 nominal simulation & 108 & 108 & 109 & 106 & 108 \\
%DIFDELCMD < M0w process-window worst case & 119/94 & 119/94 & 119/94 & 119/77 & 119/94 \\
%DIFDELCMD < \multicolumn{6}{@{}l}{\emph{Simulation and production record}} \\
%DIFDELCMD < M1 bias-corrected & 15 & 24 & 11 & 25 & 35 \\
%DIFDELCMD < Mc friction-calibrated & 27 & 39 & 14 & 40 & 38 \\
%DIFDELCMD < M2 envelope + margin & 23/0.2 & 18/8.1 & 18/2.7 & 18/9.9 & 51/16 \\
%DIFDELCMD < M3 simulation + boosting & 17/0.0 & 32/16 & 31/15 & 32/16 & 48/22 \\
%DIFDELCMD < M5 GP correction & 6.7/1.4 & 18/7.4 & 19/2.9 & 18/8.9 & 38/30 \\
%DIFDELCMD < \multicolumn{6}{@{}l}{\emph{Production record only}} \\
%DIFDELCMD < M1n median & 9.1 & 10 & 4.8 & 12 & 131 \\
%DIFDELCMD < NN nearest setting & 3.4 & 8.5 & 4.7 & 9.5 & 131 \\
%DIFDELCMD < M2n median + margin & 19/0.0 & 15/5.9 & 13/0.0 & 16/9.1 & 132/129 \\
%DIFDELCMD < M4 boosting & 11/0.6 & 13/7.9 & 15/8.5 & 13/7.6 & 133/128 \\
%DIFDELCMD < M5n GP of $q_{95}$ & 6.1/1.0 & 18/5.4 & 16/1.3 & 18/8.5 & 131/131 \\
%DIFDELCMD < \botrule
%DIFDELCMD < \end{tabular}
%DIFDELCMD < \footnotetext{Decisive/safe distance $\delta_\mathrm{d}/\delta_\mathrm{s}$ over all seven characteristics; one number for rules that always decide ($\delta_\mathrm{d}=\delta_\mathrm{s}$). New variant: a sibling at the same blank-holder force is produced (leave one alternative out within the family); new setting: all alternatives of one force are held out (interpolation at 300\,kN, extrapolation at 100 or 500\,kN); new family: calibration on the other geometry. 126 cases per rule (42 for interpolation, 84 for extrapolation). Intervals in Table~\ref{tab:distci}.}
%DIFDELCMD < \end{table}
%DIFDELCMD < %%%
\DIFdelend \DIFaddbegin \begin{table}[t]
\tabsize
\setlength{\tabcolsep}{3pt}
\caption{Decisive and safe distances (floors) by the relation of the new design to the produced evidence}\label{tab:distances}
\begin{tabular}{@{}lrrrrrrrr@{}}
\toprule
 & \multicolumn{2}{c}{New variant} & \multicolumn{4}{c}{New process setting} & \multicolumn{2}{c}{New family} \\
\cmidrule(lr){2-3}\cmidrule(lr){4-7}\cmidrule(lr){8-9}
Rule & $\delta_\mathrm{d}/\delta_\mathrm{s}$ & FA & $\delta_\mathrm{d}/\delta_\mathrm{s}$ & FA & interpolated & extrapolated & $\delta_\mathrm{d}/\delta_\mathrm{s}$ & FA \\
\midrule
M0 & 108 & 17 & 108 & 17 & 109 & 106 & 108 & 17 \\
M1 & 15 & 15 & 24 & 24 & 11 & 25 & 35 & 33 \\
M1s & 4.9 & 4.7 & 19 & 164 & 19 & 19 & 122 & 122 \\
\quad interval & 4.7--6.3 &  & 14--21 &  & 19--20 & 15--19 & 121--123 & \\
M2 & 23/0.15 & 0.30 & 18/8.1 & 8.3 & 18/2.7 & 18/9.9 & 51/16 & 11 \\
M5 & 7.2/0.75 & 0.15 & 19/6.8 & 7.6 & 19/2.8 & 19/8.9 & 38/28 & 32 \\
\quad interval & 6.6--8.4 &  & 18--20 &  & 18--20 & 18--20 & 35--39 & \\
M5n & 7.2/0.55 & 0.10 & 17/6.1 & 9.0 & 16/0 & 18/8.6 & 132/130 & 167 \\
\quad interval & 6.4--7.5 &  & 16--18 &  & 15--16 & 16--19 & 130--133 & \\
NN & 3.4 & 2.8 & 8.5 & 10 & 4.7 & 9.5 & 131 & 169 \\
\quad interval & 2.0--4.5 &  & 5.5--10 &  & 4.3--5.1 & 6.6--11 & 129--131 & \\
\botrule
\end{tabular}
\footnotetext{Over all seven characteristics; one number for rules that always decide. FA: safe distance on the false-accept side (failing alternatives). Interval: bootstrap 95\,\% interval of the decisive distance (fits fixed). 126 cases per relation (42 interpolated, 84 extrapolated); a new family rests on two transfers. All rules: ESM Table~\esm{alldist}.}
\end{table}
\DIFaddend %% END GENERATED distances

\DIFdelbegin %DIFDELCMD < \begin{figure}[t]
%DIFDELCMD < \centering
%DIFDELCMD < \includegraphics[width=\textwidth]{F3_decisions_compare}
%DIFDELCMD < \caption{Verdict shares against the requirement distance $d$ (symmetric logarithmic axis, linear within $\pm10$ floors; 126 cases per point) for four rules and three evidence relations}\label{fig:decisions}
%DIFDELCMD < \end{figure}
%DIFDELCMD < %%%
\DIFdelend \DIFaddbegin \begin{figure}[t]
\centering
\includegraphics[width=\textwidth]{F3_decisions_compare}
\caption{Verdict shares against the requirement distance $d$ (symmetric logarithmic axis, linear within $\pm$10 floors; 126 cases per panel) with the decisive and safe distances marked}\label{fig:decisions}
\end{figure}

\subsection{\DIFadd{Where the distances come from}}\label{sec:res-drivers}

\DIFadd{Three features of the evaluation design shape the distances (ESM Table~}\esm{drivers}\DIFadd{). First, the sibling relation mixes replication with variation. In 58 of the 126 new-variant cases neither sibling differs resolvably from the held-out alternative, and the nearest produced setting is then right beyond 0.85 floors, the reproducibility of $q_{95}$ between near-replicate runs; when both siblings differ resolvably (34 cases), it needs 4.9 floors and M5 9.7. Second, the extrapolation cases are confounded with stroke speed: the production-only rules lose almost only at 100\,kN, where the press also ran slower (the nearest setting decides at 11 floors against 4.9 at 500\,kN), whereas the simulation-based interval rules are unsafe at both ends. Third, the production record of the other family does not help within a family: pooling it makes most rules less decisive (M5 from 7.2 to 11 floors for a new variant, M2 from 23 to 53, M2n from 19 to 325), although it raises the coverage of M5 from 86 to 90\,\% and makes the centre of its interval the most accurate point prediction within the family (2.9 floors against 3.4 for the nearest setting).
}

\DIFadd{The relation therefore does not govern decision fitness in general. Over the eleven calibrated rules and all three relations, the relation explains 69\,\% of the variation of the log decisive distance and the rule 6\,\%, but this share comes from the new family, where the rules without the simulation fail by construction; within a family the rule explains 68\,\% and the relation 15\,\%. The best attainable distance nonetheless grows with the novelty of the design, from 3.4 floors with a sibling to 8.5 for a new setting and 35 for a new family.
}\DIFaddend 

\subsection{What the simulation and learning contribute}\label{sec:res-ml}

Pairing each calibrated rule with its counterpart without the simulation isolates what the simulation adds \DIFdelbegin \DIFdel{(Table~\ref{tab:paired}). Within a family }\DIFdelend \DIFaddbegin \DIFadd{as it is used (ESM Table~}\esm{paired}\DIFadd{). Added as an offset, }\DIFaddend it adds nothing \DIFaddbegin \DIFadd{within a family and degrades two rules}\DIFaddend : the decisive distance of M5 \DIFdelbegin \DIFdel{differs from }\DIFdelend \DIFaddbegin \DIFadd{equals }\DIFaddend that of M5n \DIFdelbegin \DIFdel{by $+0.6$ floors (95\,\% interval }\DIFdelend \DIFaddbegin \DIFadd{(difference }\DIFaddend 0.0 \DIFdelbegin \DIFdel{to $+1.6$), that }\DIFdelend \DIFaddbegin \DIFadd{floors, $-0.4$ to $+1.1$), and those }\DIFaddend of M2 \DIFdelbegin \DIFdel{from }\DIFdelend \DIFaddbegin \DIFadd{and M1 exceed those of }\DIFaddend M2n \DIFdelbegin \DIFdel{by $+3.9$ ($+0.4$ }\DIFdelend \DIFaddbegin \DIFadd{and M1n by 3.9 ($+0.5$ }\DIFaddend to $+7.4$) \DIFdelbegin \DIFdel{, and that of M1 from M1n by $+6.2$ }\DIFdelend \DIFaddbegin \DIFadd{and 6.2 floors }\DIFaddend ($+4.6$ to $+9.7$)\DIFdelbegin \DIFdel{; for a new process setting the bias correction loses 13.5 floors by using the simulation}\DIFdelend , because the simulated force effect is too large. \DIFaddbegin \DIFadd{Rescaled, it adds information within a family, where the scaled simulation decides 4.2 floors closer than the median of the produced alternatives ($-5.4$ to $-0.5$) though 1.5 floors less close than the nearest produced setting, but for a new setting the rescaled trend is worse than none. }\DIFaddend Only for a new family is the simulation decisive, by about \DIFdelbegin \DIFdel{93 floors. Learning adds little over the nearest produced setting: within the family M5 is 3.3 floorsless decisive than NN ($+2.1$ to $+5.1$)and 2 floors safer, and for a new setting 9.4 floors less decisive.
The ranking between learners rests on their }\DIFdelend \DIFaddbegin \DIFadd{94 floors, and even there production knowledge carries over where a drawing nominal transfers: for the wall angle, the design angle plus the produced family's deviation from its own design angle decides at 2.0 floors, against 7.3 for M5 (ESM Table~}\esm{transfer}\DIFadd{). A surrogate of the simulations reproduces them to within 1.8 floors for 12 of the 14 characteristics and geometries, yet its decision fitness is that of the simulation it reproduces.
}

\DIFadd{Learning adds }\DIFaddend intervals rather than \DIFdelbegin \DIFdel{on their predictions: used }\DIFdelend \DIFaddbegin \DIFadd{accuracy. Used }\DIFaddend as point rules, the centres of the M5, M5n and \DIFdelbegin \DIFdel{M4 }\DIFdelend \DIFaddbegin \DIFadd{M3n }\DIFaddend intervals decide at \DIFdelbegin \DIFdel{3.4, 2.8 and 4.8 }\DIFdelend \DIFaddbegin \DIFadd{3.3, 3.1 and 3.7 }\DIFaddend floors within the family, \DIFdelbegin \DIFdel{close to }\DIFdelend \DIFaddbegin \DIFadd{as close as }\DIFaddend the nearest produced setting\DIFdelbegin \DIFdel{(Table~\ref{tab:coverage}). The part-level models keep their ranking under other learners and settings, including quantile-loss boosting and settings selected by an inner leave-one-out (Table~\ref{tab:tuning}).
}%DIFDELCMD < 

%DIFDELCMD < %%%
\DIFdel{The intervals are calibrated only where the calibration alternatives are exchangeable with the new design (Table~\ref{tab:coverage}). Within the family the residual margin of M2 covers the true $q_{95}$ in 88.9\,\% of the cases---exactly the $n/(n+1)=8/9$ that eight residuals attain---and the part-level jackknife of M3 in 93\,\%; }\DIFdelend \DIFaddbegin \DIFadd{, so the ranking of the learned rules rests on the width of their intervals (ESM Table~}\esm{coverage}\DIFadd{). With the noise floor set by the variation between runs, }\DIFaddend the Gaussian-process intervals cover \DIFdelbegin \DIFdel{71--72}\DIFdelend \DIFaddbegin \DIFadd{85--86}\DIFaddend \,\% \DIFdelbegin \DIFdel{although their nominal level is }\DIFdelend \DIFaddbegin \DIFadd{of the true $q_{95}$ within the family against a nominal }\DIFaddend 90\,\%, \DIFdelbegin \DIFdel{and they are narrow (1.1--1.2 floors in half-width) because the fits interpolate: the noise level sits at its lower bound in 35--38}\DIFdelend \DIFaddbegin \DIFadd{the largest-residual margin 89}\DIFaddend \,\% \DIFdelbegin \DIFdel{of the fits and a length scale at a bound in 28--48}\DIFdelend \DIFaddbegin \DIFadd{(it guarantees at least $(n-1)/(n+1)=7/9$ under exchangeability, which these cases lack) and the jackknife+ interval of M3 96}\DIFaddend \,\%\DIFdelbegin \DIFdel{(Table~\ref{tab:gp}). For }\DIFdelend \DIFaddbegin \DIFadd{; for }\DIFaddend a new setting every calibrated interval covers only \DIFdelbegin \DIFdel{39--56}\DIFdelend \DIFaddbegin \DIFadd{43--55}\DIFaddend \,\%, and for a new family the Gaussian processes cover \DIFdelbegin \DIFdel{none, M2 47\,\% and M3 40}\DIFdelend \DIFaddbegin \DIFadd{at most 12}\DIFaddend \,\%. \DIFdelbegin \DIFdel{What a designer observes for a single verdict is the margin between the predicted interval and }\DIFdelend \DIFaddbegin \DIFadd{The inductive biases explain much of this: trees do not extrapolate, two force levels identify at most a linear trend, and in the median new-setting fit the length scale of }\DIFaddend the \DIFdelbegin \DIFdel{requirement. With requirements spread over the 20 floors around the truth, the share of correct verdicts among the decided ones grows with that margin for a new variant and a new setting---for the nearest produced setting from 86\,\% at any margin to 95\,\% beyond 5 floors and 99\,\% beyond 10 for a new setting---but not for a new family, where M5 is right in 72--75\,\% of its decided verdicts whatever the margin. The observed margin is a guide to a verdict only where the rule is calibrated for the evidence relation}\DIFdelend \DIFaddbegin \DIFadd{force sits at its lower bound, so the Gaussian process reverts to its mean (ESM Tables~}\esm{gp} \DIFadd{and }\esm{tuning}\DIFadd{)}\DIFaddend .

\subsection{Calibration budget\DIFdelbegin \DIFdel{and calibration design}\DIFdelend }\label{sec:res-budget}

Figure~\ref{fig:budget} and \DIFdelbegin \DIFdel{Table~\ref{tab:design} }\DIFdelend \DIFaddbegin \DIFadd{ESM Table~}\esm{design} \DIFaddend show how the distances depend on the number of produced alternatives of the family and on how the new design relates to them. A single produced sibling decides: with one alternative at the new design's force, the bias correction, the \DIFdelbegin \DIFdel{residual-margin }\DIFdelend \DIFaddbegin \DIFadd{envelope }\DIFaddend rule and the nearest setting \DIFdelbegin \DIFdel{all }\DIFdelend \DIFaddbegin \DIFadd{coincide by construction (the simulation terms cancel and a single residual gives no margin) and }\DIFaddend decide at 3.5 floors (2.3--4.6) without a wrong verdict at 10 floors\DIFdelbegin \DIFdel{, and further }\DIFdelend \DIFaddbegin \DIFadd{. Further }\DIFaddend alternatives do not improve the nearest setting\DIFdelbegin \DIFdel{. They }\DIFdelend \DIFaddbegin \DIFadd{; they }\DIFaddend make the bias correction worse (12--18 floors), \DIFdelbegin \DIFdel{because the offsets of other forces dilute the sibling's}\DIFdelend \DIFaddbegin \DIFadd{the envelope rule safer but less decisive}\DIFaddend , and the \DIFdelbegin \DIFdel{interval rules safer but not more decisive : M2 becomes safe at 0.2--2.2 floors and decisive at 15--23}\DIFdelend \DIFaddbegin \DIFadd{Gaussian processes safer and more decisive (from 12--14 floors with two alternatives to 7.2 with eight)}\DIFaddend . Without a sibling, bracketing the new force buys safety but not decisiveness: \DIFdelbegin \DIFdel{with produced alternatives onboth sides, M2, M5 and M5n }\DIFdelend \DIFaddbegin \DIFadd{from two alternatives on, the interval rules }\DIFaddend are safe at \DIFdelbegin \DIFdel{0.5--4}\DIFdelend \DIFaddbegin \DIFadd{0--4}\DIFaddend .8 floors \DIFdelbegin \DIFdel{, M2 and M5n are never wrong }\DIFdelend \DIFaddbegin \DIFadd{with no wrong verdict }\DIFaddend at 10 floors\DIFdelbegin \DIFdel{and M5 in under 0.5\,\%, but they }\DIFdelend \DIFaddbegin \DIFadd{, but }\DIFaddend decide only at 15--25 floors, \DIFdelbegin \DIFdel{while the nearest setting decides at }\DIFdelend \DIFaddbegin \DIFadd{against }\DIFaddend 4.7--6.0 \DIFdelbegin \DIFdel{floors}\DIFdelend \DIFaddbegin \DIFadd{for the nearest setting}\DIFaddend . Extrapolation stays unsafe at every budget\DIFdelbegin \DIFdel{: no rule is safe closer than 8.5 floors, and at 10 floors 2--13\,\% of the verdicts of the interval rules and the nearest setting are wrong, and 26--28\,\% of those of the bias correction. Averaged }\DIFdelend \DIFaddbegin \DIFadd{, and averaged }\DIFaddend over all subsets \DIFdelbegin \DIFdel{, }\DIFdelend the distances shrink with the number of produced alternatives mainly because the share of subsets that contain a sibling grows, from 46\,\% \DIFdelbegin \DIFdel{of the cases }\DIFdelend with two alternatives to 96\,\% with six\DIFdelbegin \DIFdel{; what a produced alternative is worth depends on its relation to the designs to come.}%DIFDELCMD < 

%DIFDELCMD < %%%
%DIF < % BEGIN GENERATED design
%DIFDELCMD < \begin{table}[t]
%DIFDELCMD < \tabsize
%DIFDELCMD < \setlength{\tabcolsep}{3pt}
%DIFDELCMD < \caption{Calibration design: distances and wrong verdicts at 10 floors by the relation of the new design to the produced alternatives}\label{tab:design}
%DIFDELCMD < \begin{tabular}{@{}llrrrrrrrr@{}}
%DIFDELCMD < \toprule
%DIFDELCMD < Relation & Rule & \multicolumn{2}{c}{$k=1$} & \multicolumn{2}{c}{$k=2$} & \multicolumn{2}{c}{$k=4$} & \multicolumn{2}{c}{$k=6$} \\
%DIFDELCMD < \cmidrule(lr){3-4}\cmidrule(lr){5-6}\cmidrule(lr){7-8}\cmidrule(lr){9-10}
%DIFDELCMD <  & & $\delta_\mathrm{d}/\delta_\mathrm{s}$ & wrong & $\delta_\mathrm{d}/\delta_\mathrm{s}$ & wrong & $\delta_\mathrm{d}/\delta_\mathrm{s}$ & wrong & $\delta_\mathrm{d}/\delta_\mathrm{s}$ & wrong \\
%DIFDELCMD < \midrule
%DIFDELCMD < sibling & M1 & 3.5 & 0 & 12 & 7 & 17 & 14 & 18 & 15 \\
%DIFDELCMD < sibling & M2 & 3.5/3.5 & 0 & 15/2.2 & 0 & 23/1.0 & 0 & 23/0.5 & 0 \\
%DIFDELCMD < sibling & M5 & -- & -- & 14/1.6 & 0 & 7.3/2.1 & 0 & 7.1/1.6 & 0 \\
%DIFDELCMD < sibling & NN & 3.5 & 0 & 3.5 & 0 & 3.4 & 0 & 3.4 & 0 \\
%DIFDELCMD < interpolation & M1 & -- & -- & 11 & 8 & 13 & 15 & 11 & 9 \\
%DIFDELCMD < interpolation & M2 & -- & -- & 15/4.8 & 0 & 25/3.8 & 0 & 18/2.7 & 0 \\
%DIFDELCMD < interpolation & M5 & -- & -- & 19/3.7 & 0 & 22/3.3 & 0 & 19/2.9 & 0 \\
%DIFDELCMD < interpolation & NN & -- & -- & 5.0 & 0 & 5.3 & 0 & 4.7 & 0 \\
%DIFDELCMD < extrapolation & M1 & 29 & 26 & 25 & 28 & 25 & 28 & 25 & 28 \\
%DIFDELCMD < extrapolation & M2 & 15/15 & 13 & 15/11 & 8 & 20/9.3 & 4 & 18/9.9 & 4 \\
%DIFDELCMD < extrapolation & M5 & -- & -- & 18/11 & 7 & 19/9.0 & 3 & 18/8.9 & 3 \\
%DIFDELCMD < extrapolation & NN & 11 & 8 & 10 & 6 & 9.5 & 4 & 9.5 & 4 \\
%DIFDELCMD < \botrule
%DIFDELCMD < \end{tabular}
%DIFDELCMD < \footnotetext{All subsets of $k$ produced alternatives of the family. Sibling: the subset holds an alternative at the new design's force; interpolation: no sibling, forces on both sides; extrapolation: no sibling, the new force outside. Wrong: share of wrong verdicts at $|d|=10$ (\%). M5 from $k=2$; interpolation needs at least two alternatives.}
%DIFDELCMD < \end{table}
%DIFDELCMD < %%%
%DIF < % END GENERATED design
\DIFdelend \DIFaddbegin \DIFadd{. The odd--even pattern of the bias correction in Fig.~\ref{fig:budget} comes from its median offset, which averages the offsets of two forces only for an even number of alternatives. With calibration series of 50, 100 or 250 parts instead of 500, no distance within a family moves by more than 1.2 floors, and for a new family by at most 4 floors (ESM Table~}\esm{robust}\DIFadd{): the qualification needs short runs of several variants, and only the floor needs a run long enough to show the drift.
}\DIFaddend 

\begin{figure}[t]
\centering
\includegraphics[width=\textwidth]{F4_budget}
\caption{Decisive and safe distances against the number of produced alternatives of the family, by the relation of the new design to them; bands: 95\,\% intervals over the held-out alternatives}\label{fig:budget}
\end{figure}

\DIFdelbegin \subsection{\DIFdel{Applicability guard}}%DIFAUXCMD
\addtocounter{subsection}{-1}%DIFAUXCMD
%DIFDELCMD < \label{sec:res-guard}
%DIFDELCMD < %%%
\DIFdelend \DIFaddbegin \subsection{\DIFadd{The procedure's decision rule}}\label{sec:res-procedure}
\DIFaddend 

\DIFdelbegin \DIFdel{A guard is useful only where it can fire, which within this design means a new process setting and a new family (Table~\ref{tab:guard}). For a new setting the range guard refuses exactly the extrapolated forces, the 84 of 126 cases whose distances are largest, and among the interpolated cases it lets through, M2, }\DIFdelend \DIFaddbegin \DIFadd{What a designer observes for a single verdict is the margin between the predicted interval and the requirement, and its value as evidence depends on the relation (Fig.~\ref{fig:reliability}): for a new variant at least 95\,\% of the decided verdicts of the interval rules are correct at every margin; for a new setting the share grows with the margin, for the nearest setting from 85\,\% at any margin to 95\,\% beyond 5 floors; for a new family }\DIFaddend M5 \DIFdelbegin \DIFdel{and }\DIFdelend \DIFaddbegin \DIFadd{is right in 71--74\,\% of its decided verdicts whatever the margin. Step 5 turns this into guard bands (ESM Table~}\esm{step6}\DIFadd{): none for the interval rules within the family and 0.25 floors for the nearest setting; 2.0 floors for an interpolated setting, with which }\DIFaddend the nearest setting \DIFdelbegin \DIFdel{are never wrong }\DIFdelend \DIFaddbegin \DIFadd{decides at 6.7 floors, is safe at 2.7 and sends no design to trial }\DIFaddend at 10 floors\DIFdelbegin \DIFdel{. The novelty guard , which compares the simulated $q_{95}$ of the new design with the calibration range , also refuses 84 cases but not the same ones; among the 80 verdicts it lets through, }\DIFdelend \DIFaddbegin \DIFadd{; 7.0 floors for an extrapolated one, with which it decides at 17 floors and sends a quarter of the designs to trial at 10 floors; and for a new family no guard band up to 20 floors reaches 95\,\%, so the procedure sends every design to trial. The guard bands are chosen on the evaluation data, which makes the trial rates optimistic.
}

\DIFadd{The applicability guards act on the same relations (ESM Table~}\esm{guard}\DIFadd{): for a new setting the range guard refuses exactly the extrapolated forces, after which }\DIFaddend M2\DIFdelbegin \DIFdel{and }\DIFdelend \DIFaddbegin \DIFadd{, }\DIFaddend M5 \DIFdelbegin \DIFdel{are wrong twice }\DIFdelend and the nearest setting \DIFdelbegin \DIFdel{never, but 72--96\,\% of the verdicts it refuses }\DIFdelend \DIFaddbegin \DIFadd{make no wrong verdict at 10 floors, whereas the novelty guard also refuses many cases that }\DIFaddend would have been \DIFdelbegin \DIFdel{correct.
For a new family the family and range guards refuse everything. The novelty guard lets 12 of 126 cases through, all of them }\DIFdelend \DIFaddbegin \DIFadd{decided correctly.
}

\emph{\DIFadd{A worked example.}} \DIFadd{A tryout engineer has to decide whether the concave cup at 300\,kN with medium lubrication, a new process setting, can be released without a trial run, with the six alternatives at 100 and 500\,kN produced and a requirement that 95\,\% of the parts have a }\DIFaddend mid-side draw-in \DIFdelbegin \DIFdel{or waviness, and among their 24 verdicts at 10 floors M5 is wrong in 6 (25}\DIFdelend \DIFaddbegin \DIFadd{of at most 15.30}\DIFaddend \,\DIFdelbegin \DIFdel{\%) and M1 in 6, whereas }\DIFdelend \DIFaddbegin \DIFadd{mm. The floor is 0.050\,mm, the relation is interpolation and no guard fires. The nominal simulation predicts 15.66\,mm and fails the design. The nearest produced setting predicts 14.99\,mm, a margin of 6.2 floors, more than its guard band of 2.0 floors, so step 6 returns meets; beyond a margin of 6 floors every decided verdict of this rule for an interpolated setting was correct. }\DIFaddend M2\DIFdelbegin \DIFdel{is never wrong and sends 12 }\DIFdelend \DIFaddbegin \DIFadd{, M5 and M5n return intervals reaching 15.37, 15.58 and 15.33\,mm and send the design }\DIFaddend to trial. \DIFdelbegin \DIFdel{A guard on the simulated value therefore protects a rule only as far as the rule's own interval does}\DIFdelend \DIFaddbegin \DIFadd{The 500 parts had a 95th percentile of 14.88\,mm, 8.4 floors below the limit: the procedure was right, the simulation wrong and the interval rules cautious}\DIFaddend .

%DIF < % BEGIN GENERATED guard
\DIFdelbegin %DIFDELCMD < \begin{table}[t]
%DIFDELCMD < \tabsize
%DIFDELCMD < \setlength{\tabcolsep}{3.5pt}
%DIFDELCMD < \caption{Applicability guards: distances and wrong verdicts among the cases a guard lets through}\label{tab:guard}
%DIFDELCMD < \begin{tabular}{@{}llrrrrrrr@{}}
%DIFDELCMD < \toprule
%DIFDELCMD <  & & & \multicolumn{2}{c}{M2} & \multicolumn{2}{c}{M5} & \multicolumn{2}{c}{NN} \\
%DIFDELCMD < \cmidrule(lr){4-5}\cmidrule(lr){6-7}\cmidrule(lr){8-9}
%DIFDELCMD < Scope & Guard & Guarded & $\delta_\mathrm{d}/\delta_\mathrm{s}$ & wrong & $\delta_\mathrm{d}/\delta_\mathrm{s}$ & wrong & $\delta_\mathrm{d}$ & wrong \\
%DIFDELCMD < \midrule
%DIFDELCMD < New setting & none & 0 & 18/8.1 & 6/239 & 18/7.4 & 4/239 & 8.5 & 6/239 \\
%DIFDELCMD <  & range & 84 & $\infty$/0.0 & 0/80 & $\infty$/0.9 & 0/80 & $\infty$ & 0/80 \\
%DIFDELCMD <  & novelty & 84 & $\infty$/0.0 & 2/80 & $\infty$/0.1 & 2/80 & $\infty$ & 0/80 \\
%DIFDELCMD < New family & none & 0 & 51/16 & 23/239 & 38/30 & 47/239 & 131 & 106/239 \\
%DIFDELCMD <  & range & 126 & $\infty$/0.0 & 0/0 & $\infty$/0.0 & 0/0 & $\infty$ & 0/0 \\
%DIFDELCMD <  & novelty & 114 & $\infty$/0.0 & 0/24 & $\infty$/0.0 & 6/24 & $\infty$ & 6/24 \\
%DIFDELCMD < \botrule
%DIFDELCMD < \end{tabular}
%DIFDELCMD < \footnotetext{Guarded: cases (of 126) sent to trial. Range guard: geometry not calibrated or blank-holder force outside the calibrated range; novelty guard: simulated $q_{95}$ outside the range of the calibration alternatives. Wrong: wrong verdicts at $|d|=10$ among the verdicts the guard lets through. The family guard equals the range guard for a new family and never fires for a new setting.}
%DIFDELCMD < \end{table}
%DIFDELCMD < %%%
%DIF < % END GENERATED guard
\DIFdelend \DIFaddbegin \begin{figure}[t]
\centering
\includegraphics[width=\textwidth]{F5_reliability}
\caption{Share of correct verdicts among the decided ones against the observable margin between interval and requirement, for requirements within 20 floors of the truth}\label{fig:reliability}
\end{figure}
\DIFaddend 

\subsection{Requirements \DIFdelbegin \DIFdel{from general }\DIFdelend \DIFaddbegin \DIFadd{anchored in capability and in }\DIFaddend tolerances}\DIFdelbegin %DIFDELCMD < \label{sec:res-scenarios}
%DIFDELCMD < %%%
\DIFdelend \DIFaddbegin \label{sec:res-requirements}
\DIFaddend 

\DIFdelbegin \DIFdel{Requirements from the angular classes of ISO~2768-1 \mbox{%DIFAUXCMD
\citep{iso1989tolerances}}\hskip0pt%DIFAUXCMD
, applied to the wall angle against its design angle and to the arm angle against flat with the row for the shorter leg of each, turn the verdicts into concrete design decisions (Table~\ref{tab:scenarios}). Of the 108 decisions, 69 concern an alternative that fails its limit, and the limit lies within 5 floors of the truth in 32 and within 10 floors in 55; deciding the worst side of a part instead of the mean of its sides would change the truth in 12 of them, since a general tolerance applies to every feature. For a new variant }\DIFdelend \DIFaddbegin \DIFadd{Upper limits at which an alternative has a performance index of 1.0, 1.33, 1.67 or 2.0 lie a median 1.0, 1.7, 2.5 and 3.1 floors above its 95th percentile, inside every decisive distance (ESM Table~}\esm{capability}\DIFadd{). Every alternative meets them, so a verdict is right, a false reject or a trial. At a performance index of 1.33 }\DIFaddend the nearest produced setting \DIFdelbegin \DIFdel{decides all 108 correctly and M5 93}\DIFdelend \DIFaddbegin \DIFadd{still says meets in 91}\DIFaddend \,\% \DIFdelbegin \DIFdel{, with no false accept; within 5 floors of the limit, }\DIFdelend \DIFaddbegin \DIFadd{of the new-variant cases, because most of its errors are far smaller than its 95\,\% distance, whereas }\DIFaddend M5 \DIFdelbegin \DIFdel{is right in 75}\DIFdelend \DIFaddbegin \DIFadd{says so in 44}\DIFaddend \,\% and sends the rest to trial\DIFdelbegin \DIFdel{, M2 is right in 25\,\%, and the bias-corrected simulation accepts three of the six failing alternatives there. For a new process setting the nearest setting is right in 97\,\% and the Gaussian process without the simulation in 92}\DIFdelend \DIFaddbegin \DIFadd{; for a new setting the two are right in 65 and 32}\DIFaddend \,\%, \DIFdelbegin \DIFdel{whereas M5 }\DIFdelend and \DIFdelbegin \DIFdel{M2 send 22--25\,\% to trial and falsely accept 4\,\% of the failing alternatives. For }\DIFdelend \DIFaddbegin \DIFadd{for }\DIFaddend a new family \DIFdelbegin \DIFdel{M5 is right in 68\,\% but rejects 46}\DIFdelend \DIFaddbegin \DIFadd{no rule exceeds 54}\DIFaddend \,\%\DIFdelbegin \DIFdel{of the adequate alternatives and accepts 9\,\% of the failing ones; M2 never accepts a failing alternative but sends 39\,\% to trial. With a false reject costing half a false accept, the cheapest rule per decision is the nearest produced setting }\DIFdelend \DIFaddbegin \DIFadd{. Release-level requirements are therefore decided at the design stage only within a family, and then by the production record. General angular tolerances of ISO~2768-1 \mbox{%DIFAUXCMD
\citep{iso1989tolerances} }\hskip0pt%DIFAUXCMD
give a coarser scenario that mostly tests whether a rule recognises gross failures: only 32 of its 108 decisions lie within 5 floors of the limit (ESM Table~}\esm{scenarios}\DIFadd{). Its cheapest rule depends on the costs (ESM Table~}\esm{costmap}\DIFadd{): }\DIFaddend for a new variant \DIFdelbegin \DIFdel{at any trial cost, the interval rules without the simulation (M2n}\DIFdelend \DIFaddbegin \DIFadd{the nearest setting at any cost}\DIFaddend , \DIFdelbegin \DIFdel{M4 and M5n, tied) }\DIFdelend for a new setting \DIFaddbegin \DIFadd{the production-only interval rules }\DIFaddend while a trial costs at most a tenth of a false accept and the nearest setting \DIFdelbegin \DIFdel{from a fifth}\DIFdelend \DIFaddbegin \DIFadd{beyond}\DIFaddend , and for a new family \DIFdelbegin \DIFdel{M2 up to a tenth and }\DIFdelend \DIFaddbegin \DIFadd{a trial of every design while a trial costs at most a tenth of a false accept, the envelope rule at a fifth and }\DIFaddend M5 \DIFdelbegin \DIFdel{from a fifth. A drawing of a deep-drawn part would normally carry profile and position tolerances to datums \mbox{%DIFAUXCMD
\citep{iso2017tolerancing, iso2021general}}\hskip0pt%DIFAUXCMD
, which the data do not allow to evaluate; the general tolerances serve as requirements with realistic margins}\DIFdelend \DIFaddbegin \DIFadd{beyond, for a false reject at half the cost of a false accept}\DIFaddend .

%DIF < % BEGIN GENERATED scenarios
\DIFdelbegin %DIFDELCMD < \begin{table}[t]
%DIFDELCMD < \tabsize
%DIFDELCMD < \setlength{\tabcolsep}{2.4pt}
%DIFDELCMD < \caption{Verdicts against requirements from general angular tolerances (ISO 2768-1)}\label{tab:scenarios}
%DIFDELCMD < \begin{tabular}{@{}lrrrrrrrrrrrr@{}}
%DIFDELCMD < \toprule
%DIFDELCMD <  & \multicolumn{4}{c}{New variant (\%)} & \multicolumn{4}{c}{New setting (\%)} & \multicolumn{4}{c}{New family (\%)} \\
%DIFDELCMD < \cmidrule(lr){2-5}\cmidrule(lr){6-9}\cmidrule(lr){10-13}
%DIFDELCMD < Rule & right & ${<}5$ & FA & trial & right & ${<}5$ & FA & trial & right & ${<}5$ & FA & trial \\
%DIFDELCMD < \midrule
%DIFDELCMD < M0 & 61 & 69 & 39 & 0 & 61 & 69 & 39 & 0 & 61 & 69 & 39 & 0 \\
%DIFDELCMD < M1 & 83 & 62 & 9 & 0 & 75 & 53 & 13 & 0 & 61 & 62 & 30 & 0 \\
%DIFDELCMD < M2 & 67 & 25 & 0 & 33 & 72 & 38 & 4 & 25 & 44 & 47 & 0 & 39 \\
%DIFDELCMD < M5 & 93 & 75 & 0 & 7 & 75 & 44 & 4 & 22 & 68 & 56 & 9 & 10 \\
%DIFDELCMD < NN & 100 & 100 & 0 & 0 & 97 & 91 & 0 & 0 & 31 & 44 & 65 & 0 \\
%DIFDELCMD < M5n & 94 & 78 & 0 & 6 & 92 & 72 & 0 & 8 & 31 & 44 & 62 & 2 \\
%DIFDELCMD < \botrule
%DIFDELCMD < \end{tabular}
%DIFDELCMD < \footnotetext{Wall and arm angle, classes f/m, c and v, 18 alternatives: 108 decisions per rule and scope, 69 with a failing alternative and 32 with the limit within 5 floors of the truth. right: correct verdicts; ${<}5$: correct verdicts among the decisions within 5 floors; FA: false accepts among the failing alternatives; trial: sent to trial.}
%DIFDELCMD < \end{table}
%DIFDELCMD < %%%
%DIF < % END GENERATED scenarios
%DIFDELCMD < 

%DIFDELCMD < %%%
\DIFdelend \subsection{Robustness}\label{sec:res-robust}

\DIFdelbegin \DIFdel{Table~\ref{tab:robust} repeats the analysis }\DIFdelend \DIFaddbegin \DIFadd{The analysis was repeated }\DIFaddend under variants of the data, the constants and the \DIFdelbegin \DIFdel{Gaussian-process specification. The variants rescale the distances but leave the conclusions in place: in every variant that includes it, the nearest produced setting remains the most decisive rule for a new variant and a new setting and M2 the safest for a new family, and the rules without the simulation fail for a new family (91 floors or more); only rules within a few floors of each other change places, such as M1, M5 and Mc for a new family. Correcting every characteristic for the }\DIFdelend \DIFaddbegin \DIFadd{model specifications (ESM Table~}\esm{robust}\DIFadd{). Adjusting for }\DIFaddend punch temperature changes \DIFdelbegin \DIFdel{the pooled floors by $-14$ to $+4$\,\% and the distances of the calibrated rules by at most }\DIFdelend \DIFaddbegin \DIFadd{no calibrated distance by more than }\DIFaddend 2.6 floors. \DIFdelbegin \DIFdel{Smoothing the simulation removes its extraction noise but not its bias: M5 is unchanged within the family and the envelope rule becomes less safe }\DIFdelend \DIFaddbegin \DIFadd{Dropping the warm-up of the first 150 parts shrinks the floors and so lengthens the distances (the nearest setting 4.1 floors with a sibling, 12 }\DIFaddend for a new \DIFdelbegin \DIFdel{family (25 instead of 16 floors). Deciding on the 90th or 99th percentile instead of the 95th moves no distance by more than 1.7 floors. Floors from batches of 25 or 100 parts rescale the distances by factors of 0.80--0.96 and 0.97--1.16}\DIFdelend \DIFaddbegin \DIFadd{setting)}\DIFaddend , and floors \DIFdelbegin \DIFdel{at the 90\,\% or 99\,\% quantile by 1.16--1.46 and 0.54--0.83. Requiring 90\,\% or 99\,\% correct verdicts instead of 95\,\% shortens the decisive distances by up to 43\,\% (within the family) or lengthens them by up to 93\,\% (for a new setting), as an operating characteristic should. The }\DIFdelend \DIFaddbegin \DIFadd{taken between series, 1.5 to 9.1 times the reference, shorten them (0.75 and 2.7 floors). Other floor protocols rescale the distances by 0.60 to 1.42, the }\DIFaddend Gaussian-process \DIFdelbegin \DIFdel{variants---a Mat\'ern kernel, a linear trend, one-hot lubrication, a floor on the length scales and the first specification without the noise floor---keep M5 at 6.4--6.9 floors for a new variant and 17.7--18.9 for a new setting. Leaving out one lubrication pattern behaves like a new variant (nearest setting 3.4, M5 6.4/1.1 floors), because the siblings at the same force remain. Refitting every rule on resampled calibration sets }\DIFdelend \DIFaddbegin \DIFadd{specifications move them by at most 2.1 floors, and in the unit of the reproducibility of $q_{95}$ the nearest setting decides at 2.8 and 6.9 floors. In every variant that includes it, the nearest produced setting remains the most decisive rule within a family and the envelope rule the safest for a new family. The refitting bootstrap }\DIFaddend widens the intervals \DIFdelbegin \DIFdel{, most for the Gaussian processes }\DIFdelend \DIFaddbegin \DIFadd{of the learned rules most }\DIFaddend (M5 \DIFaddbegin \DIFadd{4.0--11 floors }\DIFaddend for a new variant\DIFdelbegin \DIFdel{: 5.0--18.6 floors, against 2.2--7.0 for }\DIFdelend \DIFaddbegin \DIFadd{, }\DIFaddend the nearest setting \DIFdelbegin \DIFdel{), whose fits on calibration sets with repeated alternatives are unstable. Replacing each part's oil film by its pattern's median makes M3 less decisive (33 instead of 17 floors), and a conformalised Gaussian process (M6), the }\DIFdelend \DIFaddbegin \DIFadd{2.3--4.6), yet the nearest setting is the more decisive in all 200 resamples, against }\DIFaddend M5 \DIFdelbegin \DIFdel{mean with a margin on standardised leave-one-out residuals, is safe at 0.4 floors within the family but decisive only at 39: it buys safety for the over-confident Gaussian-process intervals at the cost of decisiveness}\DIFdelend \DIFaddbegin \DIFadd{and M5n and for a new setting alike (ESM Table~}\esm{paired}\DIFadd{)}\DIFaddend .

%DIF < % BEGIN GENERATED robust
\DIFdelbegin %DIFDELCMD < \begin{table}[t]
%DIFDELCMD < \tabsize
%DIFDELCMD < \setlength{\tabcolsep}{2.6pt}
%DIFDELCMD < \caption{Decisive/safe distances (floors) under the robustness variants}\label{tab:robust}
%DIFDELCMD < \begin{tabular}{@{}lrrrrrrrrr@{}}
%DIFDELCMD < \toprule
%DIFDELCMD <  & \multicolumn{4}{c}{New variant} & \multicolumn{3}{c}{New setting} & \multicolumn{2}{c}{New family} \\
%DIFDELCMD < \cmidrule(lr){2-5}\cmidrule(lr){6-8}\cmidrule(lr){9-10}
%DIFDELCMD < Variant & M1 & M2 & M5 & NN & M2 & M5 & NN & M2 & M5 \\
%DIFDELCMD < \midrule
%DIFDELCMD < Reference & 15 & 23/0.2 & 6.7/1.4 & 3.4 & 18/8.1 & 18/7.4 & 8.5 & 51/16 & 38/30 \\
%DIFDELCMD < Temperature-adjusted & 15 & 23/0.1 & 6.5/1.4 & 3.4 & 17/7.6 & 19/7.1 & 9.4 & 48/15 & 36/30 \\
%DIFDELCMD < Smoothed simulation & 15 & 25/0.3 & 6.6/1.4 & 3.4 & 18/7.6 & 19/7.2 & 8.5 & 49/25 & 34/30 \\
%DIFDELCMD < Adequacy share 0.90 & 15 & 23/0.2 & 6.7/1.6 & 3.3 & 18/7.5 & 18/7.0 & 8.4 & 51/15 & 38/31 \\
%DIFDELCMD < Adequacy share 0.99 & 15 & 23/0.0 & 7.2/1.7 & 3.3 & 18/8.3 & 20/7.1 & 8.7 & 51/16 & 38/30 \\
%DIFDELCMD < Floor: batches of 25 & 13 & 21/0.2 & 6.0/1.2 & 2.9 & 16/6.9 & 16/6.9 & 7.7 & 43/14 & 30/25 \\
%DIFDELCMD < Floor: batches of 100 & 16 & 24/0.2 & 7.6/1.5 & 3.5 & 19/8.5 & 20/7.8 & 8.9 & 57/17 & 44/31 \\
%DIFDELCMD < Floor: 90\,\% quantile & 19 & 31/0.2 & 8.3/1.7 & 4.0 & 21/10 & 22/8.6 & 10 & 61/20 & 44/37 \\
%DIFDELCMD < Floor: 99\,\% quantile & 10 & 14/0.1 & 4.7/1.0 & 2.4 & 11/4.4 & 13/4.5 & 5.7 & 37/11 & 28/17 \\
%DIFDELCMD < Target 90\,\% & 12 & 17/0.0 & 5.2/0.5 & 1.9 & 16/5.2 & 17/4.8 & 5.4 & 45/9.6 & 22/19 \\
%DIFDELCMD < Target 99\,\% & 20 & 28/1.4 & 11/5.2 & 5.1 & 27/11 & 34/11 & 11 & 60/28 & 59/35 \\
%DIFDELCMD < Interval level 0.80 &  & 23/0.2 & 6.0/1.6 &  & 18/8.1 & 16/8.1 &  & 50/17 & 37/30 \\
%DIFDELCMD < Interval level 0.95 &  & 23/0.2 & 7.4/1.0 &  & 18/8.1 & 20/6.4 &  & 51/16 & 38/29 \\
%DIFDELCMD < GP: Mat\'ern 5/2 &  &  & 6.4/1.1 &  &  & 19/6.8 &  &  & 38/30 \\
%DIFDELCMD < GP: linear trend &  &  & 6.7/1.4 &  &  & 19/6.3 &  &  & 38/30 \\
%DIFDELCMD < GP: one-hot lubrication &  &  & 6.8/2.2 &  &  & 18/6.9 &  &  & 34/32 \\
%DIFDELCMD < GP: length-scale floor &  &  & 6.9/1.2 &  &  & 18/6.5 &  &  & 38/29 \\
%DIFDELCMD < GP: no noise floor &  &  & 6.7/1.5 &  &  & 18/7.4 &  &  & 38/30 \\
%DIFDELCMD < Leave one pattern out & 15 & 22/0.3 & 6.4/1.1 & 3.4 & & & & & \\
%DIFDELCMD < \midrule Refit bootstrap, $\delta_\mathrm{d}$ & 14--31 & 19--28 & 5.0--19 & 2.2--7.0 & 16--27 & 17--23 & 5.4--14 &  & \\
%DIFDELCMD < \botrule
%DIFDELCMD < \end{tabular}
%DIFDELCMD < \footnotetext{Alternative-level rules refitted for every variant; empty cells: the variant does not affect the rule. Smoothed simulation: nominal value and envelope from a quadratic response surface over thickness and friction. Floor variants rescale the distances with other floors. Leave one pattern out: columns M1, M2, M5 and NN of the first block. Refit bootstrap: 95\,\% interval of the decisive distance when the alternatives of each geometry are resampled and every rule is refitted (200 resamples).}
%DIFDELCMD < \end{table}
%DIFDELCMD < %%%
%DIF < % END GENERATED robust
%DIFDELCMD < 

%DIFDELCMD < %%%
\DIFdelend %% ----------------------------------------------------------------------
\section{Discussion}\label{sec:discussion}

\subsection{What the evaluation establishes}

Five findings answer the \DIFdelbegin \DIFdel{three research questionsbeyond the particular numbers}\DIFdelend \DIFaddbegin \DIFadd{research questions. Two follow from the definitions and hold wherever the framework is applied}\DIFaddend . First (RQ1), \DIFdelbegin \DIFdel{production has a resolution, and the floor makes decision fitness comparable across characteristics and predictors: it is set by drift within a run rather than by the scatter of single parts or by the measurement, the decided percentile reproduces to about one floor, and the distances of very different rules can be read on one scale . It }\DIFdelend \DIFaddbegin \DIFadd{the floor puts every rule and characteristic on a common production-referenced scale, but not on a common scale of nonconforming fractions, and as a within-run quantity it }\DIFaddend is a lower bound \DIFdelbegin \DIFdel{for series production, since separate runs of one designmay differ by several floors more. Second }\DIFdelend \DIFaddbegin \DIFadd{of the variation of series production. Second, an interval rule trades decisiveness for safety, the two distances make the trade explicit, and the nominal level of an interval is reached only where the calibration alternatives are exchangeable with the new design. Three are observations on this dataset and hypotheses for confirmation. Third }\DIFaddend (RQ2), the \DIFdelbegin \DIFdel{relation of a new design to the produced evidence governs decision fitness more than the model does. With a produced sibling , the best rules calibrated on production decide within a few floors; }\DIFdelend \DIFaddbegin \DIFadd{best attainable distance grows with the novelty of the design, from 3.4 floors with a sibling (0.85 when the sibling is a near-replicate) to 8.5 }\DIFaddend for a new \DIFdelbegin \DIFdel{process setting no rule decides closer than 8.5 floors , and extrapolating in force doubles the decisive distance of the nearest produced setting and more than triples the safe distances of M2 and M5; for a new family no rule is safe closer than 16 floors. Third}\DIFdelend \DIFaddbegin \DIFadd{setting (4.7 interpolated, 9.5 extrapolated with a speed change) and to no rule safe closer than 16 floors in two transfers, while within a family the choice of rule matters more than the relation. Fourth}\DIFaddend , production evidence buys safety before decisiveness: interval rules become safe with few produced alternatives that bracket the new design\DIFdelbegin \DIFdel{but remain indecisive, and }\DIFdelend \DIFaddbegin \DIFadd{, and decisive }\DIFaddend only \DIFaddbegin \DIFadd{with }\DIFaddend near-replicates\DIFdelbegin \DIFdel{make them decisive (Section~\ref{sec:res-budget}). Fourth, interval rules trade abstention for safety, the two distances make the trade explicit, and the nominal level of an interval is reached only where the calibration alternatives are exchangeable with the new design}\DIFdelend . Fifth (RQ3), in this \DIFdelbegin \DIFdel{data regime the simulation's contribution is confined to a new family }\DIFdelend \DIFaddbegin \DIFadd{small-data regime and with this simulator, the simulation contributed to decisions only across families as an offset, a fitted scale made it useful within a family but harmful for a new setting}\DIFaddend , and learned models \DIFdelbegin \DIFdel{add calibrated abstention rather than accuracy over }\DIFdelend \DIFaddbegin \DIFadd{were as accurate as }\DIFaddend the nearest produced setting \DIFaddbegin \DIFadd{as point predictors while adding abstention}\DIFaddend .

\DIFdelbegin %DIFDELCMD < \begin{table}[t]
%DIFDELCMD < \tabsize
%DIFDELCMD < \setlength{\tabcolsep}{4pt}
%DIFDELCMD < \caption{Guidance for design-stage manufacturability decisions}\label{tab:guidelines}
%DIFDELCMD < \begin{tabular}{@{}L{4.1cm}L{8.1cm}@{}}
%DIFDELCMD < \toprule
%DIFDELCMD < \multicolumn{2}{@{}l}{\emph{From the framework (general)}} \\
%DIFDELCMD < Before trusting a predictor & Measure the floor on consecutive production of the family and report it with its protocol; treat it as a lower bound of series variation. \\
%DIFDELCMD < Qualifying a rule & Evaluate it for the evidence relation of the coming designs (grouped cross-validation), not by leave-one-out over variants that production cannot tell apart. \\
%DIFDELCMD < Choosing a rule & Compare every simulation- or learning-based rule with the nearest produced setting and with its own version without the simulation; keep it only where it is closer or safer. \\
%DIFDELCMD < Fixing a requirement & If the expected margin to the capability is below the safe distance of the best qualified rule, plan a trial or further production evidence. \\
%DIFDELCMD < Planning production evidence & Produce variants that bracket the settings of the coming designs; near-replicates make decisions easy but do not inform new settings. \\
%DIFDELCMD < A new design family & Use the simulation for gross margins only; send decisions near the requirement to trial. \\
%DIFDELCMD < \multicolumn{2}{@{}l}{\emph{Observed on this dataset (18 alternatives, 2 geometries, 7 characteristics)}} \\
%DIFDELCMD < Nominal simulation & reliable beyond about 108 floors; errs mostly by false rejects \\
%DIFDELCMD < New variant with a produced sibling & nearest produced setting right beyond 3.4 floors; calibrated interval rules safe at or below 1.4 floors \\
%DIFDELCMD < New process setting & nearest produced setting right beyond 8.5 floors (interpolation 4.7, extrapolation 9.5); margin and Gaussian-process rules safe at 5.4--8.1 floors \\
%DIFDELCMD < New family (two calibrations) & best rule safe at 16 floors; rules without the simulation fail \\
%DIFDELCMD < Changes below the minimum resolvable change & invisible in the batch centres; they may still matter in the tail of a characteristic \\
%DIFDELCMD < \botrule
%DIFDELCMD < \end{tabular}
%DIFDELCMD < \end{table}
%DIFDELCMD < %%%
\DIFdelend \DIFaddbegin \begin{table}[t]
\tabsize
\setlength{\tabcolsep}{4pt}
\caption{Guidance for design-stage manufacturability decisions, with its basis}\label{tab:guidelines}
\begin{tabular}{@{}L{3.4cm}L{6.6cm}L{2.2cm}@{}}
\toprule
Situation & Guidance & Basis \\
\midrule
\multicolumn{3}{@{}l}{\emph{From the framework (general)}} \\
Before trusting a predictor & Measure the floor on consecutive production of the family with a stated protocol and report $\sigma_\mathrm{b}$, $\sigma_\mathrm{w}$ and its lag dependence; treat it as a lower bound of series variation. & definitions \\
Qualifying a rule & Evaluate it for the evidence relation of the coming designs by grouped cross-validation, and split near-replicate siblings from resolvable ones. & definitions \\
Choosing a rule & Compare it with the nearest produced setting and with its own version without the simulation, on decisive and safe distance jointly or by expected cost. & definitions \\
Deciding a single design & Apply the rule's guard band to the observable margin; a margin inside the guard band, or a fired guard, means a trial. & definitions \\
Fixing a requirement & If the expected margin to the capability is below the safe distance of the best qualified rule, plan a trial; no decision is reliable closer than a trial reproduces the quantile. & definitions \\
\multicolumn{3}{@{}l}{\emph{Observed on this dataset (18 alternatives, 2 geometries, 7 characteristics)}} \\
New variant with a produced sibling & nearest produced setting right beyond 3.4 floors (0.85 for a near-replicate); calibrated interval rules safe below one floor & 126 cases \\
New process setting & nearest produced setting right beyond 4.7 floors interpolated and 9.5 extrapolated; interval rules safe at 5.9--16 floors & 42 + 84 cases \\
New family & best rule safe at 16 floors; rules without the simulation fail; a transferable drawing nominal helps & two transfers \\
Release-level requirements & lie within 1--3 floors of the 95th percentile, inside every decisive distance & 126 cases \\
\botrule
\end{tabular}
\end{table}
\DIFaddend 

\subsection{Using the framework in design}

Table~\ref{tab:guidelines} separates guidance that follows from the framework \DIFdelbegin \DIFdel{and plausibly holds elsewhere }\DIFdelend from observations on this dataset, with their \DIFdelbegin \DIFdel{evidence base. The general guidance turns the distances into planning questions}\DIFdelend \DIFaddbegin \DIFadd{basis}\DIFaddend . Before a requirement is fixed, the expected margin between \DIFdelbegin \DIFdel{the requirement and the capability of the design }\DIFdelend \DIFaddbegin \DIFadd{requirement and capability }\DIFaddend can be compared with the safe distance of the best rule qualified for the \DIFdelbegin \DIFdel{evidence relation of that }\DIFdelend \DIFaddbegin \DIFadd{relation of the }\DIFaddend design; if the margin is smaller, no design-stage decision can be \DIFdelbegin \DIFdel{trusted }\DIFdelend \DIFaddbegin \DIFadd{relied on, }\DIFaddend and a trial or further production evidence has to be planned. \DIFaddbegin \DIFadd{Release-level requirements make this the usual case: they lie within a few floors of the capability, where only a near-replicate production record decides reliably, which is consistent with the practice of requiring a production trial run for a process change. }\DIFaddend For a single \DIFdelbegin \DIFdel{verdict, }\DIFdelend \DIFaddbegin \DIFadd{design, step 6 turns the qualification into a verdict whose guard band a designer can see, and }\DIFaddend the \DIFdelbegin \DIFdel{margin between the predicted interval and the requirement is what a designer observes, and it is informative only where the rule is calibrated for the relation (Section~\ref{sec:res-ml})}\DIFdelend \DIFaddbegin \DIFadd{worked example shows that the most decisive rule is also the most transparent: the nearest produced setting points to the produced alternative that it copies, whereas a Gaussian-process interval does not}\DIFaddend . Which rule to prefer then depends on the costs\DIFdelbegin \DIFdel{: with the realistic margins of general tolerances, a rule that always decides was cheapest where it was accurate, and an interval rule where a false accept is expensive relative to a trial .
}\DIFdelend \DIFaddbegin \DIFadd{, and those differ by relation: a trial of a new force on an existing tool is a short press run, whereas a trial of a new family may need a new tool. In a company, steps 1 to 5 fit the start of production or the production part approval of a family, when its series are measured anyway, and step 6 the later decisions on variants and process changes.
}\DIFaddend 

\DIFdelbegin \subsection{\DIFdel{Where simulation and machine learning help}}
%DIFAUXCMD
\addtocounter{subsection}{-1}%DIFAUXCMD
\DIFdelend \DIFaddbegin \subsection{\DIFadd{Simulation and machine learning in this regime}}
\DIFaddend 

\DIFdelbegin \DIFdel{The question the call for this collection poses about learning from manufacturing data receives a sober answer here.
}\DIFdelend Within a design family \DIFdelbegin \DIFdel{, }\DIFdelend the production record \DIFdelbegin \DIFdel{carries the decision: the nearest produced setting is the most decisive rule}\DIFdelend \DIFaddbegin \DIFadd{carried the decision}\DIFaddend , and adding the simulation \DIFdelbegin \DIFdel{to }\DIFdelend \DIFaddbegin \DIFadd{as an offset did not improve }\DIFaddend a calibrated rule\DIFdelbegin \DIFdel{does not improve it and, for the bias correction, makes it worse}\DIFdelend , because the \DIFdelbegin \DIFdel{simulation's process trends---the effect of blank-holder force above all---are not credible. A physically motivated calibration of }\DIFdelend \DIFaddbegin \DIFadd{simulated force trend was 2.7 to 5.1 times too large. Rescaling the simulation repaired this within the family but extrapolated badly for a new setting, and calibrating }\DIFaddend the friction coefficient did not repair \DIFdelbegin \DIFdel{this}\DIFdelend \DIFaddbegin \DIFadd{it}\DIFaddend , which points to the friction law and the blank-holder model rather than to a parameter. \DIFdelbegin \DIFdel{Learning models , a Gaussian process or gradient boosting with or without the simulation, predict about }\DIFdelend \DIFaddbegin \DIFadd{Learned models predicted }\DIFaddend as well as the nearest produced setting; what they \DIFdelbegin \DIFdel{add is an interval , which makes a rule }\DIFdelend \DIFaddbegin \DIFadd{added was an interval that is }\DIFaddend safe where the calibration alternatives are exchangeable with the new design and over-confident where they are not. \DIFdelbegin \DIFdel{The simulation is indispensable only for a new family , where nothing in the production record of another family carries over, but even there it supports decisions only at margins that a design team would not need a model for. A model's interval is a statement about the model; the safe distancefor the evidence relation is a statement about production}\DIFdelend \DIFaddbegin \DIFadd{A multi-task model cannot help a family with no produced alternative, because the correlation between families cannot be estimated without its data; the pooled Gaussian process, which carries the geometry as an input, is the borrowing of strength that the data allow, and it improved coverage but not decisiveness. These findings concern one simulator of low fidelity for this production, used through offsets and one scale, and learners with six to nine produced settings each; a surrogate of these simulations inherits the simulator's distance, whatever its accuracy against the simulation}\DIFaddend . For the evaluation of AI-based DFM support \DIFdelbegin \DIFdel{this means that accuracy on }\DIFdelend \DIFaddbegin \DIFadd{the lesson is methodological: }\DIFaddend held-out \DIFdelbegin \DIFdel{data }\DIFdelend \DIFaddbegin \DIFadd{accuracy }\DIFaddend is not enough\DIFdelbegin \DIFdel{: }\DIFdelend \DIFaddbegin \DIFadd{, and }\DIFaddend a learned model should be compared with a production-record baseline\DIFdelbegin \DIFdel{and }\DIFdelend \DIFaddbegin \DIFadd{, }\DIFaddend scored by evidence relation \DIFdelbegin \DIFdel{, }\DIFdelend in a unit that production defines\DIFaddbegin \DIFadd{, and shown to its user with a guard band. Whether that makes designers' reliance match a rule's reliability \mbox{%DIFAUXCMD
\citep{lee2004trust}}\hskip0pt%DIFAUXCMD
, and when an abstention should go to an expert rather than to a trial \mbox{%DIFAUXCMD
\citep{madras2018predict}}\hskip0pt%DIFAUXCMD
, are hypotheses for a study with design teams}\DIFaddend .

\subsection{Implications for design research}

The framework gives \DIFdelbegin \DIFdel{three }\DIFdelend \DIFaddbegin \DIFadd{several }\DIFaddend constructs of design research a production-referenced measure. \DIFdelbegin \DIFdel{The safe distance of a rule for an evidence relation is the margin a design team has to keep between a prediction and a requirement for the verdict to be trusted, and so puts a number on what design margins are for \mbox{%DIFAUXCMD
\citep{eckert2019design, brahma2020margin}}\hskip0pt%DIFAUXCMD
. The two distances decide }\DIFdelend \DIFaddbegin \DIFadd{For variant, adaptive and platform design \mbox{%DIFAUXCMD
\citep{pahl2007engineering, simpson2006product}}\hskip0pt%DIFAUXCMD
, the distances by evidence relation quantify how far production knowledge carries from produced variants to a new one, and the nearest produced setting sets the baseline that any capability database \mbox{%DIFAUXCMD
\citep{thornton2004variation} }\hskip0pt%DIFAUXCMD
or learned model has to beat. For margins, the safe distance is the part of a margin that covers predictive uncertainty, as distinct from the part that absorbs change \mbox{%DIFAUXCMD
\citep{eckert2017safety, eckert2019design}}\hskip0pt%DIFAUXCMD
, and the guard band is its counterpart for a single decision. For testing and prototyping strategies \mbox{%DIFAUXCMD
\citep{thomke2001sequential, salado2018mathematical, camburn2017design}}\hskip0pt%DIFAUXCMD
, the distances and the reproducibility of a trial say }\DIFaddend when a test is cheaper than a decision \DIFdelbegin \DIFdel{, the trade at the core of testing strategies \mbox{%DIFAUXCMD
\citep{thomke2001sequential, salado2018mathematical}}\hskip0pt%DIFAUXCMD
}\DIFdelend \DIFaddbegin \DIFadd{and how close to a requirement a test can itself decide. For overlapped development \mbox{%DIFAUXCMD
\citep{krishnan1997model, terwiesch2002exchanging} }\hskip0pt%DIFAUXCMD
and set-based design \mbox{%DIFAUXCMD
\citep{sobek1999toyota}}\hskip0pt%DIFAUXCMD
, they state when design-stage information is fit to act on and how wide a feasibility limit has to be. And for verification and validation, the evidence relation is the question of the application domain \mbox{%DIFAUXCMD
\citep{oberkampf2010verification}}\hskip0pt%DIFAUXCMD
}\DIFaddend , and the \DIFdelbegin \DIFdel{calibration budget turns the production of variants into a planned verification activity. And they show how far simulation front-loads a decision \mbox{%DIFAUXCMD
\citep{thomke2000effect}}\hskip0pt%DIFAUXCMD
: only as far as its distance for the relation at hand}\DIFdelend \DIFaddbegin \DIFadd{distances could serve as the evidence of a risk-informed credibility assessment \mbox{%DIFAUXCMD
\citep{asme2018assessing}}\hskip0pt%DIFAUXCMD
}\DIFaddend . In the terms of the design research methodology \citep{blessing2009drm}, the evaluation reported here \DIFdelbegin \DIFdel{establishes that the framework discriminates between rules and evidence situations on real production data---its empirical performance }\DIFdelend \DIFaddbegin \DIFadd{is a support evaluation: it shows structural }\DIFaddend validity in the sense of the validation square\DIFdelbegin \DIFdel{\mbox{%DIFAUXCMD
\citep{seepersad2006validation}}\hskip0pt%DIFAUXCMD
---while its use by design teams, and the success criteria that would link the distances to fewer trials or late changes, remain to be evaluated}\DIFdelend \DIFaddbegin \DIFadd{, that the metrics discriminate between rules and relations on real data \mbox{%DIFAUXCMD
\citep{seepersad2006validation}}\hskip0pt%DIFAUXCMD
, and is a first step towards performance validity, which would require showing that decisions taken with the framework are better than decisions taken without it}\DIFaddend .

\subsection{\DIFdelbegin \DIFdel{Applying the framework elsewhere}\DIFdelend \DIFaddbegin \DIFadd{Reuse for other processes and AI-based DFM tools}\DIFaddend }\DIFaddbegin \label{sec:reuse}
\DIFaddend 

The framework needs \DIFdelbegin \DIFdel{four things: }\DIFdelend several produced variants of a design family \DIFdelbegin \DIFdel{, each with }\DIFdelend \DIFaddbegin \DIFadd{whose relations to each other span those of the coming designs, }\DIFaddend a series of consecutive parts or replicated batches of one design \DIFdelbegin \DIFdel{from which the floor can be measured ; variants whose relations to each other span the relations of the coming designs; a }\DIFdelend \DIFaddbegin \DIFadd{per variant, a }\DIFaddend design-stage \DIFdelbegin \DIFdel{predictor for every variant, which may be a simulation, a surrogate or the production record itself; and characteristics that can be measured identically on parts and predictions. None is specific to forming}\DIFdelend \DIFaddbegin \DIFadd{predictor for every variant, and characteristics measured identically on parts and predictions. On this dataset the long series served the floor, whereas calibration runs of about 50 parts per variant qualified the rules as well as full series did; how many long runs a floor needs was not tested. Table~\ref{tab:processes} maps these onto the processes named most often in DFM; attribute outcomes such as splits, porosity or short shots call for a requirement on a rate rather than on a quantile, which the framework does not yet cover. A predictor need not return an interval: any tool that returns meets, fails or trial for a stated requirement, a classifier at its operating threshold or a language-model check given the requirement value, can be scored by sweeping the requirement. For continuous design spaces, the relation becomes a distance from the produced designs in descriptor space, or the position inside or outside their hull. A minimal report of such an evaluation comprises the floor with its protocol and variance components, the decisive and safe distances by relation with refitting intervals and one-sided values, the coverage of interval rules, the nearest-produced-setting and no-simulation baselines, and the guard bands and guard performance. Distances per characteristic rest on as many cases as there are produced alternatives, so a team would pool characteristics with similar ratios of floor to part scatter, or accumulate evaluations over families, before relying on them}\DIFaddend . An attempt to apply \DIFdelbegin \DIFdel{it }\DIFdelend \DIFaddbegin \DIFadd{the framework }\DIFaddend to an open dataset of powder bed fusion of PA12 \citep{leirmo2021data} \DIFdelbegin \DIFdel{showed what the first condition demands: grouping the }\DIFdelend \DIFaddbegin \DIFadd{failed on the second requirement: grouping }\DIFaddend specimens by orientation class \DIFdelbegin \DIFdel{would have mixed several }\DIFdelend \DIFaddbegin \DIFadd{mixed }\DIFaddend orientations in a batch, and \DIFdelbegin \DIFdel{only }\DIFdelend eight identical anchor specimens per build \DIFdelbegin \DIFdel{replicate one design. Two-sided requirements follow from the absolute deviation from the target, as for the arm angle and the dome; requirements on a high conformance share are scored in the same way as long as the series estimates the quantile, and need a parametric tail beyond it}\DIFdelend \DIFaddbegin \DIFadd{were too few for a floor; a between-build floor from replicated anchors would be the remedy}\DIFaddend .

\subsection{Threats to validity}

The evaluation covers one material, one press and two geometries, and the relations within a family come from a $3\times3$ factorial in which one factor, lubrication, is mostly not resolved; the procedure is general, the numbers are not. The new family rests on two \DIFdelbegin \DIFdel{calibrations}\DIFdelend \DIFaddbegin \DIFadd{transfers between geometries that differ in three features at once, and the intervals of its distances hold the two fits fixed}\DIFaddend . Each alternative was produced once, so the floor covers drift within a run but not between days, coils or tool states\DIFdelbegin \DIFdel{; the quasi-replicate series bound that variation from above }\DIFdelend \DIFaddbegin \DIFadd{, the reference $q_{95}$ is that of one run, and the order and dates of the series are unknown; between-series variation bounds the missing component }\DIFaddend at 1.2--9.8 floors \DIFdelbegin \DIFdel{. No gauge study or acceptance test of the scanner \mbox{%DIFAUXCMD
\citep{vdi2012optical, iso2021measurement} }\hskip0pt%DIFAUXCMD
was available: the redundancy of the scans bounds the independent measurement noise, but a }\DIFdelend \DIFaddbegin \DIFadd{(Section~\ref{sec:res-robust}). A }\DIFaddend drift of the scanner cannot be separated from that of the process, \DIFaddbegin \DIFadd{so the floor is a production-and-measurement floor, and scan }\DIFaddend and \DIFdelbegin \DIFdel{systematic scan errors and differences of reference surface }\DIFdelend \DIFaddbegin \DIFadd{reference-surface differences }\DIFaddend enter the absolute offsets of the simulation\DIFdelbegin \DIFdel{, though not the floors and effects. The simulated characteristics carry extraction noise of up to several floors, and the documentation of the simulation leaves the material model open. The blank-holder force of 100\,kN ran at a slower stroke. Lubrication is a pattern whose oil amount varies between series}\DIFdelend \DIFaddbegin \DIFadd{. The extrapolation in force is confounded with stroke speed, which the simulation does not represent, so the extrapolation results describe a force and speed change together}\DIFaddend . The requirements are one-sided limits on the 95th percentile of a characteristic averaged over sides, \DIFdelbegin \DIFdel{the scenario covers two angles, and the reference $q_{95}$ has a standard error of about 0.1 floor}\DIFdelend \DIFaddbegin \DIFadd{which is below the conformance of series release}\DIFaddend . Several rules, the grouped schemes and the \DIFdelbegin \DIFdel{estimator were added after the first }\DIFdelend \DIFaddbegin \DIFadd{guard bands were added or chosen after earlier }\DIFaddend results had been seen and are evaluated on the same data, so the comparison between rules carries selection optimism \citep{cawley2010over}; the analysis was carried out by a single analyst, with code and data open for replication.

%% ----------------------------------------------------------------------
\section{Conclusions}\label{sec:conclusion}

Design-stage manufacturability decisions can be evaluated against production \DIFdelbegin \DIFdel{, in a unit }\DIFdelend \DIFaddbegin \DIFadd{on a scale }\DIFaddend that production defines. The production \DIFdelbegin \DIFdel{floor---the difference that drift and scatter produce between two batches of one design---expresses }\DIFdelend \DIFaddbegin \DIFadd{floor expresses }\DIFaddend the fitness of any predictor for a decision \DIFdelbegin \DIFdel{independently of the predictor and the characteristic, and }\DIFdelend \DIFaddbegin \DIFadd{in the units of the characteristic; }\DIFaddend the decisive and safe distances, read as operating characteristics, tell how close to a requirement a rule decides and how close it can be \DIFdelbegin \DIFdel{trusted }\DIFdelend \DIFaddbegin \DIFadd{relied on }\DIFaddend not to mislead\DIFaddbegin \DIFadd{; and a guard band turns them into a rule for a single decision}\DIFaddend .

On 18 deep-drawn design--process alternatives with matched simulations, the \DIFdelbegin \DIFdel{relation of a new designto the produced evidence governed decision fitness more than the model did. With }\DIFdelend \DIFaddbegin \DIFadd{best attainable distance grew with the novelty of the new design: 3.4 floors with }\DIFaddend a produced sibling\DIFdelbegin \DIFdel{the production record decided within a few floors, }\DIFdelend \DIFaddbegin \DIFadd{, 8.5 }\DIFaddend for a new process setting\DIFdelbegin \DIFdel{no rule decided closer than 8.5 floors, and for a new family no rule was }\DIFdelend \DIFaddbegin \DIFadd{, and no rule }\DIFaddend safe closer than 16 floors \DIFdelbegin \DIFdel{. The simulation contributed only to the new family; learned models matched the nearest produced setting in accuracy and added calibrated abstention where the calibration alternatives resembled the new design}\DIFdelend \DIFaddbegin \DIFadd{in two transfers to a new family. Within a family the choice of rule mattered more than the relation, release-level requirements lay inside every decisive distance, and in this small-data regime the simulation as used, a model with a large and partly sign-inconsistent bias, contributed only across families, while learned models added abstention rather than accuracy}\DIFaddend .

Four extensions follow: series produced on several days, coils and tool states, so that the floor covers variation between runs; design families with continuous geometric descriptors\DIFdelbegin \DIFdel{, so that new designs can be placed between produced ones; a cost model that turns the distances into decisions on trials and margins}\DIFdelend \DIFaddbegin \DIFadd{; a confirmatory evaluation of the frozen procedure on another dataset or process}\DIFaddend ; and an evaluation \DIFdelbegin \DIFdel{of the procedure }\DIFdelend with design teams \DIFaddbegin \DIFadd{of whether decisions taken with the framework are better}\DIFaddend . The released code, which reproduces every result from the public data, is the starting point for all four.

\FloatBarrier
\backmatter

\DIFaddbegin \bmhead{Supplementary information}
\DIFadd{The Electronic Supplementary Material contains Tables~S1--}\esm{last} \DIFadd{and Fig.~S1 with the further results referred to in the text.
}

\DIFaddend \bmhead{Acknowledgements}
The author thanks the authors of the RDDAC and DDACS datasets at the University of Stuttgart for publishing the data on which this work rests.

\section*{Declarations}

\bmhead{Funding}
No specific funding was received for this work.

\bmhead{Competing interests}
The author declares no competing interests.

\bmhead{Ethics approval and consent}
Not applicable.

\bmhead{Data availability}
All data are public. RDDAC: doi:10.18419/DARUS-5589 (version 2.0); DDACS: doi:10.18419/DARUS-4801 (version 3.0); both licensed CC BY 4.0. The extracted per-part and per-simulation feature tables are provided with the code.

\bmhead{Code availability}
The complete analysis pipeline is released under an open licence at \url{https://github.com/htcanseven/Claude} (directory \texttt{scripts/}\DIFdelbegin \DIFdel{; an archived release with a DOI will be deposited on Zenodo at acceptance}\DIFdelend ) and reproduces every number in this paper from the public data\DIFaddbegin \DIFadd{; a versioned archive of the code and the feature tables is deposited on Zenodo \pending{DOI}}\DIFaddend .

\bmhead{Author contributions}
H.T.C. conceived the method, designed and carried out the analysis, and wrote the paper.

\bmhead{Use of generative AI}
A large language model (Claude, Anthropic) was used under the author's direction as an assistant for code development, literature retrieval, data analysis\DIFdelbegin \DIFdel{and }\DIFdelend \DIFaddbegin \DIFadd{, }\DIFaddend drafting and editing of the text\DIFaddbegin \DIFadd{, and for simulated peer review of earlier versions}\DIFaddend . The author designed the study, verified all analyses, checked every reference against its published record, and takes full responsibility for the content. No AI-generated images are used; all figures are produced by the released code from the data.

%% ----------------------------------------------------------------------
\begin{appendices}

\section{\DIFdelbegin \DIFdel{Computation details }\DIFdelend \DIFaddbegin \DIFadd{Notation }\DIFaddend and \DIFdelbegin \DIFdel{additional tables}\DIFdelend \DIFaddbegin \DIFadd{reuse}\DIFaddend }\label{app:appendix}

\DIFdelbegin \subsection{\DIFdel{Computation details}}%DIFAUXCMD
\addtocounter{subsection}{-1}%DIFAUXCMD
%DIFDELCMD < \label{app:computation}
%DIFDELCMD < 

%DIFDELCMD < %%%
\emph{\DIFdel{Feature extraction.}} %DIFAUXCMD
\DIFdel{Scans are reduced to a grid of 2$\times$4-pixel medians after a luminescence filter (connected components of at least 20,040 pixels); calibration 0.0769\,mm per pixel across the scan line, 0.1581\,mm per line and 0.00775\,mm per height unit. The cup-bottom plane is fitted robustly to the plateau at least 10\,mm inside its edge, and all depths are taken from it. The flange comprises the points 20--40\,mm below the bottom plane, at least 8\,mm from the cup wall and 3\,mm from the outer edge. Wall positions are interpolated where the binned (0.5\,mm) radial profile reaches 10, 15 and 20\,mm depth; the wall angle is taken between 10 and 20\,mm. The arm angle is fitted to the longest run of the profile whose slope stays below 15$^\circ$, shortened by 1\,mm at both ends and at least 4\,mm long. Flange widths are medians over bands $\pm$2\,mm about the mid-lines; corner tips are the outline points within 0.3\,mm of the extreme diagonal projection. Simulated parts are mirrored from the quarter model and processed with the same definitions on a 1\,mm grid.
}%DIFDELCMD < 

%DIFDELCMD < %%%
\emph{\DIFdel{Floors and effects.}} %DIFAUXCMD
\DIFdel{Batches of $B=50$ consecutive parts; batches with fewer than 80\,\% valid values are dropped; centres are 10\,\% trimmed means; $\alpha=0.95$. Intervals: 2000 resamples of the batches within alternatives (seed 1), pairs of copies of one batch excluded, and 2000 resamples of whole alternatives. Scatter floor: 50 random permutations of the parts within each series. Components: rational subgroups of five consecutive parts; successive-difference standard deviation $\sqrt{\overline{\Delta^2}/2}$.
}%DIFDELCMD < 

%DIFDELCMD < %%%
\emph{\DIFdel{Decision rules.}} %DIFAUXCMD
\DIFdel{Adequacy share $p=0.95$; per-geometry floors; requirements below zero are left out for these non-negative characteristics. The largest-residual margin is the $\lceil (n+1)\cdot 0.9\rceil$-th smallest absolute residual of $n$ calibration alternatives, the largest when this exceeds $n$. M5 and M5n: Gaussian-process regression with a constant times anisotropic squared-exponential kernel plus white noise, normalised targets, three optimiser restarts (random state 0), length scales in $[0.1, 100]$; the noise level is bounded below by the mean squared block-bootstrap standard error of the calibration alternatives' $q_{95}$ relative to the variance of the targets; descriptors are a convex-geometry indicator, the blank-holder force in units of 100\,kN and the lubrication rank (coarse 0, medium 1, fine 2, the order of the oil amount), with descriptors that are constant over the calibration set dropped. NN: mean $q_{95}$ of the calibration alternatives of the target's geometry at the nearest blank-holder force (both neighbours when equally near; the other geometry when the family has none). Mc: for each lubrication pattern, the friction coefficient of the DDACS grid at 0.99\,mm whose simulated mid-side and corner draw-in, in floors, are closest in least squares to the centres of the calibration alternatives with that pattern (all calibration alternatives if none has it). M3 and M4: histogram gradient boosting (maximum depth 3, 200 iterations, learning rate 0.08, random state 0) on the geometry indicator, blank-holder force, two lubrication indicators, sheet thickness and oil film; 500 incoming conditions are drawn for the new alternative from calibration parts with the same lubrication pattern, with in-sample residuals pooled over the calibration alternatives resampled; offset and margin from a leave-one-out over the calibration alternatives.
}%DIFDELCMD < 

%DIFDELCMD < %%%
\emph{\DIFdel{Distances and intervals.}} %DIFAUXCMD
\DIFdel{The shares of correct and wrong verdicts are evaluated exactly on a grid of 0.05 floors up to 600 floors. Intervals: 1000 resamples of the held-out alternatives within each geometry (seeds derived from the scope), each with one of 1000 block-bootstrap replicates of every alternative's $q_{95}$ (batches of 50 resampled within the series); the resamples are shared by all rules of a scope. Refitting bootstrap: 200 resamples of the alternatives within each geometry, every alternative-level rule refitted.
}%DIFDELCMD < 

%DIFDELCMD < %%%
\emph{\DIFdel{Budget, guard, scenarios, robustness.}} %DIFAUXCMD
\DIFdel{The budget uses all subsets of $k=1,\dots,8$ calibration alternatives of the family (M5 }\DIFdelend \DIFaddbegin \DIFadd{Table~\ref{tab:notation} summarises the notation and the terms defined in Section~\ref{sec:method}, }\DIFaddend and \DIFdelbegin \DIFdel{M5n from $k=2$), bands from 200 resamples of the held-out alternatives. The guards and the tolerance scenarios use the stored intervals of every rule. ISO~2768-1 angular rows by the shorter leg: the wall (about 28\,mm) and the concave arm (about 18\,mm) over 10 up to 50\,mm, the convex arm (about 9\,mm) up to 10\,mm. The temperature correction subtracts $b\,(T-\bar T_a)$ from each characteristic, with $b$ the pooled within-alternative slope on the punch temperature $T$. The smoothed simulation replaces the simulated values of each geometry and force by a quadratic response surface in sheet thickness and friction coefficient.
}%DIFDELCMD < 

%DIFDELCMD < %%%
\emph{\DIFdel{Software.}} %DIFAUXCMD
\DIFdel{Python 3.11 with NumPy 2.4, pandas 3.0, SciPy 1.17, scikit-learn 1.9 \mbox{%DIFAUXCMD
\citep{pedregosa2011scikit} }\hskip0pt%DIFAUXCMD
and Matplotlib 3.11.
}%DIFDELCMD < 

%DIFDELCMD < %%%
\subsection{\DIFdel{Additional tables}}%DIFAUXCMD
\addtocounter{subsection}{-1}%DIFAUXCMD
%DIFDELCMD < \label{app:tables}
%DIFDELCMD < 

%DIFDELCMD < %%%
\DIFdel{Table~\ref{tab:distci} gives the bootstrap intervals of the decisive and safe distances, Table~\ref{tab:onesided} the one-sided safe distances, Table~\ref{tab:coverage} the coverage and width of the interval rules, Table~\ref{tab:paired} the paired differences between rules, Table~\ref{tab:qcdist} the distances per characteristic, Table~\ref{tab:tuning} the part-level rules under other learners, Table~\ref{tab:gp} the hyperparameters of the Gaussian-process fits, Table~\ref{tab:scan} the redundancy of the scans and Table~\ref{tab:simnoise} the extraction noise of the simulations. Two analyses are provided with the code but not used here: the sensitivities of the simulation over the geometric corners of the simulation dataset, which have no production counterpart, and the in-line verifiability of the characteristics from the press-force record}\DIFdelend \DIFaddbegin \DIFadd{Table~\ref{tab:processes} maps the framework's requirements onto other processes. Computation details are given in the ESM}\DIFaddend .

%DIF < % BEGIN GENERATED distci
\DIFdelbegin %DIFDELCMD < \begin{table}[t]
%DIFDELCMD < \tabsize
%DIFDELCMD < \setlength{\tabcolsep}{4pt}
%DIFDELCMD < \caption{Bootstrap 95\,\% intervals of the decisive and safe distances (floors)}\label{tab:distci}
%DIFDELCMD < \begin{tabular}{@{}lrrrrrr@{}}
%DIFDELCMD < \toprule
%DIFDELCMD <  & \multicolumn{2}{c}{New variant} & \multicolumn{2}{c}{New setting} & \multicolumn{2}{c}{New family} \\
%DIFDELCMD < \cmidrule(lr){2-3}\cmidrule(lr){4-5}\cmidrule(lr){6-7}
%DIFDELCMD < Rule & $\delta_\mathrm{d}$ & $\delta_\mathrm{s}$ & $\delta_\mathrm{d}$ & $\delta_\mathrm{s}$ & $\delta_\mathrm{d}$ & $\delta_\mathrm{s}$ \\
%DIFDELCMD < \midrule
%DIFDELCMD < M0 & 101--115 & 101--115 & 101--115 & 101--115 & 101--115 & 101--115 \\
%DIFDELCMD < M0w & 118--119 & 77--100 & 118--119 & 77--100 & 118--119 & 77--100 \\
%DIFDELCMD < M1 & 12--18 & 12--18 & 21--25 & 21--25 & 33--37 & 33--37 \\
%DIFDELCMD < Mc & 24--29 & 24--29 & 34--40 & 34--40 & 36--45 & 36--45 \\
%DIFDELCMD < M2 & 22--23 & 0.0--0.7 & 17--18 & 6.4--9.9 & 50--52 & 14--23 \\
%DIFDELCMD < M3 & 16--21 & 0.0--0.3 & 25--34 & 14--17 & 43--50 & 18--23 \\
%DIFDELCMD < M5 & 6.1--7.6 & 0.4--3.5 & 18--20 & 5.3--8.9 & 35--38 & 27--32 \\
%DIFDELCMD < M1n & 5.5--10 & 5.5--10 & 9.0--12 & 9.0--12 & 129--131 & 129--131 \\
%DIFDELCMD < NN & 2.0--4.5 & 2.0--4.5 & 5.5--10 & 5.5--10 & 129--131 & 129--131 \\
%DIFDELCMD < M2n & 15--21 & 0.0--0.7 & 15--16 & 3.3--9.3 & 131--133 & 127--129 \\
%DIFDELCMD < M4 & 9.7--12 & 0.0--1.6 & 12--15 & 5.7--9.2 & 132--133 & 127--129 \\
%DIFDELCMD < M5n & 5.3--6.6 & 0.4--2.6 & 17--18 & 3.7--9.0 & 130--132 & 129--131 \\
%DIFDELCMD < \botrule
%DIFDELCMD < \end{tabular}
%DIFDELCMD < \footnotetext{1000 resamples of the held-out alternatives within each geometry, each with a block-bootstrap replicate of the true $q_{95}$; the fits are held fixed (refitting: Table~\ref{tab:robust}). A new family rests on two calibrations, one per direction.}
%DIFDELCMD < \end{table}
%DIFDELCMD < %%%
%DIF < % END GENERATED distci
%DIFDELCMD < 

%DIFDELCMD < %%%
%DIF < % BEGIN GENERATED onesided
%DIFDELCMD < \begin{table}[t]
%DIFDELCMD < \tabsize
%DIFDELCMD < \setlength{\tabcolsep}{4.5pt}
%DIFDELCMD < \caption{One-sided safe distances (floors): false accepts and false rejects}\label{tab:onesided}
%DIFDELCMD < \begin{tabular}{@{}lrrrrrr@{}}
%DIFDELCMD < \toprule
%DIFDELCMD <  & \multicolumn{2}{c}{New variant} & \multicolumn{2}{c}{New setting} & \multicolumn{2}{c}{New family} \\
%DIFDELCMD < \cmidrule(lr){2-3}\cmidrule(lr){4-5}\cmidrule(lr){6-7}
%DIFDELCMD < Rule & FA side & FR side & FA side & FR side & FA side & FR side \\
%DIFDELCMD < \midrule
%DIFDELCMD < M0 & 17 & 109 & 17 & 109 & 17 & 109 \\
%DIFDELCMD < M1 & 15 & 12 & 24 & 23 & 33 & 36 \\
%DIFDELCMD < M2 & 0.3 & 0.0 & 8.3 & 6.7 & 11 & 23 \\
%DIFDELCMD < M3 & 0.0 & 0.0 & 12 & 17 & 18 & 23 \\
%DIFDELCMD < M5 & 1.1 & 1.4 & 7.5 & 7.4 & 32 & 29 \\
%DIFDELCMD < NN & 2.8 & 3.4 & 10 & 5.3 & 169 & 128 \\
%DIFDELCMD < M4 & 0.5 & 0.6 & 9.1 & 6.8 & 165 & 126 \\
%DIFDELCMD < M5n & 0.9 & 1.0 & 9.1 & 2.3 & 168 & 128 \\
%DIFDELCMD < \botrule
%DIFDELCMD < \end{tabular}
%DIFDELCMD < \footnotetext{FA side: smallest distance beyond which at most 5\,\% of the verdicts on failing alternatives ($d<0$) are false accepts; FR side: the same for false rejects of adequate alternatives ($d\ge0$). Requirements below zero are left out.}
%DIFDELCMD < \end{table}
%DIFDELCMD < %%%
%DIF < % END GENERATED onesided
%DIFDELCMD < 

%DIFDELCMD < %%%
%DIF < % BEGIN GENERATED coverage
%DIFDELCMD < \begin{table}[t]
%DIFDELCMD < \tabsize
%DIFDELCMD < \setlength{\tabcolsep}{3.5pt}
%DIFDELCMD < \caption{Interval rules: coverage of the true $q_{95}$, interval width and the accuracy of the centre}\label{tab:coverage}
%DIFDELCMD < \begin{tabular}{@{}lrrrrrrrrr@{}}
%DIFDELCMD < \toprule
%DIFDELCMD <  & \multicolumn{3}{c}{New variant} & \multicolumn{3}{c}{New setting} & \multicolumn{3}{c}{New family} \\
%DIFDELCMD < \cmidrule(lr){2-4}\cmidrule(lr){5-7}\cmidrule(lr){8-10}
%DIFDELCMD < Rule & cov. & half & centre & cov. & half & centre & cov. & half & centre \\
%DIFDELCMD < \midrule
%DIFDELCMD < M0w & 37 & 12 & 105 & 37 & 12 & 105 & 37 & 12 & 105 \\
%DIFDELCMD < M2 & 89 & 7.6 & 8.6 & 53 & 5.1 & 12 & 47 & 7.5 & 34 \\
%DIFDELCMD < M3 & 93 & 8.4 & 7.5 & 40 & 7.4 & 21 & 40 & 7.1 & 36 \\
%DIFDELCMD < M5 & 72 & 1.2 & 3.4 & 56 & 5.0 & 12 & 0 & 0.8 & 32 \\
%DIFDELCMD < M2n & 90 & 5.9 & 9.1 & 53 & 4.1 & 10 & 12 & 5.5 & 131 \\
%DIFDELCMD < M4 & 86 & 2.4 & 4.8 & 39 & 2.3 & 9.9 & 6 & 3.0 & 131 \\
%DIFDELCMD < M5n & 71 & 1.1 & 2.8 & 52 & 4.3 & 11 & 0 & 0.6 & 131 \\
%DIFDELCMD < \botrule
%DIFDELCMD < \end{tabular}
%DIFDELCMD < \footnotetext{cov.: share of cases whose true $q_{95}$ lies in the interval (\%; nominal level 90\,\%; the maximum-residual margin of M2 and M2n attains at most $n/(n+1)$ with $n$ calibration alternatives); half: median half-width (floors); centre: decisive distance of the interval centre used as a point rule (floors).}
%DIFDELCMD < \end{table}
%DIFDELCMD < %%%
%DIF < % END GENERATED coverage
%DIFDELCMD < 

%DIFDELCMD < %%%
%DIF < % BEGIN GENERATED paired
%DIFDELCMD < \begin{table}[t]
%DIFDELCMD < \tabsize
%DIFDELCMD < \setlength{\tabcolsep}{5pt}
%DIFDELCMD < \caption{Paired differences of the decisive distance between rules (floors)}\label{tab:paired}
%DIFDELCMD < \begin{tabular}{@{}lrrr@{}}
%DIFDELCMD < \toprule
%DIFDELCMD < Pair & New variant & New setting & New family \\
%DIFDELCMD < \midrule
%DIFDELCMD < M5 $-$ M5n & +0.6 (0.0, +1.6) & $-$0.2 ($-$0.6, +1.6) & $-$93.3 ($-$95.4, $-$91.9) \\
%DIFDELCMD < M2 $-$ M2n & +3.9 (+0.4, +7.4) & +2.3 (+1.4, +3.2) & $-$81.7 ($-$82.5, $-$78.5) \\
%DIFDELCMD < M1 $-$ M1n & +6.2 (+4.6, +9.7) & +13.5 (+10.6, +14.9) & $-$95.4 ($-$96.6, $-$92.5) \\
%DIFDELCMD < M5 $-$ NN & +3.3 (+2.1, +5.1) & +9.4 (+7.5, +13.5) & $-$92.5 ($-$94.8, $-$91.4) \\
%DIFDELCMD < M5 $-$ M2 & $-$16.0 ($-$16.6, $-$13.7) & +0.3 ($-$0.1, +1.7) & $-$12.6 ($-$15.1, $-$11.6) \\
%DIFDELCMD < M5 $-$ M4 & $-$4.6 ($-$5.3, $-$2.7) & +4.5 (+2.9, +6.6) & $-$95.0 ($-$97.0, $-$93.6) \\
%DIFDELCMD < M3 $-$ M4 & +6.2 (+5.1, +11.1) & +18.2 (+11.7, +20.4) & $-$84.7 ($-$89.4, $-$82.0) \\
%DIFDELCMD < Mc $-$ M1 & +11.5 (+8.9, +12.9) & +14.8 (+13.4, +16.3) & +2.4 (+0.9, +9.1) \\
%DIFDELCMD < \botrule
%DIFDELCMD < \end{tabular}
%DIFDELCMD < \footnotetext{Difference and 95\,\% interval from the common resamples of Table~\ref{tab:distci}; a negative difference means that the first rule decides closer to the requirement. M5$-$M5n, M2$-$M2n and M1$-$M1n isolate what the simulation adds to a rule.}
%DIFDELCMD < \end{table}
%DIFDELCMD < %%%
%DIF < % END GENERATED paired
%DIFDELCMD < 

%DIFDELCMD < %%%
%DIF < % BEGIN GENERATED qcdist
%DIFDELCMD < \begin{table}[t]
%DIFDELCMD < \tabsize
%DIFDELCMD < \setlength{\tabcolsep}{2.6pt}
%DIFDELCMD < \caption{Decisive/safe distance per characteristic (floors)}\label{tab:qcdist}
%DIFDELCMD < \begin{tabular}{@{}lrrrrrrrrrr@{}}
%DIFDELCMD < \toprule
%DIFDELCMD <  & \multicolumn{5}{c}{New variant} & \multicolumn{5}{c}{New setting} \\
%DIFDELCMD < \cmidrule(lr){2-6}\cmidrule(lr){7-11}
%DIFDELCMD < Characteristic & M1 & M2 & M5 & NN & M5n & M1 & M2 & M5 & NN & M5n \\
%DIFDELCMD < \midrule
%DIFDELCMD < Draw-in, mid-side & 21 & 20/0.0 & 2.8/1.0 & 2.0 & 2.7/0.6 & 27 & 18/8.3 & 19/6.8 & 6.5 & 11/4.2 \\
%DIFDELCMD < Draw-in, corner & 12 & 6.1/0.5 & 2.6/0.4 & 4.5 & 2.1/0.3 & 16 & 6.6/2.2 & 6.4/2.0 & 6.7 & 18/4.9 \\
%DIFDELCMD < Flange waviness & 8.1 & 23/0.1 & 1.4/0.2 & 0.9 & 1.4/0.2 & 9.7 & 16/11 & 19/11 & 11 & 18/10 \\
%DIFDELCMD < Wall angle & 11 & 9.3/0.2 & 2.0/0.9 & 3.4 & 3.5/0.9 & 14 & 7.6/3.2 & 7.7/3.1 & 1.9 & 3.2/1.6 \\
%DIFDELCMD < Arm angle, cut & 8.4 & 27/0.3 & 6.2/2.1 & 2.8 & 6.1/1.6 & 22 & 27/8.1 & 34/8.1 & 5.9 & 8.1/4.4 \\
%DIFDELCMD < Cup depth & 19 & 28/0.3 & 5.1/2.7 & 3.3 & 5.2/3.1 & 24 & 15/13 & 16/13 & 15 & 22/13 \\
%DIFDELCMD < Bottom dome, cut & 14 & 15/0.6 & 11/9.3 & 7.0 & 10/5.8 & 14 & 18/2.8 & 15/4.8 & 4.5 & 8.9/2.9 \\
%DIFDELCMD < \botrule
%DIFDELCMD < \end{tabular}
%DIFDELCMD < \footnotetext{18 held-out alternatives per characteristic and scope (36 verdicts per distance and side); one wrong verdict moves a share by 2.8 percentage points.}
%DIFDELCMD < \end{table}
%DIFDELCMD < %%%
%DIF < % END GENERATED qcdist
%DIFDELCMD < 

%DIFDELCMD < %%%
%DIF < % BEGIN GENERATED tuning
%DIFDELCMD < \begin{table}[t]
%DIFDELCMD < \tabsize
%DIFDELCMD < \setlength{\tabcolsep}{4pt}
%DIFDELCMD < \caption{Part-level rules M3 and M4 under other learners and settings (decisive/safe, floors)}\label{tab:tuning}
%DIFDELCMD < \begin{tabular}{@{}lrrrrrr@{}}
%DIFDELCMD < \toprule
%DIFDELCMD <  & \multicolumn{2}{c}{New variant} & \multicolumn{2}{c}{New setting} & \multicolumn{2}{c}{New family} \\
%DIFDELCMD < \cmidrule(lr){2-3}\cmidrule(lr){4-5}\cmidrule(lr){6-7}
%DIFDELCMD < Learner & M3 & M4 & M3 & M4 & M3 & M4 \\
%DIFDELCMD < \midrule
%DIFDELCMD < reference & 17/0.0 & 11/0.6 & 32/16 & 13/7.9 & 48/22 & 133/128 \\
%DIFDELCMD < regularised & 20/0.3 & 11/0.4 & 26/15 & 13/7.7 & 50/27 & 133/129 \\
%DIFDELCMD < flexible & 23/0.0 & 11/0.6 & 36/14 & 14/7.8 & 47/26 & 133/129 \\
%DIFDELCMD < ridge & 28/0.4 & 12/0.3 & 33/9.6 & 15/4.9 & 50/21 & 133/128 \\
%DIFDELCMD < quantile loss & 18/0.0 & 11/0.5 & 30/18 & 15/8.5 & 50/27 & 133/129 \\
%DIFDELCMD < selected & 16/0.3 & 11/0.6 & 26/17 & 13/8.2 & 47/26 & 133/129 \\
%DIFDELCMD < \botrule
%DIFDELCMD < \end{tabular}
%DIFDELCMD < \footnotetext{Reference: histogram gradient boosting, depth 3, 200 iterations, rate 0.08. Regularised: depth 2, 100 iterations, rate 0.05, L2 penalty 1. Flexible: depth 6, 400 iterations, rate 0.1. Ridge: linear regression on the standardised inputs. Quantile loss: boosting of the 95th percentile, no resampled residuals. Selected: the boosting setting with the smallest leave-one-out error over the calibration alternatives.}
%DIFDELCMD < \end{table}
%DIFDELCMD < %%%
%DIF < % END GENERATED tuning
%DIFDELCMD < 

%DIFDELCMD < %%%
%DIF < % BEGIN GENERATED gp
%DIFDELCMD < \begin{table}[t]
%DIFDELCMD < \tabsize
%DIFDELCMD < \setlength{\tabcolsep}{5pt}
%DIFDELCMD < \caption{Gaussian-process fits with a hyperparameter at a bound}\label{tab:gp}
%DIFDELCMD < \begin{tabular}{@{}lrrrrrr@{}}
%DIFDELCMD < \toprule
%DIFDELCMD <  & \multicolumn{3}{c}{M5} & \multicolumn{3}{c}{M5n} \\
%DIFDELCMD < \cmidrule(lr){2-4}\cmidrule(lr){5-7}
%DIFDELCMD < Scope & fits & noise (\%) & length (\%) & fits & noise (\%) & length (\%) \\
%DIFDELCMD < \midrule
%DIFDELCMD < New variant & 126 & 38 & 48 & 126 & 35 & 28 \\
%DIFDELCMD < New setting & 42 & 45 & 60 & 42 & 38 & 62 \\
%DIFDELCMD < Pooled variant & 126 & 21 & 57 & 126 & 2 & 56 \\
%DIFDELCMD < New family & 14 & 36 & 36 & 14 & 36 & 21 \\
%DIFDELCMD < \botrule
%DIFDELCMD < \end{tabular}
%DIFDELCMD < \footnotetext{Share of the fits (calibration set and characteristic) whose noise level sits at its lower bound, the sampling variance of the calibration alternatives' $q_{95}$, and with at least one length scale at a bound (0.1 or 100 in units of 100\,kN or one lubrication step).}
%DIFDELCMD < \end{table}
%DIFDELCMD < %%%
%DIF < % END GENERATED gp
%DIFDELCMD < 

%DIFDELCMD < %%%
%DIF < % BEGIN GENERATED scan
%DIFDELCMD < \begin{table}[t]
%DIFDELCMD < \tabsize
%DIFDELCMD < \setlength{\tabcolsep}{6pt}
%DIFDELCMD < \caption{Redundancy of the scans: ranges over the nine series of each geometry}\label{tab:scan}
%DIFDELCMD < \begin{tabular}{@{}lrr@{}}
%DIFDELCMD < \toprule
%DIFDELCMD < Quantity & Concave & Convex \\
%DIFDELCMD < \midrule
%DIFDELCMD < Flange width $W_x$, SD (mm) & 0.07--0.15 & 0.05--0.11 \\
%DIFDELCMD < Flange width $W_y$, SD (mm) & 0.27--0.54 & 0.43--0.92 \\
%DIFDELCMD < Correlation of $W_x$ and $W_y$ & $-$0.14 to 0.20 & $-$0.21 to 0.09 \\
%DIFDELCMD < $W_y-W_x$ (mm) & 2.5--3.5 & 2.3--3.5 \\
%DIFDELCMD < Wall angle, mean of E and W, SD ($^\circ$) & 0.04--0.10 & 0.06--0.13 \\
%DIFDELCMD < Wall angle N, SD ($^\circ$) & 0.31--0.35 & 0.35--0.38 \\
%DIFDELCMD < Wall angle N minus mean of E and W ($^\circ$) & 0.0--0.4 & $-$1.6 to $-$0.9 \\
%DIFDELCMD < Correlation of the E and W wall angles & $-$0.10 to 0.27 & $-$0.71 to $-$0.49 \\
%DIFDELCMD < \botrule
%DIFDELCMD < \end{tabular}
%DIFDELCMD < \footnotetext{$W_x$, $W_y$: flange width across the cup centre, across and along the scan lines; E, W, N: east, west and north walls; SD: standard deviation over the parts of a series; correlation over the parts of a series.}
%DIFDELCMD < \end{table}
%DIFDELCMD < %%%
%DIF < % END GENERATED scan
\DIFdelend \DIFaddbegin \begin{table}[h]
\tabsize
\setlength{\tabcolsep}{4pt}
\caption{Notation and terms}\label{tab:notation}
\begin{tabular}{@{}L{2.8cm}L{9.4cm}@{}}
\toprule
Symbol, term & Meaning \\
\midrule
$a$, $c$, $q_p(a,c)$ & alternative (design--process combination), quality characteristic, $p$-quantile of the parts of $a$ ($p=0.95$) \\
$R$, $d$ & requirement (upper limit on $c$ or $|c|$); its distance from the truth in floors, $R=q_p+dF_c$ \\
$[L_r,U_r]$, $V_r$ & interval of rule $r$ for $q_p$; verdict meets, fails or trial (Eq.~\ref{eq:verdict}) \\
$B$, $m$, $F_c$ & batch size (50); batch centre (10\,\% trimmed mean); production floor (Eq.~\ref{eq:floor}) \\
$\sigma_\mathrm{b}$, $\sigma_\mathrm{w}$ & standard deviation between batch means and between parts within a batch \\
ESR, MRC & effect-to-scatter ratio (Eq.~\ref{eq:esr}); minimum resolvable change, $F_c/|s_c|$ \\
$u$, $v$ & exclusion distances of an interval from the truth, from above and below, in floors \\
$\kappa(\delta)$, $\omega(\delta)$ & shares of correct and wrong verdicts among the feasible verdicts at $|d|=\delta$ \\
$\delta_\mathrm{d}$, $\delta_\mathrm{s}$ & decisive distance ($\kappa\ge0.95$ beyond it) and safe distance ($\omega\le0.05$ beyond it) \\
FA, FR side & one-sided distances for failing alternatives (false accepts) and adequate ones (false rejects) \\
$g$ & guard band: observable margin beyond which $\ge$95\,\% of decided verdicts were correct \\
Calibration alternatives & produced alternatives whose record calibrates a rule \\
Evidence relation & sibling (same process setting produced), new process setting (interpolation, extrapolation), new family \\
Calibration budget & distances against the number $k$ of produced alternatives, by relation \\
Guard & family, range and novelty checks that send a design outside the calibration to trial \\
\botrule
\end{tabular}
\end{table}
\DIFaddend 

%DIF < % BEGIN GENERATED simnoise
\DIFdelbegin %DIFDELCMD < \begin{table}[t]
%DIFDELCMD < \tabsize
%DIFDELCMD < \setlength{\tabcolsep}{6pt}
%DIFDELCMD < \caption{Extraction noise of the simulated characteristics (floors)}\label{tab:simnoise}
%DIFDELCMD < \begin{tabular}{@{}lrrrr@{}}
%DIFDELCMD < \toprule
%DIFDELCMD <  & \multicolumn{2}{c}{Concave} & \multicolumn{2}{c}{Convex} \\
%DIFDELCMD < \cmidrule(lr){2-3}\cmidrule(lr){4-5}
%DIFDELCMD < Characteristic & noise & envelope & noise & envelope \\
%DIFDELCMD < \midrule
%DIFDELCMD < Draw-in, mid-side & 0.1 & 0.8 & 0.1 & 0.7 \\
%DIFDELCMD < Draw-in, corner & 0.1 & 0.5 & 0.2 & 1.3 \\
%DIFDELCMD < Flange waviness & 0.4 & 2.9 & 4.0 & 4.2 \\
%DIFDELCMD < Wall angle & 3.1 & 2.0 & 1.7 & 1.5 \\
%DIFDELCMD < Arm angle, cut & 0.9 & 1.5 & 0.9 & 2.5 \\
%DIFDELCMD < Cup depth & 0.8 & 1.5 & 0.7 & 1.6 \\
%DIFDELCMD < Bottom dome, cut & 9.8 & 20.9 & 2.4 & 9.9 \\
%DIFDELCMD < \botrule
%DIFDELCMD < \end{tabular}
%DIFDELCMD < \footnotetext{Noise: standard deviation of the second difference of the simulated value along the friction grid over $\sqrt{1.5}$ (white noise of the extraction). Envelope: largest difference between the maximum of the simulated values of a geometry and force and the maximum of a quadratic response surface over thickness and friction.}
%DIFDELCMD < \end{table}
%DIFDELCMD < %%%
%DIF < % END GENERATED simnoise
\DIFdelend \DIFaddbegin \begin{table}[h]
\tabsize
\setlength{\tabcolsep}{3pt}
\caption{What the framework requires, mapped onto other processes}\label{tab:processes}
\begin{tabular}{@{}L{1.9cm}L{2.4cm}L{2.6cm}L{2.5cm}L{2.4cm}@{}}
\toprule
Process & Batch and run & Dominant variation & Evidence relations & Floor \\
\midrule
Sheet forming (this paper) & consecutive parts of one series & drift within a run, coil and die state between runs & variant, process setting, family & batch centres within a run \\
Injection moulding & consecutive shots of one cavity; one mould set-up & cavity, shot, set-up & cavity, process window, new mould & nested: within cavity, between cavities \\
Machining & consecutive parts between offsets & tool wear, compensated by offsets & tool path, cutting data, new part & between offset intervals \\
Casting & castings of one pattern and melt & melt, pattern position & alloy, gating, new pattern & between melts \\
Additive manufacturing & replicated specimens of one build & build, position, orientation & orientation, parameters, new geometry & between builds, from replicated anchors \\
\botrule
\end{tabular}
\end{table}
\DIFaddend 

\DIFdelbegin %DIFDELCMD < \FloatBarrier
%DIFDELCMD < %%%
\DIFdelend \end{appendices}

%DIF > % ----------------------------------------------------------------------
{\DIFdelbegin %DIFDELCMD < \makeatletter\def\bibfont{\reset@font\fontfamily{\rmdefault}\footnotesize\selectfont}\makeatother
%DIFDELCMD < \setlength{\bibsep}{0.05em}
%DIFDELCMD < %%%
\DIFdelend \DIFaddbegin \small
\DIFaddend \bibliography{refs}}

\end{document}
